\documentclass[11pt]{article}

\usepackage[utf8]{inputenc}
\usepackage[T1]{fontenc}
\usepackage[margin=1in]{geometry}
\usepackage{amsmath,amssymb}
\usepackage{graphicx}
\usepackage{longtable,booktabs}
\usepackage{microtype}
% hyperref goes last; it makes \cite clickable through to the bibliography.
\usepackage[colorlinks=true,linkcolor=blue,citecolor=teal,urlcolor=magenta]{hyperref}

% pandoc emits \tightlist but expects the preamble to define it.
\providecommand{\tightlist}{%
  \setlength{\itemsep}{0pt}\setlength{\parskip}{0pt}}

% If the compile times out on Overleaf, comment out \tableofcontents below first.
\title{A Comprehensive Survey on Large Multi-Modal Language Models}
\author{Generated by AutoSurvey}
\date{}

\begin{document}
\maketitle
\tableofcontents
\newpage

\section{Introduction to Large Multi-Modal Language Models}

\subsection{Background and Motivation}

The rapid advancements in large language models (LLMs) and the emergence of multi-modal capabilities have revolutionized the field of natural language processing (NLP) and opened up new avenues for artificial intelligence (AI) research and development. LLMs, such as GPT-3 \cite{ref1}, BERT \cite{ref2}, and T5 \cite{ref3}, have demonstrated remarkable performance in a wide range of language-related tasks, including text generation, question answering, and language translation.

The significance and potential of large multi-modal language models cannot be overstated. These models, often referred to as Multi-Modal Large Language Models (MM-LLMs), have shown the ability to understand and generate content that seamlessly integrates different modalities, such as text, images, and audio, enabling applications such as visual question answering, image captioning, and multimodal dialogue \cite{ref4,ref5}. By integrating language understanding with other modalities, these models can enable more natural and intuitive human-computer interaction, enhance decision-making processes, and facilitate the development of intelligent systems that can better understand and respond to the complexities of the real world \cite{ref6}.

The rapid progress in this field has been driven by several key factors, including the availability of large and diverse datasets, advancements in deep learning architectures, and the increasing computational power of modern hardware \cite{ref7}. Specifically, the development of transformer-based models, such as GPT and BERT, has been instrumental in the success of LLMs, enabling them to capture more complex language patterns and dependencies \cite{ref8}.

Moreover, the emergence of multi-modal models has been facilitated by the integration of various modalities, such as vision and audio, into the language modeling process \cite{ref9}. This has led to the creation of models that can not only understand and generate text but also perceive and generate content in other modalities, enabling more comprehensive and versatile AI systems \cite{ref10}.

The potential of large multi-modal language models extends beyond traditional language-based tasks. These models have shown promise in areas such as scientific research, where they can assist in literature review, data analysis, and knowledge extraction \cite{ref11}. Additionally, they have the potential to revolutionize industries like telecommunications, where they can enhance customer service, automate task completion, and improve network operations \cite{ref12}.

However, the development and deployment of large multi-modal language models also come with significant challenges and ethical considerations. These models require vast amounts of computational resources and energy to train and deploy, raising concerns about sustainability and accessibility \cite{ref13}. Additionally, there are concerns about the potential for bias, privacy violations, and the misuse of these models, which necessitate the development of robust ethical frameworks and governance mechanisms \cite{ref14}.

Despite these challenges, the rapid advancements in large multi-modal language models have generated widespread excitement and optimism about their transformative potential. As research in this field continues to evolve, it is crucial to understand the underlying mechanisms, capabilities, and limitations of these models, as well as their societal implications, to ensure that their development and deployment are guided by principles of responsible AI \cite{ref15}.

\subsection{Defining Multi-Modal Language Models}

Multimodal language models, also known as multi-modal language models, are a class of large language models that go beyond the traditional text-only domain, incorporating information from various modalities such as vision, audio, and even tactile or olfactory data \cite{ref16,ref17}. These models aim to achieve a more comprehensive and grounded understanding of language by leveraging the rich sensory information that humans naturally utilize in communication.

Compared to traditional language models, multi-modal language models possess several distinguishing features. Firstly, they are designed to process and integrate inputs from multiple modalities, rather than relying solely on textual data \cite{ref18}. This enables them to capture the interplay between different forms of information, such as the relationship between the visual appearance of an object and its linguistic description \cite{ref19}.

The key components that enable the integration of diverse modalities in multi-modal language models typically involve specialized architectural designs and training strategies. At the architectural level, these models often incorporate modality-specific encoders or fusion modules to effectively capture the unique characteristics of each input \cite{ref20,ref21}. For example, vision-specific encoders may utilize convolutional neural networks to extract visual features, while audio-specific encoders may employ time-frequency representations and recurrent neural networks to process audio signals.

The fusion of these modality-specific representations is a crucial aspect of multi-modal language models. Various fusion strategies have been explored, such as concatenation, attention-based mechanisms, and dedicated fusion layers \cite{ref16,ref22}. These fusion techniques aim to capture the complex interactions and interdependencies between the different modalities, enabling the model to leverage the complementary information they provide.

In addition to architectural design, the training strategies for multi-modal language models are also tailored to handle the challenges posed by multi-modal data. One common approach is joint training, where the model is trained on a large corpus of aligned multimodal data, such as image-caption pairs or video-transcript pairs \cite{ref23}. This allows the model to learn the associations and correspondences between different modalities during the training process.

Another prevalent training strategy is multimodal pre-training, where the model is first trained on a large, diverse dataset of multimodal data, and then fine-tuned on specific downstream tasks \cite{ref17,ref21}. This approach aims to capture the general patterns and relationships within and across modalities, which can then be effectively leveraged for various applications.

The integration of multimodal information in language models also raises interesting challenges and considerations. For instance, the alignment and synchronization of data from different modalities, which may have varying sampling rates or asynchronous temporal dynamics, can be a non-trivial problem \cite{ref24}. Additionally, the efficient handling of modality-specific features and the effective fusion of these features within the language model architecture are active areas of research \cite{ref25}.

Despite these challenges, the development of multi-modal language models has demonstrated impressive capabilities in a wide range of applications, from multimodal question answering and image captioning to cross-modal retrieval and multimodal reasoning \cite{ref26,ref27}. As the field continues to evolve, we can expect to see further advancements in the integration of diverse modalities, the enhancement of multimodal understanding, and the expansion of applications that leverage the unique strengths of these powerful models.

\subsection{Architectural Innovations}

The rapid advancements in large multi-modal language models have been enabled by various innovative architectural designs that effectively integrate multiple modalities, including vision, audio, and language. These architectural innovations have played a crucial role in empowering large language models to comprehend and generate content across diverse modalities.

One prominent architectural approach is the use of multi-modal fusion modules that seamlessly combine the representations from different modalities. For instance, the FLAVA model \cite{ref27} employs a holistic architecture that targets all modalities at once, including vision and language, using a single foundation model. The model utilizes cross-attention mechanisms to facilitate the integration of visual and textual features, enabling the model to excel at a wide range of vision-language tasks.

Similarly, the PaLM-E model \cite{ref6} introduces an embodied language model that directly incorporates real-world continuous sensor modalities, such as vision and audio, into the language model. By aligning the output of the language model with the input of generative models for various modalities, PaLM-E establishes a direct link between words and percepts, enabling the model to ground language in the physical world.

Another architectural innovation is the use of modality-specific encoders and decoders, which allow for more effective representation learning and information exchange across modalities. The TriBERT architecture \cite{ref28} is a prime example, where the model employs separate encoders for vision, pose, and audio, and utilizes flexible co-attention mechanisms to enable contextual feature learning across these modalities.

The UnIVAL model \cite{ref10} takes a step further by unifying multiple modalities, including text, images, video, and audio, within a single framework. The model employs a modular architecture with separate encoders for each modality, which are then aligned through a multimodal fusion module. This unified approach enables the model to effectively handle a diverse set of tasks and modalities without relying on large dataset sizes or massive model parameters.

Additionally, architectural innovations have also focused on enhancing the interpretability and transparency of large multi-modal language models. The LVLM-Intrepret framework \cite{ref29} introduces an interactive application that helps users systematically investigate the internal mechanisms of large vision-language models, allowing them to uncover the model's strengths, limitations, and failure modes.

Furthermore, researchers have explored methods for bridging the gap between different modalities, such as the use of multimodal adapters and attention gates. The PILL architecture \cite{ref30} introduces a novel approach that utilizes mixture-of-modality-adapter-experts to independently handle different modalities, enabling better adaptation to downstream tasks while preserving the expressive capability of the original model. The modality-attention-gating mechanism in PILL further enhances the model's ability to adaptively control the contribution of different modality tokens to the overall representation.

The ModaVerse framework \cite{ref16} presents a unique approach that operates directly at the level of natural language, aligning the language model's output with the input of generative models for various modalities. This I/O alignment mechanism avoids the complexities associated with latent feature alignments, simplifying the training process and reducing the computational and data costs compared to traditional multi-modal language models.

These architectural innovations have paved the way for the development of large multi-modal language models that can effectively handle diverse modalities, including vision, audio, and language, while leveraging their respective strengths and facilitating seamless cross-modal interactions. The continued advancements in architectural design, along with the integration of various modalities, have the potential to unlock new frontiers in the field of artificial intelligence, bringing us closer to the realization of true multi-modal intelligence.

\subsection{Training Strategies and Approaches}

The development of large multi-modal language models has been driven by various training strategies and approaches that aim to effectively integrate and leverage the diverse modalities of data, such as text, images, and audio. These training approaches can be broadly categorized into three main strategies: joint training, iterative fine-tuning, and self-supervised pre-training on multimodal data.

Joint training, also known as multi-task learning, involves training a single model simultaneously on multiple modalities and tasks, allowing the model to learn shared representations and capture the inherent connections between the different modalities \cite{ref31}. This approach has been widely adopted in the development of large multi-modal language models, where the model is trained on a diverse set of tasks, such as image captioning, visual question answering, and text-to-image generation, within a unified framework \cite{ref32}.

One of the key advantages of joint training is its ability to leverage the synergies between different modalities, enabling the model to learn more robust and generalizable representations \cite{ref33}. By exposing the model to a wide range of multi-modal tasks and data during training, the model can learn to effectively integrate and reason across different modalities, leading to improved performance on downstream tasks \cite{ref34}. However, joint training can also be challenging, as it requires carefully balancing the optimization of different task-specific objectives and handling the potential negative transfer between tasks \cite{ref35}. To address these challenges, researchers have explored various techniques, such as adaptive task weighting, task-specific attention mechanisms, and progressive training approaches \cite{ref36}.

Iterative fine-tuning, on the other hand, involves a two-stage training process \cite{ref37}. In the first stage, a large multi-modal language model is pre-trained on a diverse set of multimodal data, often using self-supervised objectives such as masked language modeling and contrastive learning \cite{ref38}. This pre-training stage allows the model to learn rich cross-modal representations that capture the underlying relationships between different modalities. In the second stage, the pre-trained model is fine-tuned on specific downstream tasks or datasets, enabling the model to adapt and specialize its representations to the task at hand \cite{ref39}. This iterative fine-tuning approach has been shown to be effective in improving the performance of large multi-modal language models on a wide range of tasks, while also allowing for efficient transfer learning \cite{ref40}.

In addition to joint training and iterative fine-tuning, self-supervised pre-training on multimodal data has also emerged as a popular approach for developing large multi-modal language models \cite{ref41}. In this approach, the model is pre-trained on a large corpus of multimodal data, such as image-text pairs or video-text pairs, using self-supervised objectives that encourage the model to learn cross-modal representations \cite{ref42}. The self-supervised pre-training stage allows the model to capture the inherent relationships and patterns within and across different modalities, without the need for explicit supervision \cite{ref43}. This learned multimodal understanding can then be fine-tuned or adapted to various downstream tasks, enabling the model to achieve strong performance on a wide range of multi-modal applications \cite{ref44}.

Overall, the development of large multi-modal language models has been driven by a diverse set of training strategies and approaches, each with its own strengths and trade-offs. Joint training, iterative fine-tuning, and self-supervised pre-training on multimodal data have all played crucial roles in advancing the state-of-the-art in this rapidly evolving field, and continued research in these areas is likely to further enhance the capabilities and versatility of large multi-modal language models.

\subsection{Multimodal Data and Datasets}

The emergence of large multi-modal language models (MLLMs) has been driven by the availability of comprehensive multimodal datasets for model training and evaluation. These datasets play a crucial role in enabling MLLMs to effectively process and understand a wide range of modalities, including text, images, videos, and other forms of multimedia \cite{ref31,ref45}.

One of the key advantages of these multimodal datasets is their ability to capture the complex interplay between different modalities, allowing MLLMs to learn the nuanced relationships and associations between them \cite{ref46,ref47}. This is particularly important in real-world scenarios, where information is often conveyed through a combination of modalities, such as text accompanied by images or videos.

Existing multimodal datasets vary in their size, diversity, and focus, catering to different research and application needs. For example, the WIT dataset \cite{ref47} provides a massive collection of over 37.6 million image-text pairs spanning 108 languages, enabling the development of multilingual MLLMs. In contrast, the SEED-Bench dataset \cite{ref46} focuses on evaluating the generative comprehension abilities of MLLMs, with a curated set of over 19,000 multiple-choice questions that target various aspects of multimodal understanding.

Another notable dataset is the LAMM (Language-Assisted Multi-Modal Instruction-Tuning Dataset, Framework, and Benchmark) \cite{ref48}, which provides a comprehensive set of instructions and tasks spanning 2D and 3D vision, enabling the assessment of MLLMs' ability to follow complex, multimodal instructions. The 3DBench dataset \cite{ref49} further expands the scope of multimodal datasets by introducing a 3D-centric benchmark, focusing on spatial understanding and reasoning tasks.

Beyond these curated datasets, researchers have also explored the use of automatically generated multimodal data to augment training and evaluation. For example, the MLLM-DataEngine \cite{ref50} framework proposes an iterative approach to generating high-quality multimodal data, leveraging language models to create incremental datasets that target specific weaknesses of the MLLM being trained.

The diversity and richness of these multimodal datasets have been instrumental in the development of MLLMs, allowing researchers to train and evaluate these models on a wide range of tasks and scenarios. However, the rapidly evolving nature of this field and the increasing complexity of multimodal information also pose new challenges. Researchers are continuously exploring methods to scale these datasets, improve their quality and diversity, and enable more comprehensive and robust evaluation of MLLMs \cite{ref51,ref52}.

One emerging trend is the use of synthetic data generation techniques, such as those employed in the MLLM-DataEngine \cite{ref50}, to augment existing datasets and create more diverse and challenging multimodal samples. Additionally, the integration of multimodal data from diverse sources, including web-crawled content, domain-specific repositories, and crowdsourced annotations, is becoming increasingly important to capture the breadth and complexity of real-world multimodal information \cite{ref53}.

As the field of MLLMs continues to evolve, the development and utilization of comprehensive multimodal datasets will remain a crucial aspect of enabling these models to achieve human-level understanding and reasoning across diverse modalities and applications. The continued efforts in this direction, coupled with advancements in model architectures and training strategies, will be instrumental in realizing the full potential of large multi-modal language models.

\subsection{Benchmarking and Evaluation Frameworks}

The rapid advancements in large multi-modal language models (MMLMs) have led to the emergence of a wide range of benchmarking and evaluation frameworks to assess their capabilities. These frameworks play a crucial role in understanding the performance, limitations, and potential of these models, as well as guiding their future development.

One of the primary reasons for the importance of benchmarking and evaluation frameworks for MMLMs is the sheer complexity and multifaceted nature of these models. MMLMs are designed to handle a diverse range of tasks and modalities, from language understanding and generation to visual perception and reasoning. Assessing their performance across this broad spectrum requires a comprehensive and well-designed evaluation approach \cite{ref51,ref54}.

Existing benchmarking and evaluation frameworks for MMLMs can be broadly categorized into two main approaches: task-specific benchmarks and holistic evaluation frameworks. Task-specific benchmarks focus on evaluating the performance of MMLMs on individual tasks, such as image captioning, visual question answering, or multimodal reasoning \cite{ref46,ref55}. These benchmarks provide valuable insights into the strengths and weaknesses of MMLMs in specific domains, allowing for a more targeted assessment of their capabilities.

On the other hand, holistic evaluation frameworks aim to provide a more comprehensive assessment of MMLMs by considering a wide range of tasks and metrics \cite{ref56,ref56}. These frameworks often incorporate multiple sub-tasks, ranging from language understanding and generation to visual perception and reasoning, as well as various evaluation metrics that assess the models' performance, robustness, fairness, and safety \cite{ref57}. By considering the models' performance across a diverse set of tasks and metrics, these holistic frameworks offer a more complete picture of their capabilities and limitations.

One of the key challenges in benchmarking and evaluating MMLMs is the need for standardized and robust evaluation protocols. Existing benchmarks and frameworks often rely on human-annotated datasets, which can be prone to biases and inconsistencies \cite{ref58,ref59}. To address this, researchers have explored the use of automated evaluation methods, such as language model-based evaluators, to provide more scalable and objective assessments \cite{ref60,ref61}.

Another important aspect of benchmarking and evaluation frameworks for MMLMs is the consideration of multimodal and multilingual capabilities. Many existing benchmarks and frameworks have focused primarily on English-language tasks and unimodal inputs, overlooking the challenges and opportunities presented by the integration of multiple modalities and languages \cite{ref62,ref63}.

To address this, researchers have developed benchmarks and frameworks that specifically assess the performance of MMLMs in multimodal and multilingual settings \cite{ref48,ref64}. These efforts aim to ensure that MMLMs are evaluated more comprehensively, capturing their ability to handle diverse inputs and outputs across different languages and modalities.

Looking ahead, the continued development and refinement of benchmarking and evaluation frameworks for MMLMs will be crucial for driving further advancements in this rapidly evolving field. As MMLMs become more sophisticated and widely deployed, it will be essential to have robust and reliable evaluation methods that can keep pace with the rapid progress, identify emerging capabilities and limitations, and guide the future development of these models \cite{ref65,ref66}.

This may involve the incorporation of new evaluation metrics, the creation of more diverse and challenging datasets, and the exploration of novel evaluation techniques, such as meta-evaluation approaches that assess the models' ability to learn and adapt \cite{ref67,ref68}. By continuously advancing the state of benchmarking and evaluation for MMLMs, the research community can ensure that these powerful models are developed and deployed in a responsible and impactful manner, ultimately contributing to the realization of more capable and trustworthy artificial intelligence systems.

\subsection{Applications and Use Cases}

The rapid advancements in large multi-modal language models (MMLMs) have opened up a wide range of applications across diverse domains, showcasing their transformative potential in real-world scenarios. These models have demonstrated remarkable capabilities in seamlessly integrating and comprehending information from various modalities, including text, images, audio, and even video, enabling them to tackle complex tasks that were previously beyond the reach of traditional language models.

One of the most promising application areas for large multi-modal language models is the healthcare sector. These models have shown immense potential in enhancing various aspects of medical practice, from clinical decision support to patient communication. In the realm of medical diagnostics, large multi-modal language models have demonstrated the ability to assist clinicians in interpreting and analyzing medical images, such as radiological scans and pathology slides, by providing contextual information and insights that complement the visual data \cite{ref69,ref70}. This integration of visual and textual data has the potential to improve diagnostic accuracy, expedite the decision-making process, and ultimately lead to better patient outcomes.

Moreover, large multi-modal language models have exhibited the capacity to comprehend and process complex medical jargon, patient records, and research literature, enabling them to serve as virtual assistants in healthcare settings. These models can aid in tasks such as triaging patient inquiries, summarizing medical reports, and providing personalized treatment recommendations \cite{ref71,ref72}. By leveraging their language understanding and generation capabilities, these models can enhance communication between healthcare providers and patients, improving the overall quality of care and patient engagement.

In the education sector, large multi-modal language models have demonstrated the ability to revolutionize the way knowledge is disseminated and concepts are explained. These models can generate personalized learning materials, interactive visualizations, and tailored feedback for students, catering to diverse learning styles and abilities \cite{ref73}. By integrating visual, audio, and textual modalities, large multi-modal language models can enhance the understanding and retention of complex scientific and mathematical principles, making education more engaging and accessible to learners \cite{ref74}.

The entertainment industry has also witnessed the transformative potential of large multi-modal language models. These models have been utilized in the creation of interactive storytelling experiences, where they can generate immersive narratives that seamlessly combine text, images, and even audio \cite{ref29}. Furthermore, large multi-modal language models have the potential to revolutionize content creation, enabling the generation of high-quality multimedia content, from video scripts to music compositions, by understanding the underlying relationships between different modalities \cite{ref10}.

Beyond these specific domains, large multi-modal language models have demonstrated their versatility in a wide range of applications, from sustainable urban planning and climate modeling to fitness and well-being \cite{ref75}. These models have the ability to process and integrate diverse data sources, such as satellite imagery, sensor data, and natural language input, to provide insights and decision support in complex real-world scenarios \cite{ref76}.

Despite the remarkable progress and potential of large multi-modal language models, their widespread adoption and integration into various industries and sectors are not without challenges. Ensuring the safety, reliability, and ethical alignment of these models in high-stakes domains, such as healthcare and finance, remains a critical priority \cite{ref71,ref77}. Additionally, the development of robust evaluation frameworks and benchmarks to assess the performance and capabilities of these models across different modalities and tasks is an ongoing area of research \cite{ref78,ref29}.

Nevertheless, the impressive capabilities and versatility of large multi-modal language models have undoubtedly positioned them as transformative forces in diverse industries and sectors, with the potential to revolutionize how we interact with and leverage technology to address complex real-world challenges.

\subsection{Ethical Considerations and Societal Impact}

As the development and deployment of large multi-modal language models continue to advance, it is crucial to examine their ethical implications and potential societal impact. These models possess immense capabilities that can revolutionize various domains, but they also raise significant concerns around privacy, fairness, and the potential for misuse.

One of the primary ethical concerns surrounding large multi-modal language models is the issue of privacy. These models have the potential to memorize and leak pre-training data, exposing sensitive personal information \cite{ref79}. Additionally, the integration of these models into user-facing systems raises concerns about the potential for privacy breaches, as users may inadvertently disclose personal data during interactions \cite{ref80}. Strategies to mitigate these privacy risks, such as implementing robust data protection measures and providing users with transparent control over their personal information, need to be carefully considered \cite{ref81}.

Fairness is another crucial ethical consideration in the context of large multi-modal language models. These models have the potential to perpetuate and amplify societal biases, prejudices, and discrimination \cite{ref82,ref83}. Ensuring that these models are designed and deployed in a fair and equitable manner, with mechanisms in place to detect and mitigate biases, is essential to avoid exacerbating existing societal inequalities \cite{ref84}.

The potential for misuse of large multi-modal language models is another significant ethical concern. These models can be exploited for malicious purposes, such as generating disinformation, impersonating individuals, or creating harmful content \cite{ref85,ref86}. Developing robust safeguards and mechanisms to detect and prevent such misuse is crucial to ensure the responsible development and deployment of these technologies \cite{ref87}.

To address these ethical concerns and ensure the responsible development and deployment of large multi-modal language models, a multifaceted approach is necessary. This may involve the establishment of comprehensive ethical frameworks and guidelines, the implementation of rigorous testing and evaluation protocols, and the engagement of diverse stakeholders, including policymakers, industry leaders, and civil society \cite{ref88,ref89}.

Moreover, the societal impact of large multi-modal language models extends beyond individual privacy and fairness concerns. These models have the potential to disrupt various sectors, such as education, healthcare, and the workforce, with both positive and negative implications \cite{ref90}. Careful consideration must be given to the broader societal implications of these technologies, including their impact on employment, social dynamics, and environmental sustainability \cite{ref91}.

\subsection{Enhancing Multi-Modal Reasoning through Prompting and In-Context Learning}

The emergence of large language models (LLMs) and their integration with multimodal capabilities have led to remarkable advancements in the field of artificial intelligence. One of the key areas where these developments have had a significant impact is the enhancement of multi-modal reasoning through prompting and in-context learning techniques.

In-context learning, the ability of LLMs to learn and adapt to new tasks by simply conditioning on a few demonstration examples, has been a game-changing capability \cite{ref92,ref93}. This concept has now been extended to the multimodal domain, where the integration of visual and textual information has the potential to unlock even more powerful reasoning abilities.

Several recent studies have explored the use of prompting and in-context learning to enhance the multi-modal reasoning capabilities of large language models. One such approach, known as Multimodal-CoT, incorporates both language (text) and vision (images) modalities into a two-stage framework that separates rationale generation and answer inference \cite{ref94}. By leveraging the multimodal information, Multimodal-CoT has been shown to outperform state-of-the-art LLMs on complex reasoning tasks, such as the ScienceQA benchmark, even surpassing human performance.

Another line of research has focused on the role of visual information in the in-context learning process. A study on understanding and improving in-context learning on vision-language models \cite{ref95} has revealed that the textual information in the demonstrations is the primary driver of in-context learning performance, while the visual information has a relatively limited impact. This finding has important implications for the design of effective in-context learning prompts for multimodal tasks, suggesting the need for a careful balance between textual and visual information.

To further enhance the multimodal reasoning capabilities of LLMs, researchers have explored the concept of ``visual chain-of-thought'' prompting, where the model is prompted to generate a series of intermediate reasoning steps that incorporate both visual and textual information \cite{ref94}. By separating the rationale generation and answer inference stages, this approach has been shown to significantly improve the performance of LLMs on various multimodal reasoning tasks.

In addition to prompting techniques, the integration of external tools and knowledge sources has also been explored as a way to boost the multimodal reasoning abilities of LLMs. For instance, the KAM-CoT framework \cite{ref96} combines chain-of-thought reasoning, knowledge graphs, and multiple modalities to provide a comprehensive understanding of multimodal tasks. By incorporating external knowledge during the reasoning process, KAM-CoT has been shown to outperform state-of-the-art methods on the ScienceQA dataset, achieving significantly higher accuracy.

Another approach, known as ToolkenGPT \cite{ref97}, explores the use of tool embeddings to represent a wide range of external tools and enable their integration with LLMs. This method allows LLMs to flexibly leverage a large number of tools, improving their performance on complex multimodal tasks that require diverse capabilities.

The field of enhancing multimodal reasoning through prompting and in-context learning is an active area of research, with many exciting developments on the horizon. Researchers are exploring ways to better understand the interplay between textual and visual information, develop more effective prompting strategies, and leverage external knowledge and tools to further boost the reasoning capabilities of LLMs \cite{ref98,ref99}.

As large language models continue to evolve and incorporate multimodal capabilities, the ability to effectively leverage prompting and in-context learning techniques will be crucial in unlocking their full potential for solving complex, real-world problems. The insights and advancements in this area will undoubtedly pave the way for more intelligent and versatile artificial systems that can seamlessly integrate and reason across multiple modalities.

\subsection{Challenges and Limitations}

The development and deployment of large multi-modal language models (MMMLMs) have brought forth a myriad of challenges and limitations that demand careful consideration. One of the key issues is the inherent data bias present in the training corpora. MMMLMs, like their text-only counterparts, are susceptible to encoding societal biases, stereotypes, and inequities reflected in the data used for pre-training \cite{ref100}. This can lead to the generation of outputs that perpetuate harmful prejudices and discrimination, posing significant risks when these models are applied in high-stakes domains such as healthcare, finance, and education \cite{ref101}.

Addressing data bias is a complex challenge, as it requires not only curating diverse and representative datasets but also developing debiasing techniques that can effectively mitigate the propagation of biases within the model architecture \cite{ref100}. Recent advancements in this area, such as the use of adversarial training, data augmentation, and prompting strategies, have shown promise, but more research is needed to develop robust and scalable solutions \cite{ref100}.

Another significant limitation of MMMLMs is their potential safety and security risks. These models can be vulnerable to adversarial attacks, where malicious actors exploit their weaknesses to generate harmful or misleading content \cite{ref102}. This can range from the generation of explicit or violent content to the spread of misinformation and disinformation, which can have severe societal consequences \cite{ref103}. Additionally, the potential for MMMLMs to be misused for malicious purposes, such as automated text or image generation for fraud or manipulation, poses a significant challenge \cite{ref103}.

To address these safety and security concerns, researchers have proposed various mitigation strategies, including the development of comprehensive evaluation frameworks, adversarial training, and the incorporation of safety constraints during the model's training and deployment \cite{ref104,ref105}. However, the rapid pace of innovation in this field and the complexity of the underlying systems make it challenging to keep up with emerging threats, underscoring the need for ongoing research and collaboration between the research community, industry, and policymakers \cite{ref105}.

Another significant limitation of MMMLMs is their interpretability and transparency. The inner workings of these models, particularly the multimodal integration and reasoning mechanisms, remain largely opaque, making it difficult to understand how they arrive at their outputs \cite{ref106}. This lack of interpretability can hinder the development of trust, the identification of biases and errors, and the ability to hold these models accountable for their decisions \cite{ref106}.

To address this challenge, researchers have explored various techniques for enhancing the interpretability of MMMLMs, such as the use of attention visualization, feature importance analysis, and the development of comprehensive evaluation frameworks \cite{ref106,ref107}. However, the inherent complexity of these models and the diverse range of modalities they integrate pose significant challenges, and more research is needed to develop scalable and effective interpretability solutions.

Finally, the computational and resource requirements of MMMLMs present significant limitations. The training and deployment of these models often require vast computational resources, such as large GPU clusters and extensive storage capacity \cite{ref101}. This can limit the accessibility and scalability of these models, particularly for smaller organizations and resource-constrained environments.

To address these limitations, researchers have explored techniques such as model compression, quantization, and sparse attention mechanisms to improve the efficiency of MMMLMs \cite{ref101}. Additionally, the development of specialized hardware, such as TPUs and dedicated AI accelerators, has the potential to enhance the computational performance of these models \cite{ref101}. However, these solutions often come with their own trade-offs, and further research is needed to strike a balance between model performance, efficiency, and accessibility.

\subsection{Future Outlook and Conclusion}

The rapid advancements in large multi-modal language models (MMMLs) have transformed the landscape of artificial intelligence, ushering in unprecedented capabilities that have captured the imagination of researchers, practitioners, and the general public alike. This survey has provided a comprehensive overview of the key developments, architectural innovations, and applications of these remarkable models, underscoring their immense potential to reshape various domains.

As we peer into the future, the continued evolution of MMMLs holds the promise of even more profound impacts. One of the most exciting prospects is the potential for these models to achieve greater generalization and adaptability, transcending the boundaries of specific tasks and modalities. The emergence of techniques like prompt engineering and in-context learning \cite{ref108} has demonstrated the ability of MMMLs to leverage contextual information and flexibly apply their knowledge to diverse scenarios. Researchers are actively exploring ways to further enhance the multimodal reasoning capabilities of these models, enabling them to seamlessly integrate and reason across visual, auditory, and textual data.

The integration of MMMLMs with other emerging technologies, such as robotics and virtual environments, offers tantalizing opportunities for immersive and embodied interactions \cite{ref6}. By grounding language understanding in real-world physical interactions, these models can develop a deeper comprehension of the world, leading to more natural and intuitive human-AI collaboration. The advancements in areas like simultaneous localization and mapping (SLAM), computer vision, and haptic feedback are paving the way for MMMLs to be deployed in diverse applications, from assistive technologies to remote collaboration.

Moreover, the potential of MMMLs to revolutionize scientific discovery and education cannot be overstated \cite{ref109,ref110}. By leveraging their vast knowledge and reasoning capabilities, these models can assist researchers in formulating novel hypotheses, synthesizing complex information, and generating new insights across a wide range of disciplines. In the educational domain, MMMLs can serve as personalized tutors, providing tailored learning experiences, generating engaging educational content, and empowering students to explore and understand the world in innovative ways.

However, the development and deployment of MMMLs are not without their challenges \cite{ref111}. Addressing issues related to data bias, safety, and interpretability will be crucial to ensure the responsible and ethical use of these models. Researchers are actively exploring techniques to mitigate bias, enhance transparency, and ensure the alignment of MMMLs with human values \cite{ref14}. As these models become more integrated into our daily lives, the need for robust governance frameworks and interdisciplinary collaboration will become increasingly paramount.

Looking ahead, the future of large multi-modal language models holds tremendous promise. The continued advancements in areas such as multimodal fusion, cross-modal reasoning, and language-guided perception \cite{ref15} will pave the way for even more versatile and capable models. The integration of MMMLs with other emerging technologies, such as edge computing and federated learning, can further enhance their scalability, accessibility, and privacy-preserving capabilities.

As we navigate this exciting frontier, the research community must remain vigilant in addressing the ethical and societal implications of these powerful technologies. By fostering responsible innovation, nurturing interdisciplinary collaboration, and prioritizing the well-being of individuals and communities, we can harness the transformative potential of large multi-modal language models to create a future where human-AI collaboration flourishes, scientific progress accelerates, and educational opportunities are democratized. The insights and directions outlined in this comprehensive survey serve as a springboard for the next generation of MMML research and development, ushering in a era of unprecedented possibilities.

\section{Architectural Designs and Training Approaches}

\subsection{Architectural Designs for Vision-Language Integration}

The integration of vision and language models is a critical component in the development of large multi-modal language models. Researchers have explored various architectural designs to effectively combine the representations of visual and textual modalities, with the goal of enabling more robust and comprehensive cross-modal understanding.

These architectural approaches can be broadly categorized into three main frameworks: transformer-based encoders, multi-modal fusion modules, and bridging mechanisms.

The transformer-based encoder approach leverages the powerful self-attention mechanism of the transformer architecture to jointly model the interactions between vision and language. Models like METER \cite{ref112} utilize vision encoders (e.g., CLIP-ViT, Swin Transformer) and text encoders (e.g., RoBERTa, DeBERTa) to extract modality-specific features, which are then fed into a multi-modal fusion module to enable cross-modal interactions. This fully transformer-based design has been shown to outperform previous region-feature-based approaches on various vision-language tasks.

Building upon the transformer-based approach, the BridgeTower architecture \cite{ref113} introduces multiple bridge layers that connect the top layers of the uni-modal encoders with each layer of the cross-modal encoder. This bottom-up cross-modal alignment and fusion strategy has been found to be more effective than simple concatenation or fusion of modality-specific features.

In contrast, the multi-modal fusion module approach focuses on developing specialized mechanisms to integrate the visual and textual representations. All-in-One Image-Grounded Conversational Agents \cite{ref114} combines state-of-the-art transformer and ResNeXt modules, feeding them into a novel attentive multi-modal module to leverage the individual strengths of language and vision models. Similarly, Switch-BERT \cite{ref25} extends the BERT model by introducing learnable layer-wise and cross-layer interactions, allowing the model to optimize attention from a set of attention modes that represent different types of intra-modal and inter-modal interactions.

The bridging mechanism approach, on the other hand, aims to build explicit connections between the visual and textual modalities to facilitate cross-modal understanding. Encoder Fusion Network with Co-Attention Embedding for Referring Image Segmentation \cite{ref115} transforms the visual encoder into a multi-modal feature learning network, using language to refine the multi-modal features progressively. Additionally, a co-attention mechanism is embedded to enable the parallel update of multi-modal features. The Question Aware Vision Transformer for Multimodal Reasoning \cite{ref116} takes a different approach by embedding question awareness directly within the vision encoder, resulting in dynamic visual features that focus on relevant image aspects related to the posed question.

These architectural designs highlight the importance of effectively capturing the interactions and alignments between the visual and textual modalities. The choice of approach often depends on the specific task and dataset, as well as the trade-offs between model complexity, performance, and computational efficiency. By leveraging the strengths of different architectural frameworks, researchers continue to push the boundaries of large multi-modal language models, enabling more robust and comprehensive cross-modal understanding.

\subsection{Training Approaches for Multimodal Pre-training}

The rapid progress in large language models and the emergence of multimodal capabilities have sparked significant interest in developing effective training approaches for pre-training large multi-modal language models. Researchers have explored various objectives and strategies to learn cross-modal representations that can effectively capture the rich semantic associations between different modalities, such as vision and language.

One of the widely adopted training approaches for multimodal pre-training is masked language modeling (MLM) \cite{ref117,ref118,ref119}. In MLM, a portion of the input text is randomly masked, and the model is trained to predict the masked tokens based on the surrounding context and the associated visual information. This objective encourages the model to learn cross-modal alignments between the textual and visual inputs, as the model needs to leverage both modalities to effectively recover the masked tokens.

However, several studies have highlighted the limitations of the standard MLM approach in the multimodal setting. For example, \cite{ref120} found that many masked tokens in the text can be predicted using only the language information, without fully leveraging the visual input. Additionally, \cite{ref118} noted that the standard MLM approach tends to focus more on learning local-to-local alignments, while neglecting the importance of global-to-local and global-to-global alignments between the two modalities.

To address these issues, researchers have proposed various extensions and alternative training approaches for multimodal pre-training. One such approach is to incorporate additional pre-training objectives that explicitly encourage cross-modal alignment at different levels of granularity. For instance, \cite{ref117} introduced a Word-Region Alignment (WRA) task, which uses optimal transport to align individual words in the text with the corresponding image regions, in addition to the standard MLM and Image-Text Matching (ITM) tasks.

Similarly, \cite{ref121} proposed the Global and Local Semantic Completion Learning (GLSCL) task, which combines a Masked Global Semantic Completion (MGSC) objective to learn global-level cross-modal alignments and a Masked Local Token Completion (MLTC) objective to capture fine-grained local-level alignments. By incorporating both global and local semantic alignment objectives, these approaches aim to learn more comprehensive and robust cross-modal representations.

Another line of research has explored the use of contrastive learning as a training objective for multimodal pre-training. \cite{ref122} introduced the Dense Contrastive Visual-Linguistic Pretraining (DCVLP) framework, which replaces the standard region regression and classification tasks with cross-modal region contrastive learning. By learning to align the visual regions and textual tokens in a self-supervised manner, DCVLP can overcome the limitations of noisy or sparse annotations in existing pre-training datasets.

Building on the success of contrastive learning, \cite{ref123} proposed a vision-language pre-training approach specifically tailored for scene text detection tasks. Their model incorporates an image-text contrastive learning (ITC) objective, along with masked language modeling (MLM) and a word-in-image prediction (WIP) task, to learn cross-modal representations that can effectively capture the interactions between text and visual inputs.

Furthermore, \cite{ref124} explored the application of multimodal pre-training to the audio-language domain. They introduced a Cross-modal Transformer for Audio-and-Language (CTAL) model, which is pre-trained on audio-text pairs using a masked language modeling objective and a masked cross-modal acoustic modeling objective, demonstrating the effectiveness of cross-modal pre-training for audio-language tasks.

In addition to the core pre-training objectives, researchers have also explored different strategies to enhance the learning of cross-modal representations. \cite{ref125} proposed the EPIC framework, which introduces a Saliency-based Masking Strategy and an Inconsistent Token Generation Procedure to better align the visual and textual inputs during pre-training.

Another important aspect of multimodal pre-training is the handling of multimodal data, especially in the presence of missing or noisy modalities. \cite{ref33} introduced the LANISTR framework, which can effectively learn from a combination of structured and unstructured data, including language, image, and tabular inputs, by employing a similarity-based multimodal masking loss to handle missing modalities.

As the field of multimodal pre-training continues to evolve, researchers are also exploring ways to leverage the rich semantic information available in external knowledge sources to further enhance the cross-modal representations. \cite{ref126} proposed the ROSITA framework, which integrates cross-modal and intra-modal knowledge from scene graphs to improve the semantic alignments between vision and language during pre-training.

Overall, the research on training approaches for multimodal pre-training has made significant strides in recent years, moving beyond the standard MLM objective and exploring a variety of novel strategies to learn more effective and comprehensive cross-modal representations. By incorporating diverse pre-training tasks, contrastive learning objectives, and external knowledge sources, these approaches aim to bridge the gap between different modalities and enable more robust and generalizable multimodal models for a wide range of downstream applications.

\subsection{Multimodal Dataset Curation and Augmentation}

The rapid proliferation of multimodal data, encompassing text, images, audio, and video, has significantly driven the advancement of large multimodal language models (LMLMs). However, the curation and augmentation of large-scale multimodal datasets remain a critical challenge. Traditional approaches to dataset creation often rely on manual curation and labeling, which can be time-consuming, labor-intensive, and prone to human bias \cite{ref127,ref128}.

To address these limitations, researchers have explored various techniques to automate and enhance the process of multimodal dataset curation and augmentation. One such approach is data synthesis, where artificial data is generated to augment the existing dataset \cite{ref128,ref129}. This not only increases the size of the training dataset but also introduces diversity and reduces the reliance on manual labeling. For example, \cite{ref129} proposed a Progressive Dataset Distillation (PDD) method that synthesizes multiple small sets of synthetic images, each conditioned on the previous sets, to capture the changing training dynamics of deep neural networks.

Another promising technique is dataset distillation, which aims to synthesize a compact dataset that retains the essential information from a large original dataset \cite{ref130,ref131,ref128}. By reducing the size of the dataset while preserving its informative content, dataset distillation can significantly improve the efficiency of model training. For instance, \cite{ref131} proposed a critical samples selection strategy to exploit the most valuable samples during the distillation process, leading to substantial reductions in training costs.

Task-specific data collection is another approach to enhance the diversity and quality of multimodal datasets \cite{ref132}. By carefully curating datasets that align with the specific tasks and applications of interest, researchers can ensure that the training data is well-suited to the target problem. This is particularly important in the context of LMLMs, where the diversity and representativeness of the training data can have a significant impact on the model's performance and generalization capabilities \cite{ref132}.

In addition to these techniques, researchers have also explored the use of data augmentation to further enrich multimodal datasets. Data augmentation involves applying transformations or perturbations to the existing data to create new, synthetic samples \cite{ref133,ref134,ref135}. This not only increases the size of the dataset but also helps to improve the model's robustness and generalization by exposing it to a wider range of variations during training.

For example, \cite{ref134} proposed a method called LatentAugment that leverages generative adversarial networks (GANs) to generate diverse and high-quality synthetic samples for data augmentation. By modifying the latent vectors of the GAN, LatentAugment can produce augmented samples that are both faithful to the original data distribution and highly diverse, leading to improved model performance.

Furthermore, \cite{ref135} demonstrated the effectiveness of using large multimodal models, such as CLIP and BLIP-2, to filter and modify web-crawled data for the DataComp challenge, which focused on improving the quality of training data for multimodal tasks.

The curation and augmentation of multimodal datasets is a critical and evolving area of research, as the success of LMLMs is heavily dependent on the availability of high-quality, diverse, and representative training data. By leveraging techniques such as data synthesis, dataset distillation, task-specific data collection, and data augmentation, researchers are making significant strides in addressing the challenges of multimodal dataset creation and enrichment \cite{ref130,ref136}.

As the field of LMLMs continues to advance, the importance of effective multimodal dataset curation and augmentation will only grow. Researchers and practitioners must continue to explore innovative methods and strategies to ensure that the training data used to develop these powerful models is as comprehensive, representative, and high-quality as possible. This will be crucial for unlocking the full potential of LMLMs and enabling their successful deployment across a wide range of applications.

\subsection{Modality-Aware Model Design and Parameter Sharing}

Designing effective architectures for multimodal learning is a crucial challenge, as the interactions between different modalities can profoundly impact the overall performance. To capture these intricate relationships, researchers have explored various strategies for modality-aware model design and parameter sharing across vision and language components.

One key approach is the use of modality-specific modules within the network architecture. \cite{ref25} proposes the Switch-BERT model, which extends the BERT architecture by introducing learnable layer-wise and cross-layer interactions. This modality-specific design allows the model to better capture the complex relationships between different modalities. Similarly, \cite{ref137} introduces the concept of modality heterogeneity and interaction heterogeneity to guide the design of multimodal architectures, proposing the High-Modality Multimodal Transformer (HighMMT) to leverage these insights.

Another approach focuses on leveraging the benefits of parameter sharing across modalities. \cite{ref138} presents the UNIMO-3 model, which aims to learn effective cross-modal interactions by establishing connections between different layers of the vision and language encoders. This parameter sharing strategy enables the model to better integrate information from multiple modalities. \cite{ref139} introduces the BEiT-3 model, which utilizes a Multiway Transformer architecture to enable both deep fusion and modality-specific encoding, further demonstrating the importance of modality-aware design.

The concept of modality-aware design is further explored in \cite{ref27}, which presents the FLAVA model, a holistic foundation model that targets all modalities at once. By considering the interplay between different modalities from the ground up, FLAVA aims to learn more comprehensive multimodal representations. In addition, \cite{ref140} introduces the mPLUG-2 model, which employs a multi-module composition network to allow for modality collaboration while also disentangling different modality modules, showcasing the potential of modular architectures for multimodal learning.

Furthermore, \cite{ref141} presents a novel modality interaction strategy for 3D object detection, where individual per-modality representations are learned and maintained throughout the model. This modality-aware design highlights the importance of effectively integrating information from different modalities for task-specific applications. The significance of modality-aware design and parameter sharing is also emphasized in the context of emerging multimodal applications, such as \cite{ref17}.

Overall, the research efforts in modality-aware model design and parameter sharing across vision and language components have demonstrated the potential to enhance the learning of effective multimodal representations, enabling improved performance on a wide range of multimodal tasks. These approaches highlight the importance of carefully considering the interplay between different modalities in the design of large multimodal language models.

\subsection{Efficient Multi-Modal Reasoning and Inference}

The use of prompting and in-context learning techniques has emerged as a promising approach for enhancing the multi-modal reasoning and inference capabilities of large language models (LLMs). These techniques leverage the ability of LLMs to rapidly adapt to new tasks and contexts using a few demonstration examples, rather than relying on fine-tuning or other resource-intensive adaptation methods.

One key aspect of this line of research is the exploration of how prompts and in-context learning can be effectively applied to multi-modal tasks that involve both language and visual information. Recent studies have demonstrated the potential of these techniques in improving the performance of LLMs on complex vision-language tasks, such as visual question answering, image captioning, and multi-modal reasoning \cite{ref142,ref95}.

For example, the MMICL approach \cite{ref142} introduces a novel context scheme that augments the in-context learning ability of vision-language models (VLMs) by incorporating both language and visual information into the prompts. This enables the model to better understand and reason about complex multi-modal inputs, leading to improved performance on a wide range of vision-language tasks. Similarly, the work on understanding and improving in-context learning on VLMs \cite{ref95} highlights the importance of carefully selecting effective multi-modal demonstrations to enhance the in-context learning capabilities of these models.

Another promising direction in this area is the use of task-specific adapters to improve the multi-modal reasoning and inference capabilities of LLMs. Researchers have explored the idea of designing modular architectures that can effectively combine and leverage the strengths of pre-trained vision and language models \cite{ref143,ref6}.

For instance, the work on mastering robot manipulation with multimodal prompts \cite{ref143} presents a framework that incorporates a multi-modal prompt encoder and a transformer-based mapping network to enable embodied robots to understand and execute complex instructions that involve both language and visual information. Similarly, the PaLM-E model \cite{ref6} introduces an embodied multimodal language model that directly incorporates real-world continuous sensor modalities, including vision, into the language model architecture, enabling it to perform a variety of embodied reasoning and manipulation tasks.

These task-specific adapters and modular architectures offer several advantages over traditional fine-tuning approaches. By separating the reasoning and recognition components of the model, they can better leverage the strengths of pre-trained vision and language models, while also allowing for more efficient and targeted adaptation to specific multi-modal tasks and domains \cite{ref144}. Furthermore, the use of adapters and modular designs can help mitigate the risk of catastrophic forgetting, where fine-tuning a model on a new task can lead to a loss of performance on previously learned tasks.

In addition to prompting and in-context learning, researchers have also explored the use of multi-modal chain-of-thought reasoning to enhance the multi-modal inference capabilities of LLMs \cite{ref145,ref94}. These approaches aim to emulate the step-by-step reasoning process that humans often employ when solving complex, multi-step problems that involve both language and visual information.

For example, the DDCoT method \cite{ref145} introduces a novel prompting technique that maintains a critical attitude through negative-space prompting and incorporates multi-modality into the reasoning process by first dividing the reasoning responsibility of LLMs into reasoning and recognition, and then integrating the visual recognition capability of visual models into the joint reasoning process. This approach has been shown to improve the reasoning abilities of both large and small language models on a variety of multi-modal tasks.

Similarly, the Multimodal-CoT framework \cite{ref94} proposes a two-stage approach that separates rationale generation and answer inference, allowing the model to leverage better-generated rationales that are based on multi-modal information. This approach has been demonstrated to outperform state-of-the-art LLMs on complex multi-modal reasoning tasks, and even surpass human performance on certain benchmarks.

Overall, the research on enhancing the multi-modal reasoning and inference capabilities of large language models through the use of prompting, in-context learning, and task-specific adapters highlights the importance of leveraging the complementary strengths of vision and language models to tackle complex, real-world problems that involve multiple modalities. As this field continues to evolve, we can expect to see further advancements in the development of efficient and effective multi-modal reasoning systems that can seamlessly integrate and reason about language, vision, and other modalities.

\section{Benchmarking and Evaluation Frameworks}

\subsection{Existing Benchmark Datasets and Evaluation Metrics}

The evaluation of large multi-modal language models (MLLMs) has been a topic of growing interest in the research community, with the development of various benchmark datasets and evaluation metrics. These benchmarks and metrics aim to provide a comprehensive and standardized framework for assessing the capabilities of MLLMs in diverse multimodal tasks and scenarios.

One of the pioneering efforts in this direction is the MM-BigBench \cite{ref51}, which introduces a comprehensive assessment framework for evaluating the performance of MLLMs on multimodal content comprehension tasks. The MM-BigBench dataset encompasses a wide range of metrics, including the Best Performance metric to determine the upper bound of each model's performance, the Mean Relative Gain metric to assess the overall performance, the Stability metric to measure the models' sensitivity, and the Adaptability metric to quantify the adaptability between models and instructions. By employing these diverse evaluation methods, the MM-BigBench framework provides a more holistic and nuanced understanding of the capabilities of MLLMs in understanding and reasoning about multimodal content.

Another prominent benchmark for evaluating MLLMs is SEED-Bench \cite{ref146}, which focuses on assessing the hierarchical capabilities of MLLMs, ranging from the ability to accept and generate text to the ability to accept and generate both text and images. SEED-Bench comprises 24,000 multiple-choice questions with accurate human annotations, covering 27 dimensions that evaluate both text and image generation. The use of multiple-choice questions with ground-truth options derived from human annotations enables an objective and efficient assessment of model performance, eliminating the need for human or GPT intervention during evaluation.

In addition to these comprehensive multimodal benchmarks, researchers have also proposed specialized benchmarks to evaluate the performance of MLLMs in specific domains or tasks. For instance, MMBench \cite{ref55} focuses on assessing the various vision-language capabilities of MLLMs, including spatial and temporal understanding, reasoning, and generation. The benchmark incorporates a meticulously curated dataset and a novel CircularEval strategy that leverages ChatGPT to convert free-form predictions into pre-defined choices, facilitating a more robust evaluation of the model's predictions.

Furthermore, the ChEF \cite{ref54} framework proposes a modular and standardized approach to evaluating MLLMs, allowing for the creation of diverse evaluation scenarios by combining different components, such as multimodal datasets, instruction retrieval formulas, and evaluation metrics. By providing a unified and extensible evaluation framework, ChEF aims to facilitate a more systematic and comprehensive assessment of MLLMs' capabilities.

While these benchmark datasets and evaluation frameworks have made significant contributions to the field, they also face various limitations and challenges. One of the primary concerns is the potential for dataset bias and the need to ensure the diversity and representativeness of the evaluation data \cite{ref147}. Many of the existing benchmarks focus on a limited set of languages or domains, which may not accurately reflect the real-world performance of MLLMs in more diverse and complex scenarios.

Another challenge is the development of robust and reliable evaluation metrics that can effectively capture the nuances of multimodal understanding and reasoning \cite{ref148}. The traditional metrics used in NLP, such as BLEU or ROUGE, may not be well-suited for evaluating the performance of MLLMs, as they often fail to capture the semantic and contextual aspects of multimodal information processing. There is a need for the development of more sophisticated evaluation metrics that can better align with human judgments and provide a more comprehensive assessment of the models' capabilities.

Additionally, the rapid advancements in MLLM technology have outpaced the development of evaluation benchmarks, leading to concerns about the relevance and scalability of existing datasets and metrics \cite{ref60}. There is a growing need for flexible and adaptable evaluation frameworks that can keep pace with the evolving capabilities of MLLMs and provide meaningful insights into their performance across diverse scenarios.

To address these challenges, researchers have explored the use of large language models (LLMs) themselves as evaluators, leveraging their natural language understanding and generation capabilities to assess the performance of other MLLMs \cite{ref149}. However, the reliability and alignment of LLM-based evaluators with human judgments remain a subject of ongoing investigation \cite{ref60,ref150}.

In conclusion, the existing benchmark datasets and evaluation metrics for large multi-modal language models have made significant strides in providing a more comprehensive and standardized framework for assessing the capabilities of these models. However, ongoing challenges and limitations, such as dataset bias, the need for more robust evaluation metrics, and the scalability of evaluation approaches, highlight the importance of continued research and innovation in this area. As the field of multi-modal language modeling continues to evolve, the development of flexible and adaptable evaluation frameworks will be crucial for driving the advancement of this technology and ensuring its reliable and beneficial deployment in real-world applications.

\subsection{Comprehensive Benchmarking Frameworks}

Existing benchmark datasets and evaluation protocols have played a crucial role in the rapid advancement of large multi-modal language models (LMLMs). However, many of these benchmarks have been criticized for their narrow focus on specific tasks or modalities, failing to provide a comprehensive assessment of the models' overall capabilities \cite{ref54,ref151}.

To address this limitation, researchers have proposed the development of more holistic and integrated benchmarking frameworks that can evaluate the models' multi-modal reasoning, language understanding, and generalization capabilities \cite{ref152,ref55}. These comprehensive benchmarking frameworks aim to provide a more thorough and rigorous assessment of the models' performance, going beyond traditional task-specific evaluations.

One such example is the ConTextual benchmark, which focuses on evaluating the context-sensitive text-rich visual reasoning capabilities of LMLMs \cite{ref152}. The benchmark consists of instructions designed to assess the models' ability to perform tasks that require a deeper understanding of the interactions between textual and visual elements, such as time-reading, navigation, and shopping. By emphasizing diverse real-world scenarios, ConTextual aims to provide a more realistic and challenging evaluation of the models' multi-modal reasoning abilities.

Another comprehensive benchmarking framework is MMBench, which is designed to robustly evaluate the various abilities of vision-language models \cite{ref55}. MMBench employs a meticulously curated dataset that surpasses existing benchmarks in terms of the number and variety of evaluation questions and abilities. Additionally, it introduces a novel CircularEval strategy that incorporates the use of ChatGPT to convert free-form predictions into pre-defined choices, enabling a more robust evaluation of the models' performance.

The MLLM-Bench benchmark, on the other hand, focuses on evaluating the performance of multimodal large language models (MLLMs) across a wide range of scenarios, including perception, understanding, applying, analyzing, evaluating, and creation, as well as ethical considerations \cite{ref151}. The benchmark is designed to reflect user experience more accurately and provide a more holistic assessment of model performance, addressing the limitations of existing benchmarks that rely on objective queries with standard answers.

Similarly, the SEED-Bench benchmark aims to comprehensively evaluate the generative comprehension capabilities of MLLMs \cite{ref46}. The benchmark consists of 19,000 multiple-choice questions with accurate human annotations, spanning 12 evaluation dimensions that cover both the spatial and temporal understanding of the models.

Another notable comprehensive benchmarking framework is ChEF, which provides a standardized and holistic evaluation of MLLMs \cite{ref54}. ChEF is structured as four modular components: Scenario (scalable multimodal datasets), Instruction (flexible instruction retrieving formulae), Inferencer (reliable question answering strategies), and Metric (indicative task-specific score functions). This modular design allows for the creation of new evaluation recipes by systematically selecting these components, enabling a flexible and comprehensive assessment of MLLMs.

These comprehensive benchmarking frameworks represent a significant step forward in the evaluation of large multi-modal language models, as they move beyond traditional task-specific assessments and provide a more holistic and rigorous evaluation of the models' capabilities. By employing diverse and challenging scenarios, these benchmarks aim to uncover the true strengths and limitations of LMLMs, driving the development of more robust and generalized models that can effectively handle complex multi-modal tasks.

Moreover, these comprehensive benchmarking frameworks often include advanced evaluation metrics and techniques, such as error localization capabilities, to enhance the reliability and interpretability of the assessment process. This, in turn, can facilitate the development of more trustworthy and transparent multi-modal language models that can be safely deployed in real-world applications.

Overall, the emergence of comprehensive benchmarking frameworks for large multi-modal language models represents a crucial advancement in the field, as it paves the way for a more systematic and rigorous evaluation of these powerful AI systems. By challenging the models with diverse and complex tasks, these benchmarks can help researchers and developers identify the key areas for improvement and guide the development of more capable and reliable multi-modal language models.

\subsection{Benchmarking Across Languages and Domains}

While significant progress has been made in the evaluation of large multi-modal language models (LMLLs), the majority of existing benchmarks and evaluation frameworks have been primarily focused on a limited set of languages, predominantly English. This lack of linguistic diversity in MMLL evaluation has become a growing concern, as it fails to capture the true capability and generalizability of these models across the vast linguistic landscape.

The rapid advancements in LMLLs have underscored the need for more comprehensive and inclusive benchmarking efforts that span a wide range of languages and domains \cite{ref153,ref154}. Currently, the majority of MMLL benchmarks are confined to a handful of high-resource languages, leaving a significant portion of the world's linguistic diversity unaccounted for. This narrow focus not only limits our understanding of MMLL performance but also raises concerns about the equitable development and deployment of these models, particularly in underserved linguistic communities \cite{ref155,ref62}.

To address this pressing issue, researchers have begun to explore the development of language-specific and domain-specific evaluation frameworks for LMLLs \cite{ref156,ref157}. These efforts aim to capture the nuances and unique characteristics of different languages and domains, providing a more comprehensive assessment of MMLL capabilities.

One of the key challenges in developing language-specific benchmarks lies in the availability and quality of evaluation data \cite{ref66}. Many low-resource and minority languages lack the necessary corpora and datasets required for rigorous MMLL evaluation. Researchers have addressed this challenge by leveraging crowdsourcing, machine translation, and cross-lingual transfer learning techniques to expand the linguistic coverage of existing benchmarks \cite{ref153,ref62}.

Domain-specific benchmarking presents an additional layer of complexity, as LMLLs need to demonstrate not only linguistic proficiency but also the ability to understand and reason about domain-specific knowledge and concepts \cite{ref158,ref159}. Researchers have addressed this challenge by curating datasets that capture the specialized vocabulary, reasoning patterns, and task requirements of various domains, such as healthcare, finance, and transportation.

Furthermore, the development of cross-lingual and cross-domain evaluation frameworks has emerged as a crucial research direction, as it allows for the assessment of MMLL performance in more realistic and dynamic settings \cite{ref160,ref161}. These frameworks aim to gauge the ability of LMLLs to adapt and perform well across diverse linguistic and contextual boundaries, providing a more comprehensive understanding of their overall capabilities and limitations.

The challenge of benchmarking LMLLs across languages and domains is further exacerbated by the inherent complexity and heterogeneity of these models. As the scale and architectural sophistication of LMLLs continue to evolve, the task of designing appropriate evaluation frameworks becomes increasingly challenging \cite{ref65}. Researchers have explored innovative approaches, such as the use of agent-based evaluation \cite{ref162} and the integration of human-centric assessments \cite{ref163}, to capture the nuanced and multifaceted capabilities of LMLLs.

In conclusion, the need for comprehensive and inclusive benchmarking of large multi-modal language models across diverse languages and domains has become increasingly apparent. While significant progress has been made in this direction, the research community continues to face numerous challenges, including data availability, domain-specific knowledge capture, and the inherent complexity of these models. Addressing these challenges will be crucial in ensuring the responsible development and deployment of LMLLs, ultimately enabling their equitable access and impact across the global linguistic landscape.

\subsection{Automated and Human-Centric Evaluation Approaches}

Automated and human-centric evaluation approaches have both advantages and limitations when it comes to assessing the performance of large multi-modal language models (LMLMs). Automated metrics, such as those based on text overlap (e.g., BLEU, ROUGE) or language model judgments, offer a cost-effective and scalable way to evaluate model performance \cite{ref60}. However, these automated methods have been shown to exhibit biases and poor alignment with human judgments, particularly for open-ended and creative tasks \cite{ref149,ref164}.

In contrast, human-centric evaluation approaches, where human annotators provide subjective ratings or rankings of model outputs, can capture more nuanced and holistic aspects of model performance \cite{ref165}. These human evaluations, however, are often time-consuming, expensive, and prone to inconsistencies across different annotators \cite{ref166}.

To address the limitations of both automated and human-centric approaches, researchers have explored the potential for hybrid evaluation methods that leverage the strengths of both \cite{ref167,ref168}. These hybrid approaches may involve using LLMs as initial evaluators, with human annotators providing feedback and refinement, or incorporating both automated and human-centric metrics in a comprehensive evaluation framework \cite{ref169}.

One promising hybrid approach is the use of LLMs as ``evaluator-in-the-loop'' during the evaluation process \cite{ref167}. In this paradigm, LLMs are not only the subjects of evaluation but also active participants in the evaluation process, generating initial assessments or feedback that can then be refined and validated by human annotators. This approach has been shown to enhance the efficiency of human evaluations by prioritizing the most informative data instances and reducing the number of required annotations \cite{ref170}.

Another hybrid approach involves the hierarchical decomposition of evaluation criteria, where LLMs are used to generate and refine a comprehensive set of evaluation dimensions, which are then aggregated and weighted according to human preferences \cite{ref168}. This method aims to capture the multi-faceted nature of model performance while maintaining alignment with human judgments.

In addition to these hybrid approaches, researchers have also explored the use of ``meta-evaluation'' frameworks, where LLMs are evaluated on their ability to serve as evaluators for other models \cite{ref60}. These meta-evaluation approaches can help identify the strengths, limitations, and biases of LLMs as evaluators, which is crucial for ensuring the reliability and trustworthiness of their assessments.

Furthermore, there have been efforts to develop more comprehensive and standardized evaluation frameworks that incorporate both automated and human-centric metrics \cite{ref150,ref171}. These frameworks often cover a diverse range of tasks and modalities, aiming to provide a holistic assessment of model capabilities and align with real-world applications.

Despite the progress in this area, challenges remain in developing robust and reliable evaluation methods for large multi-modal language models. Factors such as data bias, model scale, and task complexity can all impact the performance and alignment of both automated and human-centric evaluation approaches \cite{ref172,ref173}.

Addressing these challenges will require continued collaboration between researchers, industry practitioners, and end-users to establish best practices, benchmarks, and evaluation frameworks that can accurately capture the capabilities and limitations of large multi-modal language models. Hybrid approaches that leverage the strengths of both automated and human-centric evaluation methods, while addressing their limitations, hold great promise in advancing the field of LMLM assessment and, ultimately, driving the development of more capable and trustworthy multi-modal AI systems.

\subsection{Evaluating Robustness and Scalability}

Evaluating the robustness and scalability of large multi-modal language models (LMLMs) is crucial as these models continue to gain prominence in various applications. Robustness refers to the ability of these models to maintain their performance in the face of perturbations or distributional shifts, while scalability examines their capacity to handle increasing input/output complexity and maintain consistent performance as the model size grows.

Assessing the robustness of LMLMs is a multi-faceted challenge. One important aspect is their resilience to adversarial attacks, where carefully crafted inputs are designed to mislead the model and generate incorrect or undesired outputs \cite{ref174,ref102}. Researchers have explored various adversarial attack methods, such as gradient-based attacks, that can effectively exploit vulnerabilities in these models \cite{ref175,ref176}. Understanding the mechanisms behind these attacks and developing robust defenses is crucial to ensuring the reliability and safety of LMLMs in real-world applications.

Another important dimension of robustness evaluation is the models' performance under distributional shifts, where the input data deviates from the training distribution \cite{ref177}. This can happen due to changes in the language, domain, or modality of the input, and it is essential to assess how well the models can generalize and maintain their performance in such scenarios. Benchmarks like AdvGLUE \cite{ref178} have been proposed to systematically evaluate the robustness of language models to various types of distributional shifts.

In addition to robustness, the scalability of LMLMs is another crucial aspect to consider. As these models grow in size and complexity, it is important to evaluate their performance under increasing input/output complexity, such as longer sequences, higher-dimensional visual inputs, or more complex multi-modal interactions \cite{ref179}. Metrics like matrix entropy \cite{ref180} can provide insights into the models' ability to extract and compress relevant information as they scale, while measures of computational efficiency, such as the trade-off between model size and performance, can inform the practical deployment of these models \cite{ref181}.

The challenge of evaluating the robustness and scalability of LMLMs is further compounded by the evolving nature of these models and the need for continuous monitoring and assessment. As new architectures, training strategies, and applications emerge, the evaluation approaches must also adapt to capture the changing landscape of this rapidly advancing field \cite{ref182}.

\subsection{Interpretability and Explainability}

Interpretability and explainability are crucial aspects in the evaluation of large multi-modal language models (LMLMs), as they provide insights into the decision-making processes of these complex systems and help identify potential biases. The development of benchmarks and evaluation frameworks that effectively assess the interpretability and explainability of LMLMs remains a significant challenge in the field.

One key challenge in evaluating the interpretability of LMLMs is the inherent complexity of these models, which often operate as ``black boxes'' with intricate internal mechanisms. Existing approaches, such as attention visualization and feature importance analysis, have provided some level of insight, but they tend to focus on specific local explanations or global patterns, without offering a comprehensive understanding of the models' decision-making processes \cite{ref29,ref183}.

To address these limitations, researchers have explored more holistic and multi-dimensional approaches to interpreting LMLMs. For example, the Sparsity-Guided Holistic Explanation (SparseCBM) framework \cite{ref183} integrates sparsity-based techniques to elucidate the models' behaviors at the input, subnetwork, and concept levels, providing a more comprehensive understanding of their inner workings. Similarly, the CLEAR framework \cite{ref184} equips LMLMs with self-awareness and self-correction capabilities, enabling a more transparent and accountable decision-making process.

Another challenge in evaluating the interpretability and explainability of LMLMs is the lack of comprehensive and standardized benchmarks. Existing evaluation frameworks often focus on specific tasks or domains, limiting their ability to capture the full range of LMLM capabilities. To address this, researchers have proposed more diverse and holistic benchmarking approaches, such as the MM-Vet framework \cite{ref185}, which examines the integration of different core vision-language capabilities in LMLMs.

In addition to the technical challenges, the assessment of interpretability and explainability also requires consideration of the users' perspectives and needs. Different stakeholders, such as domain experts, policymakers, and the general public, may have varying requirements for the level of interpretability and explainability of LMLMs, depending on the application context. Developing evaluation frameworks that can cater to these diverse needs is crucial for ensuring the responsible and trustworthy deployment of these models \cite{ref186}.

Furthermore, the interpretability and explainability of LMLMs are closely tied to concerns about bias and fairness. As these models are trained on large and potentially biased datasets, they may exhibit and propagate societal biases, which can have significant implications in high-stakes applications \cite{ref187}. Developing benchmarks and evaluation frameworks that can effectively assess and mitigate these biases is a pressing research challenge.

To address these issues, researchers have proposed various approaches, such as the use of concept-based analysis \cite{ref183}, which can help identify and address biases by examining the models' underlying conceptual representations. Additionally, the incorporation of human-centric evaluation methods, such as the use of diverse user groups and qualitative assessments, can provide valuable insights into the interpretability and explainability of LMLMs from multiple perspectives \cite{ref186}.

In conclusion, the assessment of interpretability and explainability is a critical aspect of evaluating large multi-modal language models. Addressing the technical and user-centric challenges in this area requires a multi-faceted approach, including the development of more comprehensive benchmarking frameworks, the integration of holistic and multi-dimensional interpretation methods, and the consideration of diverse stakeholder needs and concerns. By prioritizing interpretability and explainability in the evaluation of LMLMs, the research community can ensure the responsible and trustworthy deployment of these powerful models in real-world applications.

\section{Applications and Case Studies}

\subsection{Healthcare Applications}

The rapid advancements in large multi-modal language models have unlocked a plethora of opportunities for transforming the healthcare sector. These models have demonstrated exceptional capabilities in clinical language understanding, medical diagnostic assistance, and patient care, positioning them as invaluable tools in the pursuit of enhancing healthcare delivery.

One of the key applications of large multi-modal language models in the healthcare domain is clinical language understanding. These models have shown remarkable proficiency in comprehending and interpreting the nuanced and specialized language used in clinical settings \cite{ref188,ref189,ref190}. They can effectively tackle a wide range of tasks, including named entity recognition, relation extraction, natural language inference, and question-answering, which are crucial for extracting valuable insights from clinical notes, patient records, and medical literature.

The ability of large multi-modal language models to understand clinical terminology and reasoning has also made them instrumental in providing medical diagnostic assistance. By integrating visual and textual information from medical images, patient records, and other relevant data, these models can support clinicians in making more accurate and informed diagnoses \cite{ref191,ref192,ref193}.

Moreover, the potential of large multi-modal language models extends beyond clinical settings, as they can also play a crucial role in enhancing patient care and engagement. These models can provide personalized and accessible information to patients, improving their understanding of their medical conditions, treatment options, and aftercare instructions \cite{ref194,ref195,ref196}.

The integration of large multi-modal language models in healthcare has also shown promise in streamlining clinical workflows and supporting research efforts. These models can automate various administrative tasks, such as clinical documentation, coding, and billing, freeing up healthcare professionals to dedicate more time to patient care \cite{ref70}. Additionally, large multi-modal language models can assist in the synthesis and analysis of vast amounts of medical literature, accelerating the pace of scientific discoveries and the translation of research findings into clinical practice \cite{ref197}.

However, the deployment of large multi-modal language models in healthcare is not without its challenges. Ensuring the accuracy, reliability, and safety of these models is paramount, as any errors or inconsistencies can have serious consequences for patient outcomes \cite{ref198}. Rigorous evaluation frameworks, such as those proposed in \cite{ref188} and \cite{ref193}, are crucial for assessing the performance and limitations of these models in various healthcare scenarios.

Furthermore, the integration of large multi-modal language models into healthcare workflows requires careful consideration of ethical and privacy concerns. These models must be designed and deployed in a manner that respects patient confidentiality, mitigates biases, and promotes equitable access to healthcare services \cite{ref14}. Ongoing research and collaboration between healthcare professionals, researchers, and policymakers are essential for addressing these challenges and ensuring the responsible and effective use of large multi-modal language models in the medical domain.

In conclusion, the emergence of large multi-modal language models has opened up new frontiers in healthcare, transforming clinical language understanding, medical diagnostic assistance, and patient care. These models have demonstrated immense potential, but their successful integration into the healthcare ecosystem requires a comprehensive and multifaceted approach that addresses both the technical and ethical considerations. By leveraging the capabilities of large multi-modal language models while prioritizing patient safety, privacy, and equitable access, the healthcare sector can harness the transformative power of this technology to enhance patient outcomes and drive innovation in the field.

\subsection{Finance Applications}

The finance industry has been one of the early adopters of large multi-modal language models (LMLMs), recognizing their transformative potential across various financial operations. These powerful models have found diverse applications, from enhancing investment analysis to providing personalized financial advice.

In the realm of investment analysis, LMLMs have demonstrated remarkable capabilities in extracting valuable insights from vast troves of financial data, including textual information from earnings reports, market news, and research reports. \cite{ref199} By leveraging their natural language understanding and generation abilities, LMLMs can distill key insights, identify patterns, and generate summaries that aid investment professionals in making informed decisions.

One prominent application of LMLMs in investment analysis is sentiment analysis. Researchers have explored the use of LMLMs, such as ChatGPT and GPT-4, in analyzing the sentiment of financial text, including news articles and social media posts, to gain a deeper understanding of market trends and investor sentiment. \cite{ref200} These models have exhibited superior performance in sentiment classification compared to traditional methods, showcasing their ability to capture the nuances and complexities inherent in financial language.

Beyond sentiment analysis, LMLMs have also been leveraged for portfolio optimization and risk management. By integrating LMLMs with quantitative models, researchers have developed frameworks that can analyze multi-modal data, including textual information and numerical data, to generate investment recommendations and assess potential risks. \cite{ref201} These advancements have the potential to enhance the overall efficiency and accuracy of investment decision-making processes.

In the realm of financial reporting, LMLMs have been employed to automate the generation of financial documents, such as earnings reports and analyst notes. \cite{ref202} By leveraging their natural language generation capabilities, these models can produce coherent and structured financial reports, freeing up valuable time for financial professionals to focus on higher-level strategic analysis and decision-making.

Moreover, LMLMs have been explored for their potential in personalized financial advice. By integrating these models with customer data and financial knowledge, financial institutions can develop intelligent conversational agents that can provide tailored recommendations and guidance to individual clients. \cite{ref203} These personalized financial assistants have the ability to engage in natural language dialogue, understand the specific needs and goals of each client, and offer customized financial advice.

The integration of LMLMs in the finance industry extends beyond the applications mentioned above. Researchers have also explored the use of these models in tasks such as financial document summarization, financial question-answering, and financial risk assessment. \cite{ref204} These advancements have the potential to revolutionize the way financial institutions operate, driving increased efficiency, improved decision-making, and enhanced customer experiences.

However, the adoption of LMLMs in the finance industry is not without its challenges. Concerns around data privacy, algorithmic bias, and the reliability of model outputs have emerged as critical considerations. \cite{ref205} Addressing these challenges requires the development of robust governance frameworks, ethical guidelines, and rigorous evaluation mechanisms to ensure the responsible and trustworthy deployment of these powerful models.

Despite these challenges, the finance industry's embrace of LMLMs reflects a broader trend towards the integration of advanced artificial intelligence and natural language processing technologies. As these models continue to evolve and their capabilities expand, the finance industry is poised to experience significant transformations, unlocking new opportunities for improved decision-making, enhanced customer experiences, and increased operational efficiency.

\subsection{Education Applications}

The role of large multi-modal language models (LMLMs) in transforming science, technology, engineering, and mathematics (STEM) education is a rapidly emerging research area with significant potential. These models have demonstrated impressive capabilities in a wide range of tasks, including personalized learning, interactive visualization, and concept comprehension, which can greatly enhance STEM education.

One of the key applications of LMLMs in STEM education is personalized learning. By understanding the student's prior knowledge, cognitive abilities, and learning preferences, these models can create adaptive learning experiences that dynamically adjust the content, pace, and delivery of educational materials to optimize learning outcomes \cite{ref206}. This personalized approach can be particularly beneficial in STEM subjects, where students often face challenges due to the inherent complexity and abstract nature of the concepts.

LMLMs can also play a crucial role in enhancing interactive visualization for STEM education. By combining the language understanding and generation capabilities of these models with advanced visualization techniques, educators can create immersive learning experiences that facilitate deeper understanding and retention of STEM knowledge \cite{ref73}. Interactive simulations, animations, and graphs generated by LMLMs can help students better comprehend and engage with complex STEM concepts.

Moreover, LMLMs have the potential to revolutionize the way STEM concepts are taught and understood. These models can be leveraged to generate personalized explanations, analogies, and examples that resonate with individual students, thereby enhancing their comprehension of complex STEM topics \cite{ref74}. By understanding the strengths and limitations of LMLMs, educators can strategically integrate these models into their teaching practices to create more engaging, effective, and inclusive STEM learning environments.

One particular example of how LMLMs can benefit STEM education is in the field of mathematics. Researchers have proposed the use of large language models to generate step-by-step explanations and interactive visualizations for mathematical concepts, enabling students to better understand the underlying reasoning and problem-solving processes \cite{ref207}. This approach can be especially valuable in addressing the common challenge of math anxiety, which often stems from the perceived difficulty and abstraction of mathematical concepts \cite{ref208}.

Similarly, in the realm of science education, LMLMs can be leveraged to create interactive simulations and virtual experiments, allowing students to explore and interact with scientific phenomena in a safe and engaging environment \cite{ref74}. These models can also be used to generate personalized explanations and examples that bridge the gap between scientific concepts and real-world applications, fostering a deeper understanding and appreciation of STEM disciplines.

Furthermore, the integration of LMLMs in STEM education has the potential to enhance accessibility and inclusivity. By providing personalized learning experiences and interactive visualizations, these models can cater to students with diverse learning needs, including those with disabilities or language barriers \cite{ref209}. This can help to ensure that STEM education is accessible and engaging for all students, regardless of their background or individual characteristics.

However, the integration of LMLMs in STEM education is not without its challenges. Researchers and educators must carefully navigate issues related to data bias, privacy, and ethical considerations to ensure that these models are developed and deployed in a responsible and equitable manner \cite{ref14}. Additionally, there is a need for further research to understand the long-term impact of LMLMs on student learning, engagement, and academic outcomes, as well as their potential to address systemic barriers and inequities in STEM education.

Despite these challenges, the potential of LMLMs to transform STEM education is undeniable. By leveraging the capabilities of these models to personalize learning, enhance interactive visualization, and improve concept comprehension, educators can create more engaging, effective, and inclusive STEM learning environments that inspire and empower the next generation of scientists, engineers, and innovators. As the field of STEM education continues to evolve, the integration of LMLMs will likely play a crucial role in shaping the future of teaching and learning in these critical domains.

\subsection{Public Health and Wellness Applications}

The rapidly evolving field of large multi-modal language models (LMLMs) holds immense potential for addressing critical public health challenges and promoting overall wellness. These versatile models can serve as powerful tools in a wide range of applications, from disease monitoring and outbreak prediction to personalized wellness coaching and the dissemination of vital health information.

One area where LMLMs have demonstrated significant promise is in disease monitoring and outbreak prediction. By integrating and analyzing data from various sources, such as social media, news reports, and real-time surveillance systems, these models can identify early warning signs of emerging health threats and facilitate timely intervention \cite{ref210}. The PandemicLLM framework \cite{ref210}, for instance, has shown the ability to effectively utilize textual public health policies, genomic surveillance, spatial, and epidemiological time series data to generate accurate and timely predictions of disease spread, outperforming traditional forecasting models. This capability can be invaluable in rapidly evolving public health crises, enabling policymakers and healthcare professionals to allocate resources more effectively and implement targeted interventions.

Moreover, LMLMs can play a crucial role in the dissemination of vital health information, particularly in underserved and hard-to-reach communities. By seamlessly integrating text, images, audio, and video modalities, these models can deliver tailored and culturally-sensitive health education and awareness campaigns, addressing language barriers and fostering greater engagement with the target populations \cite{ref211}. The GenSpectrum Chat \cite{ref211} chatbot, for example, leverages the capabilities of GPT-4 to enable interactive exploration of SARS-CoV-2 genomic data in multiple languages, making critical public health insights accessible to a diverse global audience.

In the realm of personalized wellness coaching, LMLMs hold immense potential. By incorporating data from wearable devices, biometric sensors, and user-provided information, these models can offer customized guidance and support for individuals seeking to improve their overall health and well-being \cite{ref195,ref1}. The Health-LLM framework \cite{ref195}, for example, combines large-scale feature extraction, medical knowledge integration, and a retrieval-augmented generation mechanism to provide personalized disease prediction and management recommendations, tailored to the individual's unique health profile and needs.

Furthermore, the multifaceted nature of LMLMs allows for their integration with various public health initiatives, such as telehealth services and remote patient monitoring programs. By seamlessly bridging the gap between patients and healthcare providers, these models can enhance accessibility, improve communication, and facilitate better-informed decision-making \cite{ref71}. The novel interface proposed in the ``Redefining Digital Health Interfaces with Large Language Models'' paper, for instance, showcases how LLM-based systems can leverage external tools to provide a unique and user-friendly interface between clinicians and digital health technologies, addressing current challenges with the direct application of LLMs in clinical settings.

Despite the immense potential of LMLMs in public health and wellness applications, there are several key challenges that must be addressed. Ensuring the accuracy, reliability, and safety of the information generated by these models is of paramount importance, particularly in sensitive healthcare domains \cite{ref212}. Rigorous evaluation, transparent development processes, and the integration of robust safety mechanisms are critical to mitigate the risks of spreading misinformation or providing inappropriate recommendations \cite{ref213}. Additionally, ethical considerations, such as data privacy, bias, and the equitable distribution of these technologies, must be carefully navigated to ensure that the benefits of LMLMs are accessible to all \cite{ref214}.

As the field of large multi-modal language models continues to evolve, their transformative impact on public health and wellness will only become more pronounced. By seamlessly integrating diverse data sources, enabling personalized interventions, and facilitating the widespread dissemination of critical health information, these models hold the potential to revolutionize the way we approach some of the most pressing public health challenges of our time. With thoughtful development, rigorous evaluation, and a commitment to ethical and equitable practices, the promise of LMLMs in public health and wellness can be fully realized, ultimately leading to improved health outcomes and enhanced quality of life for individuals and communities worldwide.

\subsection{Sustainability and Climate Applications}

The applications of large multi-modal language models (LMLMs) in the domains of sustainability and climate are wide-ranging and promising. These powerful AI systems have the potential to revolutionize how we approach critical challenges related to climate change, environmental protection, and sustainable urban development.

One of the key areas where LMLMs can have a significant impact is in the summarization and communication of climate-related information. The vast amount of scientific literature, policy documents, and real-time data on climate change can be overwhelming, making it challenging for policymakers, stakeholders, and the general public to stay informed and make informed decisions. LMLMs can be leveraged to process and distill this wealth of information, providing concise and accessible summaries that highlight the most critical insights and implications \cite{ref215}.

Furthermore, LMLMs can enhance our understanding of the complex interactions between climate, environment, and human activities. By integrating data from various modalities, such as satellite imagery, sensor networks, and textual reports, these models can offer valuable insights into patterns, trends, and causal relationships \cite{ref216}. This can inform urban planning and decision-making processes, enabling more sustainable and resilient cities \cite{ref217}.

In the context of disaster management, LMLMs can play a crucial role in rapid response and recovery efforts. By leveraging multimodal data, such as satellite imagery, social media, and emergency reports, these models can provide timely and accurate assessments of the extent and impact of natural disasters, like floods and wildfires. This information can then be used to coordinate relief efforts, allocate resources, and guide decision-making, ultimately minimizing the devastating consequences of such events \cite{ref218}.

Moreover, LMLMs can contribute to the development of innovative sustainable technologies and solutions. For example, these models can assist in the design and optimization of renewable energy systems, smart grid management, and energy-efficient buildings \cite{ref101}. By integrating diverse data sources and leveraging their natural language understanding capabilities, LMLMs can identify patterns, generate insights, and propose innovative strategies to address the pressing challenges of sustainability and climate change.

However, the integration of LMLMs in sustainability and climate applications is not without its challenges. Ensuring the accuracy, reliability, and ethical use of these models is paramount, as their outputs can have significant real-world implications \cite{ref219}. Rigorous evaluation and validation frameworks, as well as transparent and responsible development practices, are essential to mitigate the risks of biased or misleading information.

Furthermore, the successful deployment of LMLMs in these domains requires close collaboration between domain experts, AI researchers, and policymakers. By bridging the gap between these stakeholders, we can harness the full potential of LMLMs to drive sustainable and climate-resilient solutions, ultimately contributing to a more sustainable and equitable future \cite{ref220}.

In conclusion, the applications of large multi-modal language models in the domains of sustainability and climate are multifaceted and hold immense promise. From climate information summarization and urban planning to disaster management and sustainable technology development, these powerful AI systems can significantly enhance our understanding, decision-making, and problem-solving capabilities in the face of pressing environmental challenges. As we continue to advance the development and responsible deployment of LMLMs, we can unlock new possibilities for a more sustainable and resilient planet.

\subsection{Accessibility and Inclusivity Applications}

Large multimodal language models (MLLMs) possess immense potential to enhance accessibility and inclusivity, particularly for individuals with diverse needs and backgrounds. These models, which integrate language, vision, audio, and other modalities, can create multimodal interfaces that cater to a wide range of users, breaking down barriers and promoting more equitable access to information and services.

One key aspect of utilizing MLLMs for accessibility is their ability to provide alternative communication channels for individuals with disabilities. For example, MLLMs can generate personalized audio descriptions of visual content for users with visual impairments \cite{ref17}, or offer sign language interpretation for deaf or hard-of-hearing individuals \cite{ref221}. Furthermore, these models can assist users with motor impairments by enabling hands-free interaction through voice commands or gaze-based controls \cite{ref222}. By seamlessly integrating multiple modalities, MLLMs can empower users to choose the most suitable communication method, fostering greater independence and participation in various domains.

Beyond accessibility, MLLMs can also contribute to inclusivity by tailoring their outputs to diverse cultural and linguistic backgrounds. Recognizing the limitations of English-centric models, researchers have developed multilingual and culturally-aware MLLMs, such as PALO, which offers visual reasoning capabilities in 10 major languages spanning approximately 5 billion people \cite{ref223}. These models not only enhance accessibility for non-English speakers but also promote cultural preservation and representation, allowing users to engage with technology in their native languages and contexts.

Moreover, MLLMs can play a pivotal role in addressing the needs of marginalized communities and underrepresented groups. By incorporating diverse datasets and training strategies that capture the nuances of different demographics, MLLMs can be designed to recognize and respond to the unique challenges and perspectives of these groups \cite{ref224}. This can lead to more inclusive and empathetic interactions, where users from diverse backgrounds feel recognized, respected, and supported by the technology.

One particularly compelling application of MLLMs for inclusivity is in the field of education. These models can create personalized learning experiences that cater to individual needs, learning styles, and cultural backgrounds \cite{ref225}. By integrating multimodal content, such as interactive visualizations, audio-visual explanations, and culturally-relevant examples, MLLMs can enhance accessibility and engagement for students with diverse abilities and backgrounds. This can contribute to more equitable and inclusive educational environments, fostering the success of learners who may have previously faced barriers.

Furthermore, the adaptability and flexibility of MLLMs offer opportunities to address the needs of specific communities. For instance, these models can be tailored to support the unique communication requirements of neurodiverse individuals, such as those with autism spectrum disorder or cognitive impairments \cite{ref226}. By understanding and accommodating the preferences and challenges of these users, MLLMs can create more inclusive experiences that empower individuals to participate and thrive in various aspects of life.

In conclusion, the integration of large multimodal language models holds immense promise for enhancing accessibility and inclusivity across various domains. By leveraging their ability to seamlessly process and generate multimodal content, MLLMs can break down barriers, promote representation, and foster more equitable access to information and services. As the field of MLLMs continues to evolve, it is crucial to prioritize the development of models that prioritize the needs of diverse users, ensuring that the transformative potential of these technologies is realized in a way that benefits all members of society.

\section{Ethical Considerations and Societal Impact}

\subsection{Ethical Implications of Large Multi-Modal Language Models}

The rapid emergence of large multi-modal language models (LMMLs) has brought about significant advancements in natural language processing and generation, enabling a wide range of applications across various domains. However, these powerful models also raise pressing ethical concerns that need to be carefully examined and addressed.

One of the primary ethical implications of LMMLs is the potential for privacy violations. These models are trained on large and diverse datasets, which may inadvertently contain sensitive personal information \cite{ref227}. As LMMLs become increasingly adept at understanding and generating human-like text, there is a risk that they could be used to infer or even extract private details about individuals, compromising their right to privacy \cite{ref79}. This concern is particularly acute in scenarios where LMMLs are integrated into user-facing applications, such as virtual assistants or chatbots, where users may be prompted to disclose sensitive information \cite{ref228}.

Another critical ethical implication of LMMLs is the potential for perpetuating or amplifying biases and unfair discrimination \cite{ref229}. These models are trained on data that often reflects societal biases and inequalities, and their outputs may inadvertently reinforce or exacerbate these biases \cite{ref82}. This can have far-reaching consequences, particularly in domains such as healthcare, education, or employment, where the decisions made by LMMLs could have a significant impact on people's lives \cite{ref230}.

The potential for misuse of LMMLs also raises serious ethical concerns. These models can be leveraged for malicious purposes, such as generating misinformation, impersonating individuals, or creating harmful content \cite{ref85}. This could lead to the erosion of public trust, the spread of false narratives, and the exploitation of vulnerable individuals or communities \cite{ref86,ref81,ref79}.

To ensure the responsible development and deployment of LMMLs, it is crucial to implement robust ethical frameworks and oversight mechanisms. This may involve establishing clear guidelines and standards for the development, testing, and deployment of these models, with a focus on safeguarding privacy, promoting fairness, and mitigating the potential for misuse \cite{ref84}.

Additionally, the development of LMMLs should incorporate principles of transparency and accountability, allowing for the scrutiny and oversight of these models by various stakeholders, including the public, policymakers, and independent ethical review boards \cite{ref231,ref232,ref233}.

Furthermore, the integration of ethical considerations should be a fundamental aspect of the entire lifecycle of LMMLs, from the data curation and model training stages to the deployment and monitoring phases \cite{ref234,ref235}.

By proactively addressing these ethical concerns and implementing robust safeguards, the research and development community can work towards ensuring the responsible and beneficial use of LMMLs, maximizing their potential to enhance human knowledge and capabilities while minimizing the risks of unintended harm or societal disruption.

\subsection{Privacy Concerns and Mitigation Strategies}

The rapid advancement of large multi-modal language models has raised significant privacy concerns due to their ability to process and generate vast amounts of data, including potentially sensitive or personally identifiable information (PII). A key concern is the potential for these models to leak or inadvertently expose private data from their training corpus \cite{ref236,ref237}.

One of the primary privacy concerns is the potential for personal data leaks, where sensitive information such as email addresses, phone numbers, or other PII could be extracted from the model's outputs \cite{ref238,ref236}. This risk is exacerbated by the fact that large language models are often trained on a diverse range of web-crawled data, which may contain a significant amount of personal information \cite{ref239,ref240}.

Furthermore, large multi-modal language models have the capability to infer and associate different pieces of information, which raises concerns about the potential exposure of sensitive data, even if it is not directly present in the training corpus \cite{ref237}. As these models continue to grow in size and complexity, their ability to make such inferences is likely to increase, posing an even greater threat to user privacy.

To address these privacy concerns, researchers have explored various strategies and frameworks for protecting user privacy in the context of large multi-modal language models. One prominent approach is the use of differential privacy, a mathematical framework that provides provable guarantees of privacy by introducing carefully calibrated noise into the training process or the model's outputs \cite{ref241,ref242}. By applying differential privacy techniques, researchers have shown that it is possible to train language models that provide strong privacy guarantees while still maintaining a high level of utility \cite{ref243,ref244}.

Another strategy for protecting user privacy in the context of large multi-modal language models is the use of text sanitization and prompt-based approaches \cite{ref245,ref246}. These methods focus on identifying and removing or obfuscating sensitive information within the model's inputs or outputs, effectively preventing the exposure of private data.

Additionally, researchers have explored the use of multi-party computation and secure enclaves to enable the use of large multi-modal language models while maintaining the privacy of the user's data \cite{ref247,ref105}. These approaches aim to keep the user's data secure and prevent it from being accessed or exploited by the model provider or other third parties.

It is important to note that while these privacy-preserving strategies and frameworks show promise, they are not without their own challenges and limitations. Balancing the trade-off between privacy protection and model utility remains a significant challenge, as some privacy-preserving techniques may come at the cost of reduced model performance or functionality \cite{ref248,ref249}. Furthermore, the complexity of large multi-modal language models and the diversity of data sources they are trained on make it difficult to provide a one-size-fits-all solution for privacy protection.

As the adoption and applications of large multi-modal language models continue to grow, it is essential that the research community and industry work together to address the privacy concerns and develop robust solutions that prioritize user privacy while still enabling the transformative potential of these technologies. By proactively addressing these challenges, we can ensure the responsible and ethical development and deployment of large multi-modal language models for the benefit of society.

\subsection{Fairness and Bias Considerations}

Ensuring fairness and addressing biases in large multi-modal language models (LMLMs) is a critical challenge that requires a multifaceted approach. As these powerful models become increasingly integrated into various applications, it is essential to mitigate the propagation of harmful biases and promote equitable outcomes for all users.

A significant body of research has examined the challenges of fairness and bias in language models, particularly in the context of unimodal (text-only) models \cite{ref100,ref250}. However, the integration of multiple modalities, such as vision and audio, in LMLMs introduces additional complexities and unique considerations.

One crucial challenge lies in the inherent biases present in the training data used to develop these multi-modal models \cite{ref251}. The datasets used to train LMLMs often reflect societal biases and stereotypes, leading to the perpetuation of these biases in the model's outputs. This issue is particularly pronounced in domains where historical data is scarce or heavily skewed, such as in the representation of marginalized communities \cite{ref252}.

Furthermore, the integration of multiple modalities introduces the potential for cross-modal biases, where the interaction between different input sources amplifies or exacerbates existing biases \cite{ref253}. For example, in a visual question-answering task, the model's predictions may be influenced by a combination of language and visual biases, leading to compounded discriminatory outcomes.

Addressing these challenges requires a multi-pronged approach that combines the development of robust evaluation frameworks, the implementation of bias mitigation strategies, and the active involvement of diverse stakeholders \cite{ref254}.

Firstly, comprehensive evaluation frameworks are needed to assess the fairness and bias in LMLMs \cite{ref255}. These frameworks should encompass a diverse range of demographic attributes, intersectional identities, and task-specific considerations to capture the nuances of bias \cite{ref256}. Additionally, the evaluation process should incorporate both automated metrics and human-centric assessments to ensure a holistic understanding of model fairness \cite{ref257}.

Secondly, bias mitigation strategies should be developed and incorporated into the design and training of LMLMs. These strategies may include techniques such as data augmentation, prompt engineering, and fine-tuning on debiased datasets \cite{ref258,ref259}. Additionally, the exploration of causal reasoning and interpretability frameworks can provide valuable insights into the sources of bias and guide the development of more equitable models \cite{ref260}.

Lastly, the active involvement of diverse stakeholders, including domain experts, policymakers, and marginalized community representatives, is crucial in ensuring that the development and deployment of LMLMs align with societal values and promote equitable access and outcomes \cite{ref261}. This collaborative approach can help identify and address blind spots, incorporate context-specific considerations, and foster a shared understanding of the ethical implications of these powerful technologies.

By addressing the challenges of fairness and bias in LMLMs, researchers and developers can work towards creating more equitable and inclusive AI systems that serve the diverse needs of all users. This endeavor requires a sustained and multidisciplinary effort, drawing on the expertise of various fields, including machine learning, social sciences, and ethics, to ensure that the transformative potential of large multi-modal language models is harnessed in a responsible and fair manner.

\subsection{Potential for Misuse and Malicious Use Cases}

The potential for large multi-modal language models (MMLLMs) to be misused for harmful or malicious purposes is a growing concern that warrants careful examination. These powerful models have demonstrated remarkable capabilities in understanding and generating human-like text, images, and multimedia content, which can be exploited by bad actors for nefarious purposes.

One of the primary risks associated with MMLLMs is their potential to be used for the generation of disinformation. Malicious actors could leverage the model's ability to produce coherent and seemingly credible content to spread false narratives, manipulate public opinion, and undermine trust in legitimate sources of information \cite{ref262,ref263,ref264}. Additionally, MMLLMs can be exploited to generate content that directly incites harm, such as instructions for creating weapons or explosives, or detailed plans for criminal activities \cite{ref265}. The scale and speed at which these models can produce such content pose a significant challenge in terms of detection and mitigation \cite{ref266}.

Another concerning aspect of MLLM misuse is the potential for exploitation of vulnerabilities within the models themselves. Researchers have identified various security vulnerabilities in large language models, including the possibility of adversarial attacks that can manipulate the model's outputs \cite{ref267,ref268}. These vulnerabilities could be exploited by bad actors to bypass the models' safety mechanisms and generate harmful content or engage in other malicious activities \cite{ref269}.

Moreover, the integration of MMLLMs with various applications, such as virtual assistants or interactive chatbots, introduces new avenues for misuse. Malicious actors could leverage these applications to impersonate legitimate entities, obtain sensitive information, or engage in social engineering attacks \cite{ref270,ref79}.

To mitigate these risks, several measures can be implemented. Developers and researchers should continue to invest in the development of robust safety mechanisms, such as advanced content filtering, anomaly detection, and multi-layered security protocols \cite{ref271,ref86}. Additionally, comprehensive testing and evaluation frameworks should be established to identify and address vulnerabilities before deployment \cite{ref272,ref273}.

Regulatory frameworks and industry-wide standards may also play a crucial role in ensuring the responsible development and deployment of MMLLMs \cite{ref274}. Policymakers, ethicists, and security experts should collaborate to develop guidelines and guidelines that balance the potential benefits of these technologies with the necessary safeguards to mitigate the risks of misuse \cite{ref275}. Furthermore, public awareness and education campaigns are essential to empower users to identify and report suspicious or harmful MLLM-generated content \cite{ref276}. By fostering a collective effort of stakeholders, the impact of MLLM misuse can be minimized, and these powerful tools can be used for the greater good of society.

\subsection{Ethical Governance and Oversight}

The responsible development and deployment of large multi-modal language models (MM-LLMs) requires robust ethical governance and oversight mechanisms to ensure these powerful technologies are aligned with societal values and interests. Policymakers, industry leaders, and the research community all have pivotal roles to play in establishing effective governance frameworks that can navigate the complex ethical terrain of MM-LLMs.

At the policy level, there is a growing recognition of the need for comprehensive regulatory frameworks to govern the use of generative AI technologies, including MM-LLMs \cite{ref277,ref278}. Initiatives like the European Union's proposed AI Act demonstrate efforts to establish common standards and guidelines for the development, testing, and deployment of AI systems \cite{ref278}. However, the rapid pace of technological change and the diversity of applications for MM-LLMs pose significant challenges in crafting legislation that is both sufficiently flexible and enforceable.

Policymakers must work closely with industry stakeholders and the research community to develop governance models that can effectively bridge the gap between high-level ethical principles and practical implementation \cite{ref279}. One promising approach is the ``Dual Governance'' framework, which proposes a synergistic relationship between centralized regulations and decentralized, community-driven safety mechanisms \cite{ref280}. This model recognizes the limitations of top-down regulations alone and seeks to leverage the agility and innovation of crowdsourced efforts to complement and enhance formal oversight.

Industry players, such as technology companies developing and deploying MM-LLMs, also have a critical role to play in establishing robust ethical governance. The ``Hourglass Model of Organizational AI Governance'' provides a comprehensive framework for integrating ethical considerations throughout the entire lifecycle of AI systems, from environmental scanning to model deployment and monitoring \cite{ref281}. This multi-layered approach ensures that ethical principles are not merely aspirational but are actively incorporated into the design, development, and deployment of MM-LLMs.

The research community, meanwhile, can contribute to ethical governance through the development of rigorous evaluation frameworks and auditing mechanisms \cite{ref282}. By establishing standardized benchmarks and metrics for assessing the ethical alignment of MM-LLMs, researchers can provide policymakers and industry stakeholders with the tools necessary to make informed decisions about the responsible use of these technologies \cite{ref233}.

Additionally, the research community can play a vital role in fostering inclusive and participatory governance processes, ensuring that diverse perspectives, including those from marginalized communities, are represented in the development of ethical frameworks \cite{ref279}. This collaborative approach can help to mitigate the risk of ethical blind spots and ensure that the deployment of MM-LLMs promotes societal well-being rather than entrenching existing inequities \cite{ref283}.

Ultimately, the responsible development and deployment of large multi-modal language models will require a concerted effort from policymakers, industry leaders, and the research community, working in tandem to establish robust and adaptive governance frameworks. By leveraging the complementary strengths of centralized regulations, decentralized safety mechanisms, and participatory processes, these stakeholders can navigate the complex ethical landscape of MM-LLMs and ensure these transformative technologies are aligned with the values and interests of society as a whole.

\subsection{Societal Impact and Implications}

The rapid advancements in large multi-modal language models (LMLMs) have far-reaching implications for various aspects of society. These transformative technologies have the potential to revolutionize how we interact with information, learn, work, and engage with one another, leading to both promising opportunities and challenging ethical considerations.

One significant area of societal impact is the potential influence of LMLMs on employment and the job market. As discussed in \cite{ref284}, LMLMs exhibit characteristics of general-purpose technologies, suggesting that they could have considerable economic and social implications. The study estimates that around 80\% of the U.S. workforce could have at least 10\% of their work tasks affected by the introduction of LMLMs, with approximately 19\% of workers potentially seeing at least 50\% of their tasks impacted. This finding underscores the disruptive potential of these technologies, as they can automate or augment a wide range of tasks, from content creation to problem-solving.

The implications of this shift extend beyond job displacement, as LMLMs may also transform the nature of work itself. \cite{ref285} highlights the potential of LMLMs to serve as research tools, assisting scientists and professionals in tasks such as literature reviews, data analysis, and writing. While this can boost productivity and efficiency, it also raises concerns about the changes in skill requirements and the potential deskilling of certain professions. As LMLMs become more integrated into the workplace, there will be a need to rethink job descriptions, education, and training to ensure that workers are equipped with the necessary skills to collaborate effectively with these technologies.

In the realm of education, the integration of LMLMs has both promising and challenging implications. \cite{ref286} and \cite{ref90} explore the potential of LMLMs to enhance educational experiences, such as by automating certain assessment tasks, providing personalized feedback, and streamlining administrative processes. However, these technologies also raise concerns about academic integrity, as students may be tempted to misuse LMLMs for tasks like essay writing or code generation, potentially undermining learning objectives. \cite{ref287} further emphasizes the need for robust strategies to ensure that the integration of LMLMs in education fosters critical thinking, creativity, and authentic learning, rather than simply outsourcing tasks to these models.

Beyond employment and education, LMLMs have the potential to impact social dynamics and community interactions. \cite{ref89} explores the implications of LMLMs for democratic discourse, highlighting the risks of these models being perceived as authoritative sources of information and the potential for them to erode the authenticity of human reasoning and decision-making. As LMLMs become more widely adopted, there will be a growing need to educate the public on the capabilities and limitations of these technologies, empowering individuals to critically evaluate the information they receive and engage in informed, nuanced discussions.

Furthermore, the widespread use of LMLMs raises concerns about the equitable access and distribution of these technologies. \cite{ref288} reveals a gender gap in the adoption of LMLMs, with male users outpacing female users. This disparity underscores the importance of ensuring that the benefits and risks of LMLMs are evenly distributed across diverse demographics, as the uneven adoption of these technologies could exacerbate existing social and economic inequalities.

To address these societal implications, a multi-pronged approach is necessary. Policymakers, educators, and industry leaders must collaborate to develop robust governance frameworks, ethical guidelines, and educational initiatives that promote the responsible development and deployment of LMLMs. \cite{ref105} and \cite{ref173} emphasize the importance of addressing the security and privacy concerns associated with LMLMs, as well as the need to align these technologies with societal needs and values.

Additionally, \cite{ref289} highlights the importance of understanding the cultural biases and limitations of LMLMs, as these models may perpetuate or amplify existing societal biases. Ongoing research and evaluation of the societal impact of LMLMs, as well as transparent communication about their capabilities and limitations, will be crucial in ensuring that these technologies benefit society as a whole.

\section{Enhancing Multi-Modal Reasoning through Prompting and In-Context Learning}

\subsection{Prompt Design Strategies for Enhancing Multi-Modal Reasoning}

The rapid advancements in large language models (LLMs) have enabled the development of multi-modal capabilities, where these models can seamlessly integrate and process information from various modalities, including text, images, and even audio \cite{ref94}. This integration of multi-modal information has significantly enhanced the reasoning abilities of these models, allowing them to tackle complex tasks that require a comprehensive understanding of the problem context.

One of the key strategies for improving the multi-modal reasoning capabilities of LLMs is through the design of effective prompts. Prompt design has emerged as a powerful tool for harnessing the full potential of these models, and researchers have explored various approaches to enhance the multi-modal reasoning abilities of LLMs through prompting \cite{ref94,ref145}.

One such approach is to prompt the LLM with intermediate reasoning steps, where the model is provided with a series of steps or a ``chain of thought'' that guides the reasoning process \cite{ref290,ref94}. By breaking down a complex problem into smaller, intermediate steps, the LLM can better understand the underlying logic and reasoning required to arrive at the final solution. This approach has been particularly effective in tasks that involve multi-step reasoning, such as mathematical problem-solving or scientific reasoning \cite{ref145,ref96}.

Another strategy for enhancing multi-modal reasoning is the use of retrieval-based prompts, where the LLM is prompted to retrieve relevant information from external sources to aid in the reasoning process \cite{ref99}. This approach is particularly useful in scenarios where the necessary information is not readily available within the model's own parameters, such as when dealing with complex, domain-specific knowledge \cite{ref291}. By integrating external knowledge sources, the LLM can expand its understanding of the problem context and make more informed decisions during the reasoning process.

Furthermore, researchers have explored the incorporation of knowledge from external sources, such as knowledge graphs or other structured data, to further enhance the multi-modal reasoning capabilities of LLMs \cite{ref96}. By grounding the reasoning process in comprehensive and structured knowledge, LLMs can make more informed and accurate decisions, reducing the likelihood of hallucinations or other types of errors \cite{ref292}.

The impact of these prompt design strategies on the multi-modal reasoning capabilities of LLMs has been extensively studied. For instance, \cite{ref94} demonstrated that a two-stage framework that separates the generation of rationales and the inference of answers can significantly improve the performance of LLMs on scientific reasoning tasks, even surpassing human-level performance. Similarly, \cite{ref145} introduced a novel prompting approach that divides the reasoning responsibility between the LLM and the visual recognition model, leading to substantial improvements in multimodal reasoning tasks.

Moreover, \cite{ref96} proposed a framework that integrates chain-of-thought reasoning, knowledge graphs, and multiple modalities to provide a comprehensive understanding of multimodal tasks. Their approach not only outperformed state-of-the-art methods but also achieved these results with a significantly smaller model, demonstrating the cost-effectiveness and efficiency of their prompt design strategy.

Overall, the design of effective prompts has emerged as a critical factor in enhancing the multi-modal reasoning capabilities of LLMs. By leveraging intermediate reasoning steps, retrieval-based prompts, and the incorporation of external knowledge, researchers have demonstrated the potential of these models to tackle complex, multi-modal problems with increasing accuracy and robustness. As the field of multi-modal language modeling continues to evolve, the exploration of novel prompt design strategies will likely remain a key area of focus, paving the way for even more advanced and versatile LLMs.

\subsection{Leveraging Multimodal In-Context Learning}

The integration of visual and language information in in-context examples has the potential to significantly enhance the reasoning abilities of large language models (LLMs). Recent advancements in multimodal in-context learning (ICL) have revealed the complementary benefits of combining textual and visual cues to aid LLMs in tackling complex tasks \cite{ref293,ref294,ref295}.

One key aspect of leveraging multimodal ICL is the effective design and selection of in-context examples. Studies have shown that the performance of ICL is highly sensitive to the quality and composition of the in-context examples \cite{ref296,ref297}. Carefully curated in-context examples that strike a balance between visual and textual information can lead to significant performance gains compared to using either modality alone \cite{ref95}.

Researchers have explored various techniques to enhance the impact of multimodal in-context examples. One approach is to leverage visual-textual alignment and retrieval methods to select the most relevant and complementary in-context examples \cite{ref293,ref296}. By considering both the visual and textual aspects of the in-context examples, these techniques can effectively capture the crucial information required for the task at hand, leading to improved ICL performance.

Another promising direction is the use of multimodal prompting, where the in-context examples combine visual and textual information in a unified format \cite{ref294,ref298}. This approach aims to bridge the gap between the multimodal nature of the task and the mode of in-context information presentation, thereby facilitating more seamless integration and reasoning by the LLMs.

Researchers have also explored techniques to enhance the robustness and generalizability of multimodal ICL. One such approach is Bidirectional Alignment (BiAlign), which aligns the input preferences of smaller and larger models to effectively transfer the powerful ICL capabilities of larger models to more compact models \cite{ref295}. By considering both the visual and textual aspects of the in-context examples, BiAlign can improve the ICL performance of smaller models, making them more efficient and deployable in real-world scenarios.

The incorporation of visual information into in-context examples has also been shown to benefit specific tasks, such as in-context visual reasoning \cite{ref299}. By visualizing the reasoning steps or intermediate outputs, LLMs can better understand the task requirements and leverage their visual processing capabilities to enhance their overall reasoning abilities.

Furthermore, the use of multimodal in-context examples has been explored in the context of video understanding tasks \cite{ref300}. By presenting LLMs with pairs of instructions and corresponding high-level visual programs, researchers have demonstrated the ability to harness the reasoning power of LLMs to tackle complex video-based tasks.

Despite the promising results, the integration of visual and textual information in in-context examples is not without its challenges. One key challenge is the potential mismatch between the modalities, where the visual information may not directly align with the textual information, leading to confusion or suboptimal performance \cite{ref297}. Addressing this challenge requires careful design of multimodal in-context examples and potentially the development of more robust multimodal integration techniques.

Another challenge is the computational and memory requirements associated with the inclusion of visual information in the in-context examples. As the size and complexity of the in-context examples increase, the computational and memory demands on the LLMs also grow, which can limit their practical applicability \cite{ref301}. Developing efficient techniques for representing and processing multimodal in-context examples is an active area of research.

Overall, the integration of visual and textual information in in-context examples has shown great promise in enhancing the reasoning abilities of LLMs. By leveraging the complementary strengths of the two modalities, researchers have demonstrated improvements in a wide range of tasks, from visual reasoning to video understanding. However, further advancements are needed to address the challenges associated with multimodal ICL, such as modality mismatch, computational efficiency, and robustness, to unlock the full potential of this approach.

\subsection{Mitigating Biases and Limitations in Multimodal In-Context Learning}

One of the key limitations of current multimodal in-context learning (M-ICL) approaches is their susceptibility to biases and limitations that can undermine the robustness and generalizability of the reasoning process. Researchers have identified several potential biases and limitations in M-ICL that require careful consideration and mitigation strategies.

One prominent issue is the overreliance on textual information in M-ICL. Several studies have found that M-ICL models tend to heavily rely on the textual modality, often disregarding or underutilizing the visual information available \cite{ref297}. This bias towards textual information can lead to suboptimal performance on tasks that require a more balanced integration of visual and textual cues, such as visual question answering \cite{ref251}.

To address this limitation, researchers have explored methods to encourage M-ICL models to better leverage the visual modality. One approach is the use of multimodal prompting strategies, where the in-context examples incorporate both textual and visual information \cite{ref302}. By exposing the M-ICL model to a more diverse set of multimodal examples during the in-context learning phase, the model can learn to more effectively integrate and reason with both modalities, reducing the overreliance on textual information.

Another potential limitation of M-ICL is the sensitivity to the quality and distribution of the in-context examples. Studies have shown that the performance of M-ICL can be highly dependent on the specific examples provided in the prompt, with certain examples leading to significantly better or worse performance \cite{ref303,ref304}. This sensitivity to the in-context examples can undermine the reliability and generalizability of M-ICL, as the model's behavior may be heavily influenced by the idiosyncrasies of the particular examples used.

To address this limitation, researchers have explored methods for actively selecting the most informative and representative in-context examples to include in the prompt. One approach is to use influence-based methods to identify examples that are likely to have a positive or negative impact on the M-ICL model's performance \cite{ref303}. By selectively including the most influential examples, the model can be guided towards more robust and generalizable reasoning. Additionally, the use of reinforcement learning algorithms to learn policies for selecting optimal in-context examples has shown promise in improving the stability and performance of M-ICL \cite{ref305}.

Another important limitation of M-ICL is the potential for the model to exhibit biases, such as those related to gender, race, or other societal constructs. These biases can be particularly problematic in sensitive applications, such as healthcare or education, where the model's decisions can have significant real-world consequences \cite{ref306}. Researchers have explored various debiasing techniques, including prompt engineering, fine-tuning, and the use of causal reasoning frameworks, to mitigate the impact of these biases in M-ICL \cite{ref307,ref251}.

In addition to these biases, M-ICL models can also suffer from limitations in their ability to handle out-of-distribution or novel scenarios. Studies have found that M-ICL performance can degrade significantly when faced with distributional shifts or unseen combinations of modalities \cite{ref308}. Addressing this limitation requires the development of more robust and adaptable M-ICL models that can generalize beyond the specific training distributions.

Overall, the mitigation of biases and limitations in multimodal in-context learning is a crucial area of research that requires a multifaceted approach. Strategies such as multimodal prompting, active example selection, debiasing techniques, and improved out-of-distribution generalization can all contribute to enhancing the robustness and reliability of M-ICL models. By addressing these challenges, researchers can pave the way for the widespread and responsible deployment of M-ICL in a wide range of real-world applications.

\subsection{Multimodal Prompting for Enhanced Reasoning}

The rapid advancements in large language models (LLMs) have ushered in a new era of in-context learning, where models can adapt to new tasks and scenarios by leveraging a few demonstration examples provided as a prompt \cite{ref301,ref309}. However, as these models are predominantly trained on textual data, their reasoning capabilities are primarily biased towards the language modality. To further enhance the reasoning abilities of LLMs, researchers have explored the incorporation of visual information through the development of multimodal language models (MLLMs) \cite{ref310}.

One prominent issue is the overreliance on textual information in M-ICL. Several studies have found that M-ICL models tend to heavily rely on the textual modality, often disregarding or underutilizing the visual information available \cite{ref297}. This bias towards textual information can lead to suboptimal performance on tasks that require a more balanced integration of visual and textual cues, such as visual question answering \cite{ref251}.

To address this limitation, researchers have explored various multimodal prompting techniques, which combine both visual and textual information to guide the in-context learning process. This approach aims to capitalize on the complementary strengths of the two modalities, where the visual information can provide additional context and grounding to support the reasoning performed by the language model.

In \cite{ref299}, the authors propose a ``Visualization-of-Thought'' (VoT) prompting strategy, where they generate visual representations of the reasoning steps to enhance the spatial reasoning capabilities of LLMs. By visualizing the intermediate reasoning steps, the model is better able to understand and follow the logical flow of the problem-solving process, leading to improved performance on tasks that require spatial reasoning.

Similarly, \cite{ref311} introduces a novel approach called ``Multi-modal Prompt Learning'' (MaPle), which learns prompts for both the visual and language branches of a MLLM. This approach promotes strong coupling between the vision and language prompts, ensuring that the model can effectively leverage the information from both modalities to improve its performance on downstream tasks.

Additionally, \cite{ref312} proposes a ``Joint Visual and Text Prompting'' (VTPrompt) method, which combines fine-grained visual information with textual prompts to enhance the object-centric perception capabilities of MLLMs. By providing the model with both visual cues and textual information, VTPrompt helps the model better understand and integrate the relevant visual and linguistic elements, leading to improved performance on tasks that require a deep understanding of the visual and textual content.

Beyond these specific prompting techniques, researchers have also explored more general frameworks for leveraging multimodal information to improve the reasoning abilities of LLMs. For instance, \cite{ref313} introduces a ``Causal Context Generation'' (Causal-CoG) approach, which prompts the model to generate contextual information (i.e., text descriptions of an image) and then uses this context to enhance the model's performance on visual question answering tasks.

Similarly, \cite{ref314} proposes a ``browse-and-concentrate'' paradigm, where the model first ``browses'' through the multimodal inputs to gain essential insights, and then ``concentrates'' on the crucial details to achieve a more comprehensive understanding of the multimodal content.

These multimodal prompting techniques, along with the broader frameworks for leveraging multimodal information, demonstrate the potential of integrating visual and textual cues to enhance the reasoning capabilities of large language models. By harnessing the complementary strengths of the two modalities, these approaches can help LLMs overcome the limitations of text-only reasoning and unlock new levels of multimodal understanding and problem-solving.

However, the field of multimodal prompting for enhanced reasoning is still relatively nascent, and there are several open challenges and future research directions to explore. For example, researchers could investigate how to design more effective and generalizable multimodal prompts, explore the interplay between different modalities in the context of in-context learning, and study the underlying mechanisms by which multimodal prompting can improve the model's reasoning abilities.

Additionally, the integration of multimodal prompting with other advanced techniques, such as chain-of-thought reasoning \cite{ref94} and visual-language modeling \cite{ref315}, could further enhance the reasoning capabilities of LLMs and unlock new possibilities for intelligent systems that can seamlessly navigate and reason about the multimodal world.

\subsection{Evaluation Frameworks for Multimodal Reasoning}

Assessing the multimodal reasoning capabilities of large language models (LLMs) has become a crucial endeavor in the pursuit of advancing artificial intelligence. To this end, the development of comprehensive evaluation frameworks and benchmarks plays a pivotal role in understanding the strengths, limitations, and potential of these models.

One of the key challenges in evaluating multimodal reasoning lies in the design of diverse and challenging tasks that require the seamless integration of visual and textual information. Traditional benchmarks, such as visual question answering (VQA) and image captioning, have provided valuable insights into the language-vision capabilities of LLMs. However, these tasks often focus on relatively straightforward associations between image content and textual descriptions, failing to capture the nuances of complex multimodal reasoning \cite{ref185,ref152}.

To address this limitation, researchers have introduced more sophisticated evaluation frameworks that challenge LLMs to engage in reasoning that goes beyond simple associations. For instance, the MM-BigBench benchmark \cite{ref51} incorporates a diverse range of metrics, including Best Performance, Mean Relative Gain, and Stability, to assess the performance of LLMs on a wide spectrum of multimodal content comprehension tasks. This framework not only evaluates the overall performance of LLMs but also examines their adaptability to different instructions, providing a more comprehensive understanding of their multimodal reasoning abilities.

Another innovative approach is the ChEF framework \cite{ref54}, which introduces a modular and scalable design that allows for the creation of customized evaluation recipes. These recipes encompass various components, such as Scenario (multimodal datasets), Instruction (flexible instruction retrieving formulae), Inferencer (reliable question answering strategies), and Metric (indicative task-specific score functions). By leveraging this framework, researchers can design targeted evaluations that assess specific capabilities of LLMs, such as calibration, in-context learning, and robustness, providing a more nuanced understanding of their multimodal reasoning skills.

Another notable benchmark, NPHardEval4V \cite{ref316}, focuses specifically on evaluating the pure reasoning abilities of LLMs by converting textual description of questions into image representations. This approach aims to disentangle the effect of various factors, such as image recognition and instruction following, from the overall performance of the models, allowing for a more focused assessment of their reasoning capabilities.

In addition to these benchmarks, researchers have also explored the use of synthetic data generation to create evaluation datasets that are tailored to specific multimodal reasoning tasks \cite{ref317}. By leveraging techniques such as high-resolution text-to-image generation, these approaches can generate diverse and challenging evaluation scenarios that address the limitations of naturally occurring datasets.

Furthermore, the importance of benchmarking multimodal reasoning capabilities across different languages and domains has been recognized, leading to the development of frameworks like PALO \cite{ref223}, which aims to evaluate the vision-language reasoning capabilities of LLMs in 10 major languages. This effort highlights the need for inclusive and comprehensive evaluation approaches that cater to the diverse linguistic and cultural landscape of multimodal interactions.

Beyond task-specific benchmarks, researchers have also explored the concept of unified evaluation frameworks that can assess the overall capabilities of LLMs in a more holistic manner. For instance, the ReForm-Eval benchmark \cite{ref318} proposes a unified approach to leverage existing task-oriented benchmarks by reformulating them into a format compatible with LLMs, enabling a more comprehensive and objective evaluation of their performance.

The development of these comprehensive evaluation frameworks and benchmarks has not only shed light on the current limitations of LLMs in multimodal reasoning but has also inspired the exploration of novel techniques to enhance their capabilities. For example, the Cantor framework \cite{ref319} leverages the cognitive capabilities of multimodal LLMs to perform as multifaceted experts, deriving higher-level information and enhancing the chain-of-thought generation process for improved multimodal reasoning.

As the field of multimodal AI continues to evolve, the importance of comprehensive and standardized evaluation frameworks cannot be overstated. These frameworks not only help researchers and practitioners assess the progress of LLMs but also provide valuable insights into the underlying challenges and guide the development of more advanced models capable of seamless multimodal reasoning. By continuously refining and expanding these evaluation approaches, the research community can drive the advancement of multimodal intelligence and bring us closer to the ultimate goal of artificial general intelligence.

\section{Challenges and Limitations}

\subsection{Data Bias}

Data bias is a pervasive challenge that plagues large multi-modal language models, leading to biased and harmful outputs that can have far-reaching societal implications. These models are trained on vast amounts of data from the internet and other sources, which often reflect the biases and prejudices inherent in human-generated content \cite{ref100}.

One of the main sources of data bias in large multi-modal language models is the inherent biases present in the training data itself. Much of the data used to train these models is sourced from the internet, which can be rife with biases, stereotypes, and harmful narratives \cite{ref306}. For example, the language used in online forums, social media, and news articles may perpetuate negative stereotypes about certain demographic groups, such as gender, race, or ethnicity. When these models are trained on this data, they can learn and amplify these biases, leading to the generation of content that discriminates against or marginalizes these groups \cite{ref320}.

Another source of data bias in large multi-modal language models is the lack of diversity and representation in the data used for training. Many of these models are trained on data that is primarily generated by and reflects the perspectives of dominant or privileged groups, leaving out the experiences and voices of marginalized communities \cite{ref321}. This can lead to models that are biased towards the norms and values of the dominant group, and fail to capture the nuances and complexities of diverse lived experiences.

Furthermore, the process of data curation and annotation can also introduce biases into the training data. Decisions made by the individuals or teams responsible for selecting, cleaning, and labeling the data can inadvertently reflect their own biases and assumptions \cite{ref322}. For example, the choice of which data to include or exclude, or the way in which certain attributes or categories are defined, can all contribute to the biases present in the final dataset.

The consequences of data bias in large multi-modal language models can be severe and far-reaching. Biased outputs from these models can reinforce harmful stereotypes, perpetuate discrimination, and lead to the exclusion of certain individuals or groups from accessing or benefiting from the services and applications powered by these models \cite{ref323}. In high-stakes domains such as healthcare, finance, or criminal justice, these biases can have direct and tangible impacts on people's lives, leading to unfair treatment, denied opportunities, and even life-altering decisions.

Addressing the challenge of data bias in large multi-modal language models requires a multi-faceted approach that involves both technical and societal interventions. On the technical front, researchers and developers are exploring a range of strategies to mitigate these biases, such as \cite{ref324}:

\begin{enumerate}
\def\labelenumi{\arabic{enumi}.}
\tightlist
\item
  Improving data collection and curation practices to ensure greater diversity and representation in the training data.
\item
  Developing algorithmic techniques for debiasing the training data, such as data augmentation, adversarial training, or constrained optimization.
\item
  Incorporating bias detection and mitigation mechanisms directly into the model architecture, such as using specialized modules or training objectives to identify and reduce biases.
\item
  Implementing comprehensive evaluation and testing frameworks to assess the presence and magnitude of biases in the model's outputs, and using this information to guide further improvements.
\end{enumerate}

On the societal front, there is a growing recognition of the need for greater awareness, accountability, and collaboration to address the challenges of bias in large multi-modal language models. This includes \cite{ref255}:

\begin{enumerate}
\def\labelenumi{\arabic{enumi}.}
\tightlist
\item
  Engaging with diverse stakeholders, including marginalized communities, to better understand the nature and impact of biases in these models.
\item
  Developing ethical guidelines and governance frameworks to ensure the responsible development and deployment of large multi-modal language models, with a focus on fairness, transparency, and accountability.
\item
  Fostering interdisciplinary collaboration between technologists, policymakers, social scientists, and other relevant experts to create more holistic and comprehensive solutions to the problem of data bias.
\item
  Advocating for increased funding and resources to support research and development in this area, as well as the implementation of bias mitigation strategies in real-world applications.
\end{enumerate}

Ultimately, addressing the challenge of data bias in large multi-modal language models is a complex and ongoing endeavor that requires a sustained and multifaceted effort. By combining technical innovations with societal engagement and collaboration, researchers and developers can work towards creating more equitable and inclusive language models that serve the diverse needs of all individuals and communities \cite{ref187}.

\subsection{Safety Concerns}

The development of large multi-modal language models (LMLMs) has raised significant safety concerns, as these powerful models have the potential to generate inappropriate, misleading, or even dangerous content. One of the primary safety challenges is the risk of LMLMs perpetuating discrimination, exclusion, and toxicity by generating content that reinforces harmful stereotypes or promotes unethical behavior \cite{ref103}. The model's ability to accurately infer sensitive information, such as personal data, also poses a risk of privacy violations. Furthermore, LMLMs can contribute to the spread of misinformation and the erosion of trust in shared information, particularly in sensitive domains like healthcare or finance \cite{ref103}.

Another significant safety concern is the potential for malicious actors to exploit LMLMs for harmful purposes, such as generating content to incite violence, facilitate fraud, or enable the creation of malware \cite{ref325}. The inherent instruction-following capabilities of LMLMs make them vulnerable to jailbreak attacks, where carefully crafted prompts can elicit harmful responses, even from models that have been aligned for safety \cite{ref326}.

Furthermore, the potential risks of multi-turn dialogues, where LMLMs may not intend to reject cautionary or borderline unsafe queries, leading to the incremental generation of harmful content, underscores the need for comprehensive safety measures that can effectively mitigate risks across various interaction scenarios \cite{ref327}.

The safety concerns associated with LMLMs also extend to the human-computer interaction aspect, as the use of these models in conversational agents or other user-facing applications could lead to manipulation, deception, or unsafe interactions, particularly for vulnerable populations \cite{ref103}.

Addressing these safety challenges is crucial for the responsible development and deployment of LMLMs. Researchers have proposed various mitigation strategies, including the development of comprehensive safety benchmarks and evaluation frameworks \cite{ref328}, as well as the concept of safety detectors, which can serve as efficient and reliable alternatives to directly imposing safety constraints on deployed models \cite{ref104}.

Furthermore, researchers have explored techniques for enhancing the safety alignment of LMLMs, such as fine-tuning models on datasets that prioritize safety \cite{ref329}, incorporating prompting strategies to encourage safe behavior \cite{ref330}, and leveraging specialized training approaches like self-refinement \cite{ref326}. However, as highlighted in the \cite{ref331} paper, alignment alone may not be sufficient to prevent LMLMs from generating harmful information.

Addressing these challenges will require a multifaceted approach, including the development of more sophisticated safety evaluation frameworks, the integration of robust security measures to mitigate the risk of malicious exploitation, and the establishment of ethical governance frameworks to guide the responsible development and deployment of LMLMs \cite{ref277}. As the field of LMLMs continues to evolve, it is crucial that the research community, industry, and policymakers work together to address these safety concerns and ensure that these powerful models are developed and utilized in a way that maximizes their benefits while minimizing their potential for harm.

\subsection{Interpretability and Transparency}

The interpretability and transparency of large multi-modal language models (LMLMs) remain significant challenges that must be addressed to ensure their safe and responsible deployment. The internal decision-making processes of these complex models are often opaque, making it difficult to understand how they arrive at their outputs \cite{ref332}.

One of the key issues with LMLMs is their tendency to generate plausible-sounding but inaccurate or misleading content, a phenomenon known as ``hallucination'' \cite{ref333}. This lack of transparency can undermine trust in the model's outputs, particularly in high-stakes applications where decisions have significant real-world consequences.

Efforts to improve the interpretability of LMLMs have focused on developing techniques that can shed light on their internal decision-making processes. Approaches such as attention visualization, which highlights the parts of the input that the model is focusing on, and concept-based analyses, which identify the key concepts that the model is utilizing, have provided some insights \cite{ref332}. However, these methods often focus on either local or global explanations within a single dimension, and they may still fall short in providing a comprehensive understanding of the model's reasoning \cite{ref183}.

To address these limitations, researchers have proposed more holistic approaches to model interpretability. One such approach is the use of sparsity-guided techniques, which aim to elucidate the model's behavior at multiple levels, including the input, subnetwork, and concept levels \cite{ref183}. These techniques can provide a more comprehensive understanding of the model's reasoning, enabling better identification and mitigation of issues such as hallucinations.

Another promising approach for improving the transparency of LMLMs is the incorporation of interpretable inference-time intervention. This involves the ability to dynamically adjust the model's behavior during deployment, based on an understanding of its internal decision-making processes \cite{ref183}. By providing a more direct mechanism for controlling the model's outputs, this approach can enhance the accountability and trustworthiness of these systems.

Additionally, research has explored the use of large language models themselves as post hoc explainers for other complex predictive models \cite{ref334}. By leveraging the in-context learning capabilities of LLMs, researchers have developed approaches that can explain the predictions made by other models, potentially offering a more accessible and scalable solution for model interpretability.

However, the challenges of achieving interpretability and transparency in LMLMs go beyond just developing technical solutions. There are also significant ethical and societal implications that must be considered. For example, the potential for these models to be misused or to perpetuate biases and discrimination is a critical concern that requires careful attention \cite{ref233}.

To address these challenges, researchers have proposed frameworks for auditing and governing the development and deployment of LMLMs \cite{ref233}. These frameworks involve a multi-layered approach, with governance audits of the technology providers, model audits of the LMLMs themselves, and application audits of the specific use cases. By considering the ethical and societal implications of these models at multiple levels, these frameworks aim to ensure their responsible development and deployment.

In conclusion, the interpretability and transparency of large multi-modal language models remain crucial challenges that must be addressed to realize the full potential of these transformative technologies. Continued research into techniques for understanding and controlling the internal decision-making processes of LMLMs, as well as a broader consideration of the ethical and societal implications of these systems, will be essential for ensuring their safe and responsible deployment in a wide range of applications.

\subsection{Scalability and Computational Constraints}

The development and deployment of large multi-modal language models (MMLMs) face significant scalability and computational challenges that must be addressed to harness their full potential. As these models continue to grow in size and complexity, the computational resources required for training and inference become increasingly prohibitive, posing barriers to their widespread adoption.

One of the primary scalability challenges is the immense memory and computational requirements of MMLMs. These models typically consist of billions or even trillions of parameters, which need to be stored and processed during both training and inference \cite{ref335}. The memory demands are further exacerbated by the integration of multiple modalities, such as text, images, and audio, which require additional storage and processing capabilities \cite{ref336}.

The computational complexity of MMLMs also poses a significant challenge. The attention mechanisms and transformer-based architectures employed in these models are computationally intensive, requiring substantial processing power and time to perform the necessary operations \cite{ref337}. This computational burden becomes increasingly problematic as the model size and the number of input modalities grow, leading to significant latency and reduced throughput during inference.

To address these scalability and computational challenges, researchers have explored various strategies and solutions. One promising approach is the development of efficient hardware accelerators specifically designed for MMLMs \cite{ref7}. These specialized hardware solutions, such as custom-designed chips or accelerators, can significantly improve the performance and energy efficiency of MMLMs by leveraging parallel processing, reduced memory requirements, and tailored architectures \cite{ref338}.

Another strategy is the exploration of parameter-efficient fine-tuning (PEFT) techniques, which aim to reduce the number of trainable parameters in MMLMs without compromising their performance \cite{ref339}. By employing techniques like adapter modules, prompt tuning, or sparse attention mechanisms, researchers have demonstrated the ability to adapt MMLMs to specific tasks or domains while significantly reducing the computational and memory requirements \cite{ref340}.

Additionally, the research community is investigating model compression and distillation techniques to reduce the size and complexity of MMLMs without sacrificing their capabilities \cite{ref341}. These methods, which include techniques like knowledge distillation, weight pruning, and quantization, can enable the deployment of MMLMs on resource-constrained devices or edge environments \cite{ref342}.

Furthermore, the integration of efficient multi-modal reasoning and inference mechanisms, such as task-specific adapters or in-context learning, can help mitigate the computational burden of MMLMs \cite{ref108}. By leveraging these techniques, the model can focus on the most relevant information and selectively process the input modalities, leading to improved efficiency and performance.

Despite these efforts, the scalability and computational constraints of MMLMs remain a significant challenge, particularly as the models continue to grow in size and complexity. Addressing these challenges will require a multifaceted approach, including the development of novel hardware solutions, efficient training and inference algorithms, and effective model compression techniques. Continued research and collaboration between the AI community, hardware manufacturers, and industry partners will be crucial in overcoming these barriers and enabling the widespread deployment of large multi-modal language models across a wide range of applications \cite{ref337}.

\subsection{Future Research Directions}

The rapid advancements in large multi-modal language models (LMMsLM), as showcased in the preceding sections of this survey, have opened up a wide array of exciting future research directions. As the field continues to progress, researchers will need to address several key challenges and limitations to fully realize the potential of these transformative technologies.

One promising avenue for future research is the further enhancement of multimodal reasoning capabilities. While LMMsLM have demonstrated impressive abilities in integrating visual, textual, and other modalities, there is still room for improvement in their capacity to deeply understand and reason across these diverse inputs \cite{ref78}. Techniques such as prompting and in-context learning have shown promise in boosting multimodal reasoning \cite{ref94}, and further advancements in these areas could unlock even more powerful multimodal intelligence.

Additionally, the development of more robust and generalizable LMMsLM architectures is a crucial area for future research. Current models often struggle with domain-specific challenges and limitations, and the ability to build models that can seamlessly adapt to a wide range of tasks and environments would be transformative \cite{ref343}. Exploring novel architectural designs, integration methods, and training strategies that can foster more versatile and adaptable LMMsLM is a pressing research priority.

Another important direction for future research is the enhancement of LMMsLM's safety, reliability, and interpretability. As these models become more influential in various domains, including healthcare, finance, and education, it is crucial to address concerns related to data bias, model misuse, and the potential for harmful outputs \cite{ref105}. Developing robust evaluation frameworks, interpretability techniques, and safety-oriented training approaches could help mitigate these challenges and ensure the responsible development and deployment of LMMsLM \cite{ref344}.

The integration of LMMsLM with other AI technologies, such as robotics and reinforcement learning, presents another exciting avenue for future research \cite{ref345}. By combining the language understanding and generation capabilities of LMMsLM with the physical interaction and decision-making abilities of other AI systems, researchers could create powerful hybrid agents capable of tackling complex, real-world problems \cite{ref109}. Additionally, the exploration of LMMsLM in specialized domains, such as scientific research, education, and sustainability, could lead to transformative breakthroughs \cite{ref346,ref347,ref110}.

Furthermore, the continued development of multimodal datasets and benchmarks will be crucial for advancing LMMsLM research \cite{ref48}. As the field evolves, the creation of high-quality, diverse, and challenging datasets that capture the nuances of real-world multimodal interactions will be essential for driving progress and ensuring the robust evaluation of these models.

Finally, the ethical and societal implications of LMMsLM must be rigorously explored and addressed \cite{ref91,ref348}. As these models become more ubiquitous, it is crucial to develop comprehensive frameworks for ensuring their responsible development and deployment, mitigating risks related to privacy, fairness, and potential misuse. Engaging with a wide range of stakeholders, including policymakers, industry leaders, and the broader public, will be crucial for shaping the future of LMMsLM in a manner that maximizes societal benefits while minimizing potential harms.

\section{Conclusion and Future Outlook}

\subsection{Emerging Trends in Large Multi-Modal Language Models}

The field of large multi-modal language models (LMLMs) has witnessed a rapid surge of advancements in recent years, with researchers and practitioners alike pushing the boundaries of what is possible in the realm of general-purpose artificial intelligence. As this dynamic field continues to evolve, several emerging trends have the potential to shape the future trajectory of LMLMs.

One of the most prominent trends is the development of increasingly sophisticated architectural designs that enable seamless integration and interaction between diverse modalities. A prime example is the Multiway Transformers introduced in \cite{ref139}, which leverages a modular architecture to facilitate both deep fusion and modality-specific encoding. This innovative approach allows for the effective capture of cross-modal relationships while preserving the unique characteristics of individual modalities, paving the way for more robust and versatile LMLMs.

Another key trend is the exploration of novel training strategies and objectives that aim to enhance the multimodal understanding and reasoning capabilities of these models. The concept of ``Imglish,'' as presented in \cite{ref139}, wherein images are treated as a ``foreign language'' to be learned alongside text, represents a radical departure from traditional training approaches. By unifying the learning of various modalities under a common framework, this strategy enables LMLMs to develop a more holistic understanding of the world, potentially unlocking new avenues for cross-modal reasoning and generalization.

Furthermore, researchers are increasingly focusing on the development of specialized datasets and benchmarks that can effectively assess the multimodal capabilities of LMLMs. The introduction of the \cite{ref49} exemplifies this trend, providing a comprehensive platform for evaluating the spatial understanding and expressive abilities of these models across a wide range of 3D tasks. Such targeted evaluation frameworks are crucial for driving the field forward, as they enable the precise identification of strengths, weaknesses, and areas for improvement in the design and training of LMLMs.

Closely tied to the advancement of benchmarking efforts is the growing emphasis on enhancing the interpretability and transparency of LMLMs. The development of tools like \cite{ref29} aims to shed light on the internal mechanisms and decision-making processes of these models, enabling a deeper understanding of their strengths and limitations. As the capabilities of LMLMs continue to expand, the need for such interpretability tools will become increasingly important, as they can inform the design of more robust and trustworthy models.

Alongside these architectural and training-centric trends, researchers are also exploring the application of LMLMs in diverse domains, pushing the boundaries of what these models can accomplish. The introduction of \cite{ref349} and \cite{ref350} highlights the potential of LMLMs to tackle domain-specific challenges, demonstrating their versatility and adaptability beyond general-purpose tasks.

Furthermore, the concept of ``modality plug-and-play,'' as explored in \cite{ref340}, represents a promising direction for enhancing the deployability and scalability of LMLMs. By enabling the adaptive integration of different modalities at runtime, this approach addresses the challenges of incorporating an ever-expanding array of input modalities, particularly in resource-constrained environments such as edge devices.

Looking ahead, the continued advancements in large multi-modal language models are poised to have a transformative impact on the field of artificial intelligence. As researchers and practitioners continue to push the boundaries of what is possible, we can expect to see the emergence of even more powerful and versatile models capable of seamlessly bridging the gaps between various modalities and tackling increasingly complex real-world challenges. The future of LMLMs holds immense promise, with the potential to unlock new frontiers in human-machine interaction, decision-making, and the pursuit of artificial general intelligence.

\subsection{Towards Generalized Multi-Modal Intelligence}

The emergence of large multi-modal language models (LMLMs) holds immense promise in advancing the quest for generalized multi-modal intelligence. These models, which seamlessly integrate diverse modalities such as text, vision, and audio, represent a significant step towards bridging the gap between artificial and human-like cognitive capabilities \cite{ref351,ref352,ref6}.

At the core of this endeavor lies the fundamental challenge of enabling machines to comprehend and navigate the richly interconnected, multimodal nature of the real world, much like how human beings perceive and interact with their environment. Unlike traditional unimodal language models, LMLMs possess the ability to process and reason about a wide range of information sources, from natural language to visual cues and audio signals \cite{ref353,ref221}.

The pursuit of generalized multi-modal intelligence involves several key research directions and opportunities. Firstly, the advancement of architectural designs and training strategies for LMLMs is crucial. Researchers have explored various approaches, such as modality-specific encoders, multi-modal fusion modules, and cross-modal attention mechanisms, to effectively integrate diverse modalities within a unified model \cite{ref6,ref354,ref17}. However, there remains a need to develop more efficient and scalable architectures that can seamlessly handle a growing array of modalities without compromising performance or computational efficiency.

Secondly, the curation and utilization of comprehensive multimodal datasets are essential for training and evaluating generalized multi-modal intelligence. Existing datasets have primarily focused on specific task-oriented scenarios, such as image-text understanding or audio-visual interaction \cite{ref353,ref351}. To truly foster the development of generalized multi-modal intelligence, researchers must prioritize the creation of diverse, open-ended datasets that capture the complexity and nuances of real-world multi-modal interactions \cite{ref355,ref356}.

Furthermore, the advancement of multi-modal reasoning and in-context learning techniques is crucial for enabling LMLMs to engage in more sophisticated problem-solving and task execution. Approaches like multimodal chain-of-thought prompting, which integrate language and visual information to enhance reasoning capabilities, have demonstrated promising results \cite{ref357,ref358}. However, further research is needed to develop robust and generalizable multi-modal reasoning frameworks that can handle a wide range of tasks and scenarios \cite{ref359,ref360}.

Another key area of exploration is the development of multimodal grounding and perceptual understanding within LMLMs. Current models often struggle to establish coherent connections between language and other modalities, leading to inconsistencies and failures in real-world applications \cite{ref361,ref26}. Advancing techniques for multimodal grounding, such as language-assisted vision understanding and modality-aware feature representations, can significantly improve the overall performance and reliability of generalized multi-modal intelligence \cite{ref362,ref16}.

Lastly, the ethical and societal implications of generalized multi-modal intelligence must be carefully considered. As these advanced models become more prevalent, concerns regarding privacy, fairness, and the potential for misuse must be addressed \cite{ref102}. Developing robust governance frameworks, transparency mechanisms, and responsible deployment strategies will be crucial in ensuring the safe and beneficial integration of generalized multi-modal intelligence into various domains \cite{ref363,ref364}.

In conclusion, the pursuit of generalized multi-modal intelligence represents a transformative frontier in artificial intelligence research. By seamlessly integrating and understanding diverse modalities, LMLMs hold the potential to bridge the gap between artificial and human-like cognitive capabilities, ushering in a new era of intelligent systems that can comprehend and interact with the world in a more naturalistic and holistic manner. The research challenges and opportunities outlined above highlight the exciting possibilities and the critical need for continued advancements in this rapidly evolving field.

\subsection{Ethical Considerations and Societal Impact}

The widespread adoption of large multi-modal language models raises significant ethical concerns that must be carefully considered to ensure their responsible development and deployment. One of the primary challenges is the potential for these models to perpetuate and amplify societal biases, as they are often trained on datasets that reflect existing inequalities \cite{ref229,ref88}. If these biases are not adequately addressed, they can be propagated and even exacerbated through the outputs of these models, posing serious risks to fairness and inclusivity \cite{ref87,ref365}.

Furthermore, the increasing integration of multi-modal language models into various applications, such as virtual assistants, content generation, and decision-making systems, raises concerns about privacy \cite{ref232,ref366}. These models have the potential to expose sensitive information about individuals, either through direct data leakage or by inferring personal characteristics from the generated content \cite{ref86,ref79}. Developing robust privacy safeguards and ensuring that users have control over their personal data is crucial to building trust in these technologies \cite{ref81}.

Additionally, the potential for misuse of large multi-modal language models is a significant concern. These models can be used to generate disinformation, impersonate individuals, or create content for malicious purposes \cite{ref103,ref367}. The ease with which these models can be accessed and the difficulty in distinguishing machine-generated content from human-generated content pose serious challenges in terms of accountability and the prevention of harmful activities \cite{ref89,ref91}.

To ensure the responsible development and deployment of large multi-modal language models, a multi-pronged approach is necessary. This includes the establishment of ethical governance frameworks and oversight mechanisms \cite{ref368} that can guide the development and deployment of these technologies in a way that prioritizes ethical principles, such as fairness, transparency, and accountability \cite{ref235,ref369}. Engaging with diverse stakeholders, including policymakers, industry leaders, researchers, and civil society organizations, is crucial to developing a shared understanding of the ethical implications and collaboratively devising solutions \cite{ref370,ref371}.

Furthermore, the responsible development and deployment of large multi-modal language models must be accompanied by efforts to address the broader societal impacts of these technologies \cite{ref372,ref373}. This includes understanding the potential effects on employment, education, and social dynamics, and developing strategies to ensure that the benefits of these technologies are distributed equitably \cite{ref288,ref283}.

\subsection{Future Outlook and Research Directions}

The rapid advancements in large multi-modal language models (LMLMs) have paved the way for a future where these models will play an increasingly pivotal role in our daily lives. As we look ahead, several emerging areas of research and potential application domains hold the promise to unlock the full transformative potential of this technology.

One promising area of research is the seamless integration of multi-modal reasoning capabilities within LMLMs. While current models have demonstrated impressive performance in tasks that involve the interplay between text and visual information \cite{ref6}, there remains significant room for improvement in their ability to reason about and comprehend the complex relationships between diverse modalities, including audio, video, and even tactile feedback. Researchers are actively exploring novel architectural designs and training strategies that can foster more robust and generalizable multi-modal reasoning abilities, enabling LMLMs to tackle increasingly complex real-world problems that require holistic understanding of the environment \cite{ref48}.

Another crucial area of focus is the enhancement of LMLMs' ability to engage in meaningful and natural interactions with humans, moving beyond the current limitations of text-based dialogues. The development of large language and speech models (LLaSMs) that can seamlessly integrate speech, vision, and language \cite{ref221} is a significant step in this direction, paving the way for more intuitive and immersive human-AI interactions. As these models become more advanced, they will be able to understand and respond to multimodal cues, facilitating more natural and contextual communication, and opening up new applications in areas such as personal assistants, virtual agents, and even robotics \cite{ref374}.

Closely related to the goal of natural interaction is the drive to develop LMLMs that can truly exhibit autonomous and goal-oriented behavior. The integration of LMLMs with reinforcement learning and other decision-making frameworks \cite{ref375} has the potential to create intelligent agents that can independently perceive, reason, and act in complex environments, moving closer to the vision of artificial general intelligence (AGI). By leveraging LMLMs' language understanding and generation capabilities, coupled with their ability to interact with and manipulate the physical world, these autonomous agents could revolutionize fields like scientific research, disaster response, and urban planning \cite{ref101,ref376}.

Furthermore, the application of LMLMs is poised to expand into an even broader range of domains, including healthcare \cite{ref193}, finance \cite{ref377}, and education \cite{ref110}. As these models become more specialized and tailored to specific industries and use cases, they will unlock novel solutions to complex problems, enhancing productivity, decision-making, and personalization across various sectors \cite{ref378}.

However, the rapid development and adoption of LMLMs also raise critical concerns that must be addressed. Ethical considerations, such as privacy, bias, and the potential for misuse, remain at the forefront of the research agenda \cite{ref102}. Developing robust frameworks for ensuring the responsible and transparent deployment of LMLMs, while preserving user trust and societal well-being, is a paramount challenge that the research community must tackle \cite{ref346}.

Additionally, the scalability and computational efficiency of LMLMs pose significant hurdles, as the ever-increasing model size and complexity strain computational resources and energy consumption \cite{ref336}. Innovative architectural designs, training techniques, and hardware optimizations will be crucial in addressing these challenges and enabling the widespread adoption of LMLMs \cite{ref354}.

In conclusion, the future of large multi-modal language models holds immense promise, with the potential to revolutionize a wide range of industries and transform the way we interact with technology. By addressing the key research challenges and ethical considerations, the research community can pave the way for the development of truly intelligent and adaptable systems that can seamlessly integrate multiple modalities, engage in natural and autonomous interactions, and deliver tailored solutions to complex real-world problems. As we continue to push the boundaries of what is possible with LMLMs, we stand at the cusp of a transformative era in the evolution of artificial intelligence.


\begin{thebibliography}{378}
\addcontentsline{toc}{section}{References}
\bibitem{ref1} Language Models are Few-Shot Learners. arXiv:2005.14165v4, 2020. \url{https://arxiv.org/abs/2005.14165v4}
\bibitem{ref2} BERT Pre-training of Deep Bidirectional Transformers for Language Understanding. arXiv:1810.04805v2, 2018. \url{https://arxiv.org/abs/1810.04805v2}
\bibitem{ref3} Exploring the Limits of Transfer Learning with a Unified Text-to-Text Transformer. arXiv:1910.10683v4, 2019. \url{https://arxiv.org/abs/1910.10683v4}
\bibitem{ref4} A Review of Multi-Modal Large Language and Vision Models. arXiv:2404.01322v1, 2024. \url{https://arxiv.org/abs/2404.01322v1}
\bibitem{ref5} MM-LLMs Recent Advances in MultiModal Large Language Models. arXiv:2401.13601v4, 2024. \url{https://arxiv.org/abs/2401.13601v4}
\bibitem{ref6} PaLM-E An Embodied Multimodal Language Model. arXiv:2303.03378v1, 2023. \url{https://arxiv.org/abs/2303.03378v1}
\bibitem{ref7} A Survey on Hardware Accelerators for Large Language Models. arXiv:2401.09890v1, 2024. \url{https://arxiv.org/abs/2401.09890v1}
\bibitem{ref8} The Curious Case of Nonverbal Abstract Reasoning with Multi-Modal Large Language Models. arXiv:2401.12117v2, 2024. \url{https://arxiv.org/abs/2401.12117v2}
\bibitem{ref9} AudioGPT Understanding and Generating Speech, Music, Sound, and Talking Head. arXiv:2304.12995v1, 2023. \url{https://arxiv.org/abs/2304.12995v1}
\bibitem{ref10} UnIVAL Unified Model for Image, Video, Audio and Language Tasks. arXiv:2307.16184v2, 2023. \url{https://arxiv.org/abs/2307.16184v2}
\bibitem{ref11} An Interdisciplinary Outlook on Large Language Models for Scientific Research. arXiv:2311.04929v1, 2023. \url{https://arxiv.org/abs/2311.04929v1}
\bibitem{ref12} Large Language Models for Telecom Forthcoming Impact on the Industry. arXiv:2308.06013v2, 2023. \url{https://arxiv.org/abs/2308.06013v2}
\bibitem{ref13} LLeMpower Understanding Disparities in the Control and Access of Large Language Models. arXiv:2404.09356v1, 2024. \url{https://arxiv.org/abs/2404.09356v1}
\bibitem{ref14} On the Ethical Considerations of Text Simplification. arXiv:2204.09565v1, 2022. \url{https://arxiv.org/abs/2204.09565v1}
\bibitem{ref15} Exploring the Frontier of Vision-Language Models A Survey of Current Methodologies and Future Directions. arXiv:2404.07214v2, 2024. \url{https://arxiv.org/abs/2404.07214v2}
\bibitem{ref16} ModaVerse Efficiently Transforming Modalities with LLMs. arXiv:2401.06395v2, 2024. \url{https://arxiv.org/abs/2401.06395v2}
\bibitem{ref17} ChatBridge Bridging Modalities with Large Language Model as a Language Catalyst. arXiv:2305.16103v1, 2023. \url{https://arxiv.org/abs/2305.16103v1}
\bibitem{ref18} TEAL Tokenize and Embed ALL for Multi-modal Large Language Models. arXiv:2311.04589v3, 2023. \url{https://arxiv.org/abs/2311.04589v3}
\bibitem{ref19} GIT-Mol A Multi-modal Large Language Model for Molecular Science with Graph, Image, and Text. arXiv:2308.06911v3, 2023. \url{https://arxiv.org/abs/2308.06911v3}
\bibitem{ref20} SpeechGPT Empowering Large Language Models with Intrinsic Cross-Modal Conversational Abilities. arXiv:2305.11000v2, 2023. \url{https://arxiv.org/abs/2305.11000v2}
\bibitem{ref21} FSMR A Feature Swapping Multi-modal Reasoning Approach with Joint Textual and Visual Clues. arXiv:2403.20026v1, 2024. \url{https://arxiv.org/abs/2403.20026v1}
\bibitem{ref22} VATLM Visual-Audio-Text Pre-Training with Unified Masked Prediction for Speech Representation Learning. arXiv:2211.11275v2, 2022. \url{https://arxiv.org/abs/2211.11275v2}
\bibitem{ref23} Multimodal Machine Learning A Survey and Taxonomy. arXiv:1705.09406v2, 2017. \url{https://arxiv.org/abs/1705.09406v2}
\bibitem{ref24} Multimodal Transformer for Unaligned Multimodal Language Sequences. arXiv:1906.00295v1, 2019. \url{https://arxiv.org/abs/1906.00295v1}
\bibitem{ref25} Switch-BERT Learning to Model Multimodal Interactions by Switching Attention and Input. arXiv:2306.14182v1, 2023. \url{https://arxiv.org/abs/2306.14182v1}
\bibitem{ref26} GroundingGPT Language Enhanced Multi-modal Grounding Model. arXiv:2401.06071v5, 2024. \url{https://arxiv.org/abs/2401.06071v5}
\bibitem{ref27} FLAVA A Foundational Language And Vision Alignment Model. arXiv:2112.04482v3, 2021. \url{https://arxiv.org/abs/2112.04482v3}
\bibitem{ref28} TriBERT Full-body Human-centric Audio-visual Representation Learning for Visual Sound Separation. arXiv:2110.13412v1, 2021. \url{https://arxiv.org/abs/2110.13412v1}
\bibitem{ref29} LVLM-Intrepret An Interpretability Tool for Large Vision-Language Models. arXiv:2404.03118v1, 2024. \url{https://arxiv.org/abs/2404.03118v1}
\bibitem{ref30} PILL Plug Into LLM with Adapter Expert and Attention Gate. arXiv:2311.02126v1, 2023. \url{https://arxiv.org/abs/2311.02126v1}
\bibitem{ref31} Large-scale Multi-Modal Pre-trained Models A Comprehensive Survey. arXiv:2302.10035v3, 2023. \url{https://arxiv.org/abs/2302.10035v3}
\bibitem{ref32} Jointly Training Large Autoregressive Multimodal Models. arXiv:2309.15564v2, 2023. \url{https://arxiv.org/abs/2309.15564v2}
\bibitem{ref33} LANISTR Multimodal Learning from Structured and Unstructured Data. arXiv:2305.16556v3, 2023. \url{https://arxiv.org/abs/2305.16556v3}
\bibitem{ref34} PaLM Scaling Language Modeling with Pathways. arXiv:2204.02311v5, 2022. \url{https://arxiv.org/abs/2204.02311v5}
\bibitem{ref35} Different Strokes for Different Folks Investigating Appropriate Further Pre-training Approaches for Diverse Dialogue Tasks. arXiv:2109.06524v1, 2021. \url{https://arxiv.org/abs/2109.06524v1}
\bibitem{ref36} PaCE Unified Multi-modal Dialogue Pre-training with Progressive and Compositional Experts. arXiv:2305.14839v2, 2023. \url{https://arxiv.org/abs/2305.14839v2}
\bibitem{ref37} Continual Instruction Tuning for Large Multimodal Models. arXiv:2311.16206v1, 2023. \url{https://arxiv.org/abs/2311.16206v1}
\bibitem{ref38} CAPT Contrastive Pre-Training for Learning Denoised Sequence Representations. arXiv:2010.06351v4, 2020. \url{https://arxiv.org/abs/2010.06351v4}
\bibitem{ref39} Muppet Massive Multi-task Representations with Pre-Finetuning. arXiv:2101.11038v1, 2021. \url{https://arxiv.org/abs/2101.11038v1}
\bibitem{ref40} MAGMA -- Multimodal Augmentation of Generative Models through Adapter-based Finetuning. arXiv:2112.05253v2, 2021. \url{https://arxiv.org/abs/2112.05253v2}
\bibitem{ref41} A Survey of Vision-Language Pre-training from the Lens of Multimodal Machine Translation. arXiv:2306.07198v1, 2023. \url{https://arxiv.org/abs/2306.07198v1}
\bibitem{ref42} mSLAM Massively multilingual joint pre-training for speech and text. arXiv:2202.01374v1, 2022. \url{https://arxiv.org/abs/2202.01374v1}
\bibitem{ref43} Vision Learners Meet Web Image-Text Pairs. arXiv:2301.07088v2, 2023. \url{https://arxiv.org/abs/2301.07088v2}
\bibitem{ref44} Revisiting Pre-training in Audio-Visual Learning. arXiv:2302.03533v2, 2023. \url{https://arxiv.org/abs/2302.03533v2}
\bibitem{ref45} MM1 Methods, Analysis \& Insights from Multimodal LLM Pre-training. arXiv:2403.09611v4, 2024. \url{https://arxiv.org/abs/2403.09611v4}
\bibitem{ref46} SEED-Bench Benchmarking Multimodal LLMs with Generative Comprehension. arXiv:2307.16125v2, 2023. \url{https://arxiv.org/abs/2307.16125v2}
\bibitem{ref47} WIT Wikipedia-based Image Text Dataset for Multimodal Multilingual Machine Learning. arXiv:2103.01913v2, 2021. \url{https://arxiv.org/abs/2103.01913v2}
\bibitem{ref48} LAMM Language-Assisted Multi-Modal Instruction-Tuning Dataset, Framework, and Benchmark. arXiv:2306.06687v3, 2023. \url{https://arxiv.org/abs/2306.06687v3}
\bibitem{ref49} 3DBench A Scalable 3D Benchmark and Instruction-Tuning Dataset. arXiv:2404.14678v1, 2024. \url{https://arxiv.org/abs/2404.14678v1}
\bibitem{ref50} MLLM-DataEngine An Iterative Refinement Approach for MLLM. arXiv:2308.13566v2, 2023. \url{https://arxiv.org/abs/2308.13566v2}
\bibitem{ref51} MM-BigBench Evaluating Multimodal Models on Multimodal Content Comprehension Tasks. arXiv:2310.09036v1, 2023. \url{https://arxiv.org/abs/2310.09036v1}
\bibitem{ref52} MULTI Multimodal Understanding Leaderboard with Text and Images. arXiv:2402.03173v2, 2024. \url{https://arxiv.org/abs/2402.03173v2}
\bibitem{ref53} WanJuan A Comprehensive Multimodal Dataset for Advancing English and Chinese Large Models. arXiv:2308.10755v3, 2023. \url{https://arxiv.org/abs/2308.10755v3}
\bibitem{ref54} ChEF A Comprehensive Evaluation Framework for Standardized Assessment of Multimodal Large Language Models. arXiv:2311.02692v1, 2023. \url{https://arxiv.org/abs/2311.02692v1}
\bibitem{ref55} MMBench Is Your Multi-modal Model an All-around Player. arXiv:2307.06281v3, 2023. \url{https://arxiv.org/abs/2307.06281v3}
\bibitem{ref56} Holistic Evaluation of Language Models. arXiv:2211.09110v2, 2022. \url{https://arxiv.org/abs/2211.09110v2}
\bibitem{ref57} Evaluating Large Language Models A Comprehensive Survey. arXiv:2310.19736v3, 2023. \url{https://arxiv.org/abs/2310.19736v3}
\bibitem{ref58} A User-Centric Benchmark for Evaluating Large Language Models. arXiv:2404.13940v2, 2024. \url{https://arxiv.org/abs/2404.13940v2}
\bibitem{ref59} METAL Metamorphic Testing Framework for Analyzing Large-Language Model Qualities. arXiv:2312.06056v1, 2023. \url{https://arxiv.org/abs/2312.06056v1}
\bibitem{ref60} Can Large Language Models be Trusted for Evaluation Scalable Meta-Evaluation of LLMs as Evaluators via Agent Debate. arXiv:2401.16788v1, 2024. \url{https://arxiv.org/abs/2401.16788v1}
\bibitem{ref61} Are Large Language Model-based Evaluators the Solution to Scaling Up Multilingual Evaluation. arXiv:2309.07462v2, 2023. \url{https://arxiv.org/abs/2309.07462v2}
\bibitem{ref62} Crossing Linguistic Horizons Finetuning and Comprehensive Evaluation of Vietnamese Large Language Models. arXiv:2403.02715v1, 2024. \url{https://arxiv.org/abs/2403.02715v1}
\bibitem{ref63} SeaEval for Multilingual Foundation Models From Cross-Lingual Alignment to Cultural Reasoning. arXiv:2309.04766v4, 2023. \url{https://arxiv.org/abs/2309.04766v4}
\bibitem{ref64} Peacock A Family of Arabic Multimodal Large Language Models and Benchmarks. arXiv:2403.01031v1, 2024. \url{https://arxiv.org/abs/2403.01031v1}
\bibitem{ref65} Beyond Static Datasets A Deep Interaction Approach to LLM Evaluation. arXiv:2309.04369v1, 2023. \url{https://arxiv.org/abs/2309.04369v1}
\bibitem{ref66} Evaluating the Elementary Multilingual Capabilities of Large Language Models with MultiQ. arXiv:2403.03814v1, 2024. \url{https://arxiv.org/abs/2403.03814v1}
\bibitem{ref67} PATCH -- Psychometrics-AssisTed benCHmarking of Large Language Models A Case Study of Mathematics Proficiency. arXiv:2404.01799v1, 2024. \url{https://arxiv.org/abs/2404.01799v1}
\bibitem{ref68} CriticBench Evaluating Large Language Models as Critic. arXiv:2402.13764v3, 2024. \url{https://arxiv.org/abs/2402.13764v3}
\bibitem{ref69} Large Language Models are Few-Shot Health Learners. arXiv:2305.15525v1, 2023. \url{https://arxiv.org/abs/2305.15525v1}
\bibitem{ref70} Almanac Retrieval-Augmented Language Models for Clinical Medicine. arXiv:2303.01229v2, 2023. \url{https://arxiv.org/abs/2303.01229v2}
\bibitem{ref71} Redefining Digital Health Interfaces with Large Language Models. arXiv:2310.03560v3, 2023. \url{https://arxiv.org/abs/2310.03560v3}
\bibitem{ref72} ClinicalGPT Large Language Models Finetuned with Diverse Medical Data and Comprehensive Evaluation. arXiv:2306.09968v1, 2023. \url{https://arxiv.org/abs/2306.09968v1}
\bibitem{ref73} Elucidating STEM Concepts through Generative AI A Multi-modal Exploration of Analogical Reasoning. arXiv:2308.10454v1, 2023. \url{https://arxiv.org/abs/2308.10454v1}
\bibitem{ref74} Taking the Next Step with Generative Artificial Intelligence The Transformative Role of Multimodal Large Language Models in Science Education. arXiv:2401.00832v1, 2024. \url{https://arxiv.org/abs/2401.00832v1}
\bibitem{ref75} From Text to Transformation A Comprehensive Review of Large Language Models' Versatility. arXiv:2402.16142v1, 2024. \url{https://arxiv.org/abs/2402.16142v1}
\bibitem{ref76} LLMs-Healthcare Current Applications and Challenges of Large Language Models in various Medical Specialties. arXiv:2311.12882v3, 2023. \url{https://arxiv.org/abs/2311.12882v3}
\bibitem{ref77} The Impact of ChatGPT and LLMs on Medical Imaging Stakeholders Perspectives and Use Cases. arXiv:2306.06767v2, 2023. \url{https://arxiv.org/abs/2306.06767v2}
\bibitem{ref78} Benchmarking Sequential Visual Input Reasoning and Prediction in Multimodal Large Language Models. arXiv:2310.13473v1, 2023. \url{https://arxiv.org/abs/2310.13473v1}
\bibitem{ref79} Beyond the Safeguards Exploring the Security Risks of ChatGPT. arXiv:2305.08005v1, 2023. \url{https://arxiv.org/abs/2305.08005v1}
\bibitem{ref80} Human-Centered Privacy Research in the Age of Large Language Models. arXiv:2402.01994v1, 2024. \url{https://arxiv.org/abs/2402.01994v1}
\bibitem{ref81} The Ethics of Interaction Mitigating Security Threats in LLMs. arXiv:2401.12273v1, 2024. \url{https://arxiv.org/abs/2401.12273v1}
\bibitem{ref82} Protected group bias and stereotypes in Large Language Models. arXiv:2403.14727v1, 2024. \url{https://arxiv.org/abs/2403.14727v1}
\bibitem{ref83} She had Cobalt Blue Eyes Prompt Testing to Create Aligned and Sustainable Language Models. arXiv:2310.18333v3, 2023. \url{https://arxiv.org/abs/2310.18333v3}
\bibitem{ref84} Informed AI Regulation Comparing the Ethical Frameworks of Leading LLM Chatbots Using an Ethics-Based Audit to Assess Moral Reasoning and Normative Values. arXiv:2402.01651v1, 2024. \url{https://arxiv.org/abs/2402.01651v1}
\bibitem{ref85} Red teaming ChatGPT via Jailbreaking Bias, Robustness, Reliability and Toxicity. arXiv:2301.12867v4, 2023. \url{https://arxiv.org/abs/2301.12867v4}
\bibitem{ref86} Fortifying Ethical Boundaries in AI Advanced Strategies for Enhancing Security in Large Language Models. arXiv:2402.01725v1, 2024. \url{https://arxiv.org/abs/2402.01725v1}
\bibitem{ref87} Ethical ChatGPT Concerns, Challenges, and Commandments. arXiv:2305.10646v1, 2023. \url{https://arxiv.org/abs/2305.10646v1}
\bibitem{ref88} Mapping the Ethics of Generative AI A Comprehensive Scoping Review. arXiv:2402.08323v1, 2024. \url{https://arxiv.org/abs/2402.08323v1}
\bibitem{ref89} Vox Populi, Vox ChatGPT Large Language Models, Education and Democracy. arXiv:2311.06207v1, 2023. \url{https://arxiv.org/abs/2311.06207v1}
\bibitem{ref90} The teachers are confused as well A Multiple-Stakeholder Ethics Discussion on Large Language Models in Computing Education. arXiv:2401.12453v1, 2024. \url{https://arxiv.org/abs/2401.12453v1}
\bibitem{ref91} Voluminous yet Vacuous Semantic Capital in an Age of Large Language Models. arXiv:2306.01773v1, 2023. \url{https://arxiv.org/abs/2306.01773v1}
\bibitem{ref92} Bidirectional Language Models Are Also Few-shot Learners. arXiv:2209.14500v2, 2022. \url{https://arxiv.org/abs/2209.14500v2}
\bibitem{ref93} Speak Like a Native Prompting Large Language Models in a Native Style. arXiv:2311.13538v3, 2023. \url{https://arxiv.org/abs/2311.13538v3}
\bibitem{ref94} Multimodal Chain-of-Thought Reasoning in Language Models. arXiv:2302.00923v4, 2023. \url{https://arxiv.org/abs/2302.00923v4}
\bibitem{ref95} Understanding and Improving In-Context Learning on Vision-language Models. arXiv:2311.18021v1, 2023. \url{https://arxiv.org/abs/2311.18021v1}
\bibitem{ref96} KAM-CoT Knowledge Augmented Multimodal Chain-of-Thoughts Reasoning. arXiv:2401.12863v1, 2024. \url{https://arxiv.org/abs/2401.12863v1}
\bibitem{ref97} ToolkenGPT Augmenting Frozen Language Models with Massive Tools via Tool Embeddings. arXiv:2305.11554v4, 2023. \url{https://arxiv.org/abs/2305.11554v4}
\bibitem{ref98} MMICT Boosting Multi-Modal Fine-Tuning with In-Context Examples. arXiv:2312.06363v2, 2023. \url{https://arxiv.org/abs/2312.06363v2}
\bibitem{ref99} Retrieval-augmented Multi-modal Chain-of-Thoughts Reasoning for Large Language Models. arXiv:2312.01714v2, 2023. \url{https://arxiv.org/abs/2312.01714v2}
\bibitem{ref100} Bias and Fairness in Large Language Models A Survey. arXiv:2309.00770v2, 2023. \url{https://arxiv.org/abs/2309.00770v2}
\bibitem{ref101} Opportunities and Challenges of Applying Large Language Models in Building Energy Efficiency and Decarbonization Studies An Exploratory Overview. arXiv:2312.11701v1, 2023. \url{https://arxiv.org/abs/2312.11701v1}
\bibitem{ref102} Unbridled Icarus A Survey of the Potential Perils of Image Inputs in Multimodal Large Language Model Security. arXiv:2404.05264v1, 2024. \url{https://arxiv.org/abs/2404.05264v1}
\bibitem{ref103} Ethical and social risks of harm from Language Models. arXiv:2112.04359v1, 2021. \url{https://arxiv.org/abs/2112.04359v1}
\bibitem{ref104} Detectors for Safe and Reliable LLMs Implementations, Uses, and Limitations. arXiv:2403.06009v1, 2024. \url{https://arxiv.org/abs/2403.06009v1}
\bibitem{ref105} Securing Large Language Models Threats, Vulnerabilities and Responsible Practices. arXiv:2403.12503v1, 2024. \url{https://arxiv.org/abs/2403.12503v1}
\bibitem{ref106} Towards Consistent Language Models Using Declarative Constraints. arXiv:2312.15472v1, 2023. \url{https://arxiv.org/abs/2312.15472v1}
\bibitem{ref107} Evaluation and Enhancement of Semantic Grounding in Large Vision-Language Models. arXiv:2309.04041v2, 2023. \url{https://arxiv.org/abs/2309.04041v2}
\bibitem{ref108} Enhancing Human-like Multi-Modal Reasoning A New Challenging Dataset and Comprehensive Framework. arXiv:2307.12626v2, 2023. \url{https://arxiv.org/abs/2307.12626v2}
\bibitem{ref109} Large Language Models are Zero Shot Hypothesis Proposers. arXiv:2311.05965v1, 2023. \url{https://arxiv.org/abs/2311.05965v1}
\bibitem{ref110} Large Language Models for Education A Survey and Outlook. arXiv:2403.18105v2, 2024. \url{https://arxiv.org/abs/2403.18105v2}
\bibitem{ref111} The Limitations of Optimization from Samples. arXiv:1512.06238v3, 2015. \url{https://arxiv.org/abs/1512.06238v3}
\bibitem{ref112} An Empirical Study of Training End-to-End Vision-and-Language Transformers. arXiv:2111.02387v3, 2021. \url{https://arxiv.org/abs/2111.02387v3}
\bibitem{ref113} BridgeTower Building Bridges Between Encoders in Vision-Language Representation Learning. arXiv:2206.08657v6, 2022. \url{https://arxiv.org/abs/2206.08657v6}
\bibitem{ref114} All-in-One Image-Grounded Conversational Agents. arXiv:1912.12394v2, 2019. \url{https://arxiv.org/abs/1912.12394v2}
\bibitem{ref115} Encoder Fusion Network with Co-Attention Embedding for Referring Image Segmentation. arXiv:2105.01839v1, 2021. \url{https://arxiv.org/abs/2105.01839v1}
\bibitem{ref116} Question Aware Vision Transformer for Multimodal Reasoning. arXiv:2402.05472v1, 2024. \url{https://arxiv.org/abs/2402.05472v1}
\bibitem{ref117} UNITER UNiversal Image-TExt Representation Learning. arXiv:1909.11740v3, 2019. \url{https://arxiv.org/abs/1909.11740v3}
\bibitem{ref118} Semantics-enhanced Cross-modal Masked Image Modeling for Vision-Language Pre-training. arXiv:2403.00249v1, 2024. \url{https://arxiv.org/abs/2403.00249v1}
\bibitem{ref119} Masked Contrastive Reconstruction for Cross-modal Medical Image-Report Retrieval. arXiv:2312.15840v2, 2023. \url{https://arxiv.org/abs/2312.15840v2}
\bibitem{ref120} Vision-and-Language or Vision-for-Language On Cross-Modal Influence in Multimodal Transformers. arXiv:2109.04448v1, 2021. \url{https://arxiv.org/abs/2109.04448v1}
\bibitem{ref121} Global and Local Semantic Completion Learning for Vision-Language Pre-training. arXiv:2306.07096v2, 2023. \url{https://arxiv.org/abs/2306.07096v2}
\bibitem{ref122} Dense Contrastive Visual-Linguistic Pretraining. arXiv:2109.11778v1, 2021. \url{https://arxiv.org/abs/2109.11778v1}
\bibitem{ref123} Vision-Language Pre-Training for Boosting Scene Text Detectors. arXiv:2204.13867v1, 2022. \url{https://arxiv.org/abs/2204.13867v1}
\bibitem{ref124} CTAL Pre-training Cross-modal Transformer for Audio-and-Language Representations. arXiv:2109.00181v1, 2021. \url{https://arxiv.org/abs/2109.00181v1}
\bibitem{ref125} Leveraging per Image-Token Consistency for Vision-Language Pre-training. arXiv:2211.15398v2, 2022. \url{https://arxiv.org/abs/2211.15398v2}
\bibitem{ref126} ROSITA Enhancing Vision-and-Language Semantic Alignments via Cross- and Intra-modal Knowledge Integration. arXiv:2108.07073v1, 2021. \url{https://arxiv.org/abs/2108.07073v1}
\bibitem{ref127} Oasis Data Curation and Assessment System for Pretraining of Large Language Models. arXiv:2311.12537v1, 2023. \url{https://arxiv.org/abs/2311.12537v1}
\bibitem{ref128} Importance-Aware Adaptive Dataset Distillation. arXiv:2401.15863v1, 2024. \url{https://arxiv.org/abs/2401.15863v1}
\bibitem{ref129} Data Distillation Can Be Like Vodka Distilling More Times For Better Quality. arXiv:2310.06982v1, 2023. \url{https://arxiv.org/abs/2310.06982v1}
\bibitem{ref130} A Survey on Dataset Distillation Approaches, Applications and Future Directions. arXiv:2305.01975v3, 2023. \url{https://arxiv.org/abs/2305.01975v3}
\bibitem{ref131} Distill Gold from Massive Ores Efficient Dataset Distillation via Critical Samples Selection. arXiv:2305.18381v3, 2023. \url{https://arxiv.org/abs/2305.18381v3}
\bibitem{ref132} Representation Matters Assessing the Importance of Subgroup Allocations in Training Data. arXiv:2103.03399v2, 2021. \url{https://arxiv.org/abs/2103.03399v2}
\bibitem{ref133} Data Augmentation for Neural Online Chat Response Selection. arXiv:1809.00428v1, 2018. \url{https://arxiv.org/abs/1809.00428v1}
\bibitem{ref134} LatentAugment Data Augmentation via Guided Manipulation of GAN's Latent Space. arXiv:2307.11375v1, 2023. \url{https://arxiv.org/abs/2307.11375v1}
\bibitem{ref135} Leveraging Image-Text Similarity and Caption Modification for the DataComp Challenge Filtering Track and BYOD Track. arXiv:2310.14581v1, 2023. \url{https://arxiv.org/abs/2310.14581v1}
\bibitem{ref136} A Comprehensive Survey of Dataset Distillation. arXiv:2301.05603v4, 2023. \url{https://arxiv.org/abs/2301.05603v4}
\bibitem{ref137} High-Modality Multimodal Transformer Quantifying Modality \& Interaction Heterogeneity for High-Modality Representation Learning. arXiv:2203.01311v4, 2022. \url{https://arxiv.org/abs/2203.01311v4}
\bibitem{ref138} UNIMO-3 Multi-granularity Interaction for Vision-Language Representation Learning. arXiv:2305.13697v1, 2023. \url{https://arxiv.org/abs/2305.13697v1}
\bibitem{ref139} Image as a Foreign Language BEiT Pretraining for All Vision and Vision-Language Tasks. arXiv:2208.10442v2, 2022. \url{https://arxiv.org/abs/2208.10442v2}
\bibitem{ref140} mPLUG-2 A Modularized Multi-modal Foundation Model Across Text, Image and Video. arXiv:2302.00402v1, 2023. \url{https://arxiv.org/abs/2302.00402v1}
\bibitem{ref141} DeepInteraction 3D Object Detection via Modality Interaction. arXiv:2208.11112v4, 2022. \url{https://arxiv.org/abs/2208.11112v4}
\bibitem{ref142} MMICL Empowering Vision-language Model with Multi-Modal In-Context Learning. arXiv:2309.07915v3, 2023. \url{https://arxiv.org/abs/2309.07915v3}
\bibitem{ref143} Mastering Robot Manipulation with Multimodal Prompts through Pretraining and Multi-task Fine-tuning. arXiv:2310.09676v1, 2023. \url{https://arxiv.org/abs/2310.09676v1}
\bibitem{ref144} Assessing GPT4-V on Structured Reasoning Tasks. arXiv:2312.11524v1, 2023. \url{https://arxiv.org/abs/2312.11524v1}
\bibitem{ref145} DDCoT Duty-Distinct Chain-of-Thought Prompting for Multimodal Reasoning in Language Models. arXiv:2310.16436v2, 2023. \url{https://arxiv.org/abs/2310.16436v2}
\bibitem{ref146} SEED-Bench-2 Benchmarking Multimodal Large Language Models. arXiv:2311.17092v1, 2023. \url{https://arxiv.org/abs/2311.17092v1}
\bibitem{ref147} Beyond Static Models and Test Sets Benchmarking the Potential of Pre-trained Models Across Tasks and Languages. arXiv:2205.06356v2, 2022. \url{https://arxiv.org/abs/2205.06356v2}
\bibitem{ref148} Unveiling LLM Evaluation Focused on Metrics Challenges and Solutions. arXiv:2404.09135v1, 2024. \url{https://arxiv.org/abs/2404.09135v1}
\bibitem{ref149} Lost in the Source Language How Large Language Models Evaluate the Quality of Machine Translation. arXiv:2401.06568v1, 2024. \url{https://arxiv.org/abs/2401.06568v1}
\bibitem{ref150} METAL Towards Multilingual Meta-Evaluation. arXiv:2404.01667v1, 2024. \url{https://arxiv.org/abs/2404.01667v1}
\bibitem{ref151} MLLM-Bench, Evaluating Multi-modal LLMs using GPT-4V. arXiv:2311.13951v1, 2023. \url{https://arxiv.org/abs/2311.13951v1}
\bibitem{ref152} ConTextual Evaluating Context-Sensitive Text-Rich Visual Reasoning in Large Multimodal Models. arXiv:2401.13311v1, 2024. \url{https://arxiv.org/abs/2401.13311v1}
\bibitem{ref153} Better Answers to Real Questions. arXiv:1501.05098v1, 2015. \url{https://arxiv.org/abs/1501.05098v1}
\bibitem{ref154} EXAMS A Multi-Subject High School Examinations Dataset for Cross-Lingual and Multilingual Question Answering. arXiv:2011.03080v1, 2020. \url{https://arxiv.org/abs/2011.03080v1}
\bibitem{ref155} HAE-RAE Bench Evaluation of Korean Knowledge in Language Models. arXiv:2309.02706v5, 2023. \url{https://arxiv.org/abs/2309.02706v5}
\bibitem{ref156} ReadMe++ Benchmarking Multilingual Language Models for Multi-Domain Readability Assessment. arXiv:2305.14463v2, 2023. \url{https://arxiv.org/abs/2305.14463v2}
\bibitem{ref157} CPSDBench A Large Language Model Evaluation Benchmark and Baseline for Chinese Public Security Domain. arXiv:2402.07234v3, 2024. \url{https://arxiv.org/abs/2402.07234v3}
\bibitem{ref158} DrBenchmark A Large Language Understanding Evaluation Benchmark for French Biomedical Domain. arXiv:2402.13432v1, 2024. \url{https://arxiv.org/abs/2402.13432v1}
\bibitem{ref159} TransportationGames Benchmarking Transportation Knowledge of (Multimodal) Large Language Models. arXiv:2401.04471v1, 2024. \url{https://arxiv.org/abs/2401.04471v1}
\bibitem{ref160} IGLUE A Benchmark for Transfer Learning across Modalities, Tasks, and Languages. arXiv:2201.11732v2, 2022. \url{https://arxiv.org/abs/2201.11732v2}
\bibitem{ref161} BiToD A Bilingual Multi-Domain Dataset For Task-Oriented Dialogue Modeling. arXiv:2106.02787v1, 2021. \url{https://arxiv.org/abs/2106.02787v1}
\bibitem{ref162} Benchmark Self-Evolving A Multi-Agent Framework for Dynamic LLM Evaluation. arXiv:2402.11443v1, 2024. \url{https://arxiv.org/abs/2402.11443v1}
\bibitem{ref163} Towards Objectively Benchmarking Social Intelligence for Language Agents at Action Level. arXiv:2404.05337v1, 2024. \url{https://arxiv.org/abs/2404.05337v1}
\bibitem{ref164} Large Language Models are Not Yet Human-Level Evaluators for Abstractive Summarization. arXiv:2305.13091v2, 2023. \url{https://arxiv.org/abs/2305.13091v2}
\bibitem{ref165} Experts, Errors, and Context A Large-Scale Study of Human Evaluation for Machine Translation. arXiv:2104.14478v1, 2021. \url{https://arxiv.org/abs/2104.14478v1}
\bibitem{ref166} Correction of Errors in Preference Ratings from Automated Metrics for Text Generation. arXiv:2306.03866v1, 2023. \url{https://arxiv.org/abs/2306.03866v1}
\bibitem{ref167} Collaborative Evaluation Exploring the Synergy of Large Language Models and Humans for Open-ended Generation Evaluation. arXiv:2310.19740v1, 2023. \url{https://arxiv.org/abs/2310.19740v1}
\bibitem{ref168} HD-Eval Aligning Large Language Model Evaluators Through Hierarchical Criteria Decomposition. arXiv:2402.15754v1, 2024. \url{https://arxiv.org/abs/2402.15754v1}
\bibitem{ref169} Estimating Subjective Crowd-Evaluations as an Additional Objective to Improve Natural Language Generation. arXiv:2104.05224v1, 2021. \url{https://arxiv.org/abs/2104.05224v1}
\bibitem{ref170} Which Prompts Make The Difference Data Prioritization For Efficient Human LLM Evaluation. arXiv:2310.14424v1, 2023. \url{https://arxiv.org/abs/2310.14424v1}
\bibitem{ref171} MERA A Comprehensive LLM Evaluation in Russian. arXiv:2401.04531v2, 2024. \url{https://arxiv.org/abs/2401.04531v2}
\bibitem{ref172} Style Over Substance Evaluation Biases for Large Language Models. arXiv:2307.03025v3, 2023. \url{https://arxiv.org/abs/2307.03025v3}
\bibitem{ref173} Rethinking Model Evaluation as Narrowing the Socio-Technical Gap. arXiv:2306.03100v3, 2023. \url{https://arxiv.org/abs/2306.03100v3}
\bibitem{ref174} Misusing Tools in Large Language Models With Visual Adversarial Examples. arXiv:2310.03185v1, 2023. \url{https://arxiv.org/abs/2310.03185v1}
\bibitem{ref175} A Differentiable Language Model Adversarial Attack on Text Classifiers. arXiv:2107.11275v1, 2021. \url{https://arxiv.org/abs/2107.11275v1}
\bibitem{ref176} Scaling Laws for Adversarial Attacks on Language Model Activations. arXiv:2312.02780v1, 2023. \url{https://arxiv.org/abs/2312.02780v1}
\bibitem{ref177} Revisit Input Perturbation Problems for LLMs A Unified Robustness Evaluation Framework for Noisy Slot Filling Task. arXiv:2310.06504v1, 2023. \url{https://arxiv.org/abs/2310.06504v1}
\bibitem{ref178} Adversarial GLUE A Multi-Task Benchmark for Robustness Evaluation of Language Models. arXiv:2111.02840v2, 2021. \url{https://arxiv.org/abs/2111.02840v2}
\bibitem{ref179} Scaling Laws for Generative Mixed-Modal Language Models. arXiv:2301.03728v1, 2023. \url{https://arxiv.org/abs/2301.03728v1}
\bibitem{ref180} Large Language Model Evaluation via Matrix Entropy. arXiv:2401.17139v1, 2024. \url{https://arxiv.org/abs/2401.17139v1}
\bibitem{ref181} Confident Adaptive Language Modeling. arXiv:2207.07061v2, 2022. \url{https://arxiv.org/abs/2207.07061v2}
\bibitem{ref182} Limits of Detecting Text Generated by Large-Scale Language Models. arXiv:2002.03438v1, 2020. \url{https://arxiv.org/abs/2002.03438v1}
\bibitem{ref183} Sparsity-Guided Holistic Explanation for LLMs with Interpretable Inference-Time Intervention. arXiv:2312.15033v1, 2023. \url{https://arxiv.org/abs/2312.15033v1}
\bibitem{ref184} Tuning-Free Accountable Intervention for LLM Deployment -- A Metacognitive Approach. arXiv:2403.05636v1, 2024. \url{https://arxiv.org/abs/2403.05636v1}
\bibitem{ref185} MM-Vet Evaluating Large Multimodal Models for Integrated Capabilities. arXiv:2308.02490v3, 2023. \url{https://arxiv.org/abs/2308.02490v3}
\bibitem{ref186} Establishing Trustworthiness Rethinking Tasks and Model Evaluation. arXiv:2310.05442v2, 2023. \url{https://arxiv.org/abs/2310.05442v2}
\bibitem{ref187} Cognitive Bias in High-Stakes Decision-Making with LLMs. arXiv:2403.00811v1, 2024. \url{https://arxiv.org/abs/2403.00811v1}
\bibitem{ref188} CLUE A Clinical Language Understanding Evaluation for LLMs. arXiv:2404.04067v2, 2024. \url{https://arxiv.org/abs/2404.04067v2}
\bibitem{ref189} PMC-LLaMA Towards Building Open-source Language Models for Medicine. arXiv:2304.14454v3, 2023. \url{https://arxiv.org/abs/2304.14454v3}
\bibitem{ref190} Do Large Language Models understand Medical Codes. arXiv:2403.10822v2, 2024. \url{https://arxiv.org/abs/2403.10822v2}
\bibitem{ref191} ChatCAD Interactive Computer-Aided Diagnosis on Medical Image using Large Language Models. arXiv:2302.07257v1, 2023. \url{https://arxiv.org/abs/2302.07257v1}
\bibitem{ref192} RJUA-MedDQA A Multimodal Benchmark for Medical Document Question Answering and Clinical Reasoning. arXiv:2402.14840v1, 2024. \url{https://arxiv.org/abs/2402.14840v1}
\bibitem{ref193} Evaluating LLM -- Generated Multimodal Diagnosis from Medical Images and Symptom Analysis. arXiv:2402.01730v1, 2024. \url{https://arxiv.org/abs/2402.01730v1}
\bibitem{ref194} EHRTutor Enhancing Patient Understanding of Discharge Instructions. arXiv:2310.19212v1, 2023. \url{https://arxiv.org/abs/2310.19212v1}
\bibitem{ref195} Health-LLM Personalized Retrieval-Augmented Disease Prediction System. arXiv:2402.00746v6, 2024. \url{https://arxiv.org/abs/2402.00746v6}
\bibitem{ref196} DISC-MedLLM Bridging General Large Language Models and Real-World Medical Consultation. arXiv:2308.14346v1, 2023. \url{https://arxiv.org/abs/2308.14346v1}
\bibitem{ref197} MedLM Exploring Language Models for Medical Question Answering Systems. arXiv:2401.11389v2, 2024. \url{https://arxiv.org/abs/2401.11389v2}
\bibitem{ref198} Language models are susceptible to incorrect patient self-diagnosis in medical applications. arXiv:2309.09362v1, 2023. \url{https://arxiv.org/abs/2309.09362v1}
\bibitem{ref199} Revolutionizing Finance with LLMs An Overview of Applications and Insights. arXiv:2401.11641v1, 2024. \url{https://arxiv.org/abs/2401.11641v1}
\bibitem{ref200} Transforming Sentiment Analysis in the Financial Domain with ChatGPT. arXiv:2308.07935v1, 2023. \url{https://arxiv.org/abs/2308.07935v1}
\bibitem{ref201} Are Large Language Models Rational Investors. arXiv:2402.12713v1, 2024. \url{https://arxiv.org/abs/2402.12713v1}
\bibitem{ref202} BloombergGPT A Large Language Model for Finance. arXiv:2303.17564v3, 2023. \url{https://arxiv.org/abs/2303.17564v3}
\bibitem{ref203} Can GPT models be Financial Analysts An Evaluation of ChatGPT and GPT-4 on mock CFA Exams. arXiv:2310.08678v1, 2023. \url{https://arxiv.org/abs/2310.08678v1}
\bibitem{ref204} Evaluating LLMs' Mathematical Reasoning in Financial Document Question Answering. arXiv:2402.11194v2, 2024. \url{https://arxiv.org/abs/2402.11194v2}
\bibitem{ref205} Enhancing Financial Inclusion and Regulatory Challenges A Critical Analysis of Digital Banks and Alternative Lenders Through Digital Platforms, Machine Learning, and Large Language Models Integration. arXiv:2404.11898v1, 2024. \url{https://arxiv.org/abs/2404.11898v1}
\bibitem{ref206} Personality-aware Student Simulation for Conversational Intelligent Tutoring Systems. arXiv:2404.06762v1, 2024. \url{https://arxiv.org/abs/2404.06762v1}
\bibitem{ref207} MathVerse Does Your Multi-modal LLM Truly See the Diagrams in Visual Math Problems. arXiv:2403.14624v1, 2024. \url{https://arxiv.org/abs/2403.14624v1}
\bibitem{ref208} Cognitive network science reveals bias in GPT-3, ChatGPT, and GPT-4 mirroring math anxiety in high-school students. arXiv:2305.18320v1, 2023. \url{https://arxiv.org/abs/2305.18320v1}
\bibitem{ref209} Towards Automated Accessibility Report Generation for Mobile Apps. arXiv:2310.00091v2, 2023. \url{https://arxiv.org/abs/2310.00091v2}
\bibitem{ref210} Advancing Real-time Pandemic Forecasting Using Large Language Models A COVID-19 Case Study. arXiv:2404.06962v1, 2024. \url{https://arxiv.org/abs/2404.06962v1}
\bibitem{ref211} GenSpectrum Chat Data Exploration in Public Health Using Large Language Models. arXiv:2305.13821v1, 2023. \url{https://arxiv.org/abs/2305.13821v1}
\bibitem{ref212} Self-Diagnosis and Large Language Models A New Front for Medical Misinformation. arXiv:2307.04910v1, 2023. \url{https://arxiv.org/abs/2307.04910v1}
\bibitem{ref213} The Last JITAI The Unreasonable Effectiveness of Large Language Models in Issuing Just-in-Time Adaptive Interventions Fostering Physical Activity in a Prospective Cardiac Rehabilitation Setting. arXiv:2402.08658v2, 2024. \url{https://arxiv.org/abs/2402.08658v2}
\bibitem{ref214} A Toolbox for Surfacing Health Equity Harms and Biases in Large Language Models. arXiv:2403.12025v1, 2024. \url{https://arxiv.org/abs/2403.12025v1}
\bibitem{ref215} ClimateGPT Towards AI Synthesizing Interdisciplinary Research on Climate Change. arXiv:2401.09646v1, 2024. \url{https://arxiv.org/abs/2401.09646v1}
\bibitem{ref216} Learning from Multimodal and Multitemporal Earth Observation Data for Building Damage Mapping. arXiv:2009.06200v1, 2020. \url{https://arxiv.org/abs/2009.06200v1}
\bibitem{ref217} UrbanCLIP Learning Text-enhanced Urban Region Profiling with Contrastive Language-Image Pretraining from the Web. arXiv:2310.18340v2, 2023. \url{https://arxiv.org/abs/2310.18340v2}
\bibitem{ref218} FloodBrain Flood Disaster Reporting by Web-based Retrieval Augmented Generation with an LLM. arXiv:2311.02597v1, 2023. \url{https://arxiv.org/abs/2311.02597v1}
\bibitem{ref219} Surveying Attitudinal Alignment Between Large Language Models Vs. Humans Towards 17 Sustainable Development Goals. arXiv:2404.13885v1, 2024. \url{https://arxiv.org/abs/2404.13885v1}
\bibitem{ref220} Towards Human-AI Collaborative Urban Science Research Enabled by Pre-trained Large Language Models. arXiv:2305.11418v1, 2023. \url{https://arxiv.org/abs/2305.11418v1}
\bibitem{ref221} LLaSM Large Language and Speech Model. arXiv:2308.15930v3, 2023. \url{https://arxiv.org/abs/2308.15930v3}
\bibitem{ref222} MLLM-Tool A Multimodal Large Language Model For Tool Agent Learning. arXiv:2401.10727v2, 2024. \url{https://arxiv.org/abs/2401.10727v2}
\bibitem{ref223} PALO A Polyglot Large Multimodal Model for 5B People. arXiv:2402.14818v2, 2024. \url{https://arxiv.org/abs/2402.14818v2}
\bibitem{ref224} Improving Diversity of Demographic Representation in Large Language Models via Collective-Critiques and Self-Voting. arXiv:2310.16523v1, 2023. \url{https://arxiv.org/abs/2310.16523v1}
\bibitem{ref225} The Dawn of LMMs Preliminary Explorations with GPT-4V(ision). arXiv:2309.17421v2, 2023. \url{https://arxiv.org/abs/2309.17421v2}
\bibitem{ref226} Modality-invariant and Specific Prompting for Multimodal Human Perception Understanding. arXiv:2311.10791v1, 2023. \url{https://arxiv.org/abs/2311.10791v1}
\bibitem{ref227} The Ethics of ChatGPT in Medicine and Healthcare A Systematic Review on Large Language Models (LLMs). arXiv:2403.14473v1, 2024. \url{https://arxiv.org/abs/2403.14473v1}
\bibitem{ref228} It's a Fair Game , or Is It Examining How Users Navigate Disclosure Risks and Benefits When Using LLM-Based Conversational Agents. arXiv:2309.11653v2, 2023. \url{https://arxiv.org/abs/2309.11653v2}
\bibitem{ref229} Should ChatGPT be Biased Challenges and Risks of Bias in Large Language Models. arXiv:2304.03738v3, 2023. \url{https://arxiv.org/abs/2304.03738v3}
\bibitem{ref230} Advancing AI with Integrity Ethical Challenges and Solutions in Neural Machine Translation. arXiv:2404.01070v1, 2024. \url{https://arxiv.org/abs/2404.01070v1}
\bibitem{ref231} Regulating Large Language Models A Roundtable Report. arXiv:2403.15397v1, 2024. \url{https://arxiv.org/abs/2403.15397v1}
\bibitem{ref232} Proposing an Interactive Audit Pipeline for Visual Privacy Research. arXiv:2111.03984v2, 2021. \url{https://arxiv.org/abs/2111.03984v2}
\bibitem{ref233} Auditing large language models a three-layered approach. arXiv:2302.08500v2, 2023. \url{https://arxiv.org/abs/2302.08500v2}
\bibitem{ref234} FAIR Enough How Can We Develop and Assess a FAIR-Compliant Dataset for Large Language Models' Training. arXiv:2401.11033v4, 2024. \url{https://arxiv.org/abs/2401.11033v4}
\bibitem{ref235} Applying Standards to Advance Upstream \& Downstream Ethics in Large Language Models. arXiv:2306.03503v2, 2023. \url{https://arxiv.org/abs/2306.03503v2}
\bibitem{ref236} Analyzing Leakage of Personally Identifiable Information in Language Models. arXiv:2302.00539v4, 2023. \url{https://arxiv.org/abs/2302.00539v4}
\bibitem{ref237} Quantifying Association Capabilities of Large Language Models and Its Implications on Privacy Leakage. arXiv:2305.12707v2, 2023. \url{https://arxiv.org/abs/2305.12707v2}
\bibitem{ref238} Are Large Pre-Trained Language Models Leaking Your Personal Information. arXiv:2205.12628v2, 2022. \url{https://arxiv.org/abs/2205.12628v2}
\bibitem{ref239} Membership Inference Attacks and Privacy in Topic Modeling. arXiv:2403.04451v1, 2024. \url{https://arxiv.org/abs/2403.04451v1}
\bibitem{ref240} Privacy Implications of Retrieval-Based Language Models. arXiv:2305.14888v1, 2023. \url{https://arxiv.org/abs/2305.14888v1}
\bibitem{ref241} Context-Aware Differential Privacy for Language Modeling. arXiv:2301.12288v1, 2023. \url{https://arxiv.org/abs/2301.12288v1}
\bibitem{ref242} Selective Differential Privacy for Language Modeling. arXiv:2108.12944v3, 2021. \url{https://arxiv.org/abs/2108.12944v3}
\bibitem{ref243} Differentially Private Attention Computation. arXiv:2305.04701v1, 2023. \url{https://arxiv.org/abs/2305.04701v1}
\bibitem{ref244} Privacy-preserving Fine-tuning of Large Language Models through Flatness. arXiv:2403.04124v1, 2024. \url{https://arxiv.org/abs/2403.04124v1}
\bibitem{ref245} Hide and Seek (HaS) A Lightweight Framework for Prompt Privacy Protection. arXiv:2309.03057v1, 2023. \url{https://arxiv.org/abs/2309.03057v1}
\bibitem{ref246} Reducing Privacy Risks in Online Self-Disclosures with Language Models. arXiv:2311.09538v2, 2023. \url{https://arxiv.org/abs/2311.09538v2}
\bibitem{ref247} Privacy in Large Language Models Attacks, Defenses and Future Directions. arXiv:2310.10383v1, 2023. \url{https://arxiv.org/abs/2310.10383v1}
\bibitem{ref248} Planting and Mitigating Memorized Content in Predictive-Text Language Models. arXiv:2212.08619v1, 2022. \url{https://arxiv.org/abs/2212.08619v1}
\bibitem{ref249} Last One Standing A Comparative Analysis of Security and Privacy of Soft Prompt Tuning, LoRA, and In-Context Learning. arXiv:2310.11397v1, 2023. \url{https://arxiv.org/abs/2310.11397v1}
\bibitem{ref250} Towards detecting unanticipated bias in Large Language Models. arXiv:2404.02650v1, 2024. \url{https://arxiv.org/abs/2404.02650v1}
\bibitem{ref251} Quantifying and Mitigating Unimodal Biases in Multimodal Large Language Models A Causal Perspective. arXiv:2403.18346v3, 2024. \url{https://arxiv.org/abs/2403.18346v3}
\bibitem{ref252} Bias and unfairness in machine learning models a systematic literature review. arXiv:2202.08176v4, 2022. \url{https://arxiv.org/abs/2202.08176v4}
\bibitem{ref253} Bias and Fairness on Multimodal Emotion Detection Algorithms. arXiv:2205.08383v1, 2022. \url{https://arxiv.org/abs/2205.08383v1}
\bibitem{ref254} Fairness Evaluation in Text Classification Machine Learning Practitioner Perspectives of Individual and Group Fairness. arXiv:2303.00673v1, 2023. \url{https://arxiv.org/abs/2303.00673v1}
\bibitem{ref255} ROBBIE Robust Bias Evaluation of Large Generative Language Models. arXiv:2311.18140v1, 2023. \url{https://arxiv.org/abs/2311.18140v1}
\bibitem{ref256} A Group Fairness Lens for Large Language Models. arXiv:2312.15478v1, 2023. \url{https://arxiv.org/abs/2312.15478v1}
\bibitem{ref257} Mapping the Multilingual Margins Intersectional Biases of Sentiment Analysis Systems in English, Spanish, and Arabic. arXiv:2204.03558v1, 2022. \url{https://arxiv.org/abs/2204.03558v1}
\bibitem{ref258} Locating and Mitigating Gender Bias in Large Language Models. arXiv:2403.14409v1, 2024. \url{https://arxiv.org/abs/2403.14409v1}
\bibitem{ref259} ChatGPT Based Data Augmentation for Improved Parameter-Efficient Debiasing of LLMs. arXiv:2402.11764v1, 2024. \url{https://arxiv.org/abs/2402.11764v1}
\bibitem{ref260} Steering LLMs Towards Unbiased Responses A Causality-Guided Debiasing Framework. arXiv:2403.08743v1, 2024. \url{https://arxiv.org/abs/2403.08743v1}
\bibitem{ref261} Toward Operationalizing Pipeline-aware ML Fairness A Research Agenda for Developing Practical Guidelines and Tools. arXiv:2309.17337v1, 2023. \url{https://arxiv.org/abs/2309.17337v1}
\bibitem{ref262} Deception and Manipulation in Generative AI. arXiv:2401.11335v1, 2024. \url{https://arxiv.org/abs/2401.11335v1}
\bibitem{ref263} Can LLM-Generated Misinformation Be Detected. arXiv:2309.13788v5, 2023. \url{https://arxiv.org/abs/2309.13788v5}
\bibitem{ref264} Disinformation Detection An Evolving Challenge in the Age of LLMs. arXiv:2309.15847v1, 2023. \url{https://arxiv.org/abs/2309.15847v1}
\bibitem{ref265} From Text to MITRE Techniques Exploring the Malicious Use of Large Language Models for Generating Cyber Attack Payloads. arXiv:2305.15336v1, 2023. \url{https://arxiv.org/abs/2305.15336v1}
\bibitem{ref266} Exploiting Large Language Models (LLMs) through Deception Techniques and Persuasion Principles. arXiv:2311.14876v1, 2023. \url{https://arxiv.org/abs/2311.14876v1}
\bibitem{ref267} Survey of Vulnerabilities in Large Language Models Revealed by Adversarial Attacks. arXiv:2310.10844v1, 2023. \url{https://arxiv.org/abs/2310.10844v1}
\bibitem{ref268} Breaking Down the Defenses A Comparative Survey of Attacks on Large Language Models. arXiv:2403.04786v2, 2024. \url{https://arxiv.org/abs/2403.04786v2}
\bibitem{ref269} Sandwich attack Multi-language Mixture Adaptive Attack on LLMs. arXiv:2404.07242v1, 2024. \url{https://arxiv.org/abs/2404.07242v1}
\bibitem{ref270} Attacks, Defenses and Evaluations for LLM Conversation Safety A Survey. arXiv:2402.09283v3, 2024. \url{https://arxiv.org/abs/2402.09283v3}
\bibitem{ref271} RigorLLM Resilient Guardrails for Large Language Models against Undesired Content. arXiv:2403.13031v1, 2024. \url{https://arxiv.org/abs/2403.13031v1}
\bibitem{ref272} A Security Risk Taxonomy for Large Language Models. arXiv:2311.11415v1, 2023. \url{https://arxiv.org/abs/2311.11415v1}
\bibitem{ref273} Risk and Response in Large Language Models Evaluating Key Threat Categories. arXiv:2403.14988v1, 2024. \url{https://arxiv.org/abs/2403.14988v1}
\bibitem{ref274} AI Governance and Ethics Framework for Sustainable AI and Sustainability. arXiv:2210.08984v1, 2022. \url{https://arxiv.org/abs/2210.08984v1}
\bibitem{ref275} (\$S\$,\$N\$,\$T\$)-Implications. arXiv:2106.15746v1, 2021. \url{https://arxiv.org/abs/2106.15746v1}
\bibitem{ref276} Synthetic Data Methods, Use Cases, and Risks. arXiv:2303.01230v3, 2023. \url{https://arxiv.org/abs/2303.01230v3}
\bibitem{ref277} Ethical Considerations and Policy Implications for Large Language Models Guiding Responsible Development and Deployment. arXiv:2308.02678v1, 2023. \url{https://arxiv.org/abs/2308.02678v1}
\bibitem{ref278} Compromise in Multilateral Negotiations and the Global Regulation of Artificial Intelligence. arXiv:2309.17158v1, 2023. \url{https://arxiv.org/abs/2309.17158v1}
\bibitem{ref279} Bridging Deliberative Democracy and Deployment of Societal-Scale Technology. arXiv:2303.10831v2, 2023. \url{https://arxiv.org/abs/2303.10831v2}
\bibitem{ref280} Dual Governance The intersection of centralized regulation and crowdsourced safety mechanisms for Generative AI. arXiv:2308.04448v1, 2023. \url{https://arxiv.org/abs/2308.04448v1}
\bibitem{ref281} Putting AI Ethics into Practice The Hourglass Model of Organizational AI Governance. arXiv:2206.00335v2, 2022. \url{https://arxiv.org/abs/2206.00335v2}
\bibitem{ref282} Towards a Responsible AI Metrics Catalogue A Collection of Metrics for AI Accountability. arXiv:2311.13158v3, 2023. \url{https://arxiv.org/abs/2311.13158v3}
\bibitem{ref283} Use large language models to promote equity. arXiv:2312.14804v1, 2023. \url{https://arxiv.org/abs/2312.14804v1}
\bibitem{ref284} GPTs are GPTs An Early Look at the Labor Market Impact Potential of Large Language Models. arXiv:2303.10130v5, 2023. \url{https://arxiv.org/abs/2303.10130v5}
\bibitem{ref285} Apprentices to Research Assistants Advancing Research with Large Language Models. arXiv:2404.06404v1, 2024. \url{https://arxiv.org/abs/2404.06404v1}
\bibitem{ref286} Future-proofing Education A Prototype for Simulating Oral Examinations Using Large Language Models. arXiv:2401.06160v1, 2023. \url{https://arxiv.org/abs/2401.06160v1}
\bibitem{ref287} Large Language Models in Education Vision and Opportunities. arXiv:2311.13160v1, 2023. \url{https://arxiv.org/abs/2311.13160v1}
\bibitem{ref288} Gender, Age, and Technology Education Influence the Adoption and Appropriation of LLMs. arXiv:2310.06556v1, 2023. \url{https://arxiv.org/abs/2310.06556v1}
\bibitem{ref289} Friend or Foe Exploring the Implications of Large Language Models on the Science System. arXiv:2306.09928v1, 2023. \url{https://arxiv.org/abs/2306.09928v1}
\bibitem{ref290} Chain-of-Thought Prompting Elicits Reasoning in Large Language Models. arXiv:2201.11903v6, 2022. \url{https://arxiv.org/abs/2201.11903v6}
\bibitem{ref291} Towards Robust Multi-Modal Reasoning via Model Selection. arXiv:2310.08446v2, 2023. \url{https://arxiv.org/abs/2310.08446v2}
\bibitem{ref292} Thought Propagation An Analogical Approach to Complex Reasoning with Large Language Models. arXiv:2310.03965v2, 2023. \url{https://arxiv.org/abs/2310.03965v2}
\bibitem{ref293} Visual In-Context Learning for Large Vision-Language Models. arXiv:2402.11574v1, 2024. \url{https://arxiv.org/abs/2402.11574v1}
\bibitem{ref294} Towards More Unified In-context Visual Understanding. arXiv:2312.02520v2, 2023. \url{https://arxiv.org/abs/2312.02520v2}
\bibitem{ref295} Improving In-context Learning via Bidirectional Alignment. arXiv:2312.17055v1, 2023. \url{https://arxiv.org/abs/2312.17055v1}
\bibitem{ref296} How Does the Textual Information Affect the Retrieval of Multimodal In-Context Learning. arXiv:2404.12866v1, 2024. \url{https://arxiv.org/abs/2404.12866v1}
\bibitem{ref297} What Makes Multimodal In-Context Learning Work. arXiv:2404.15736v2, 2024. \url{https://arxiv.org/abs/2404.15736v2}
\bibitem{ref298} IMProv Inpainting-based Multimodal Prompting for Computer Vision Tasks. arXiv:2312.01771v1, 2023. \url{https://arxiv.org/abs/2312.01771v1}
\bibitem{ref299} Visualization-of-Thought Elicits Spatial Reasoning in Large Language Models. arXiv:2404.03622v1, 2024. \url{https://arxiv.org/abs/2404.03622v1}
\bibitem{ref300} VURF A General-purpose Reasoning and Self-refinement Framework for Video Understanding. arXiv:2403.14743v2, 2024. \url{https://arxiv.org/abs/2403.14743v2}
\bibitem{ref301} Structured Prompting Scaling In-Context Learning to 1,000 Examples. arXiv:2212.06713v1, 2022. \url{https://arxiv.org/abs/2212.06713v1}
\bibitem{ref302} Efficient Multimodal Fusion via Interactive Prompting. arXiv:2304.06306v2, 2023. \url{https://arxiv.org/abs/2304.06306v2}
\bibitem{ref303} In-context Example Selection with Influences. arXiv:2302.11042v2, 2023. \url{https://arxiv.org/abs/2302.11042v2}
\bibitem{ref304} How are Prompts Different in Terms of Sensitivity. arXiv:2311.07230v1, 2023. \url{https://arxiv.org/abs/2311.07230v1}
\bibitem{ref305} Active Example Selection for In-Context Learning. arXiv:2211.04486v1, 2022. \url{https://arxiv.org/abs/2211.04486v1}
\bibitem{ref306} Towards Understanding and Mitigating Social Biases in Language Models. arXiv:2106.13219v1, 2021. \url{https://arxiv.org/abs/2106.13219v1}
\bibitem{ref307} Fairness-guided Few-shot Prompting for Large Language Models. arXiv:2303.13217v3, 2023. \url{https://arxiv.org/abs/2303.13217v3}
\bibitem{ref308} On the Out-Of-Distribution Generalization of Multimodal Large Language Models. arXiv:2402.06599v1, 2024. \url{https://arxiv.org/abs/2402.06599v1}
\bibitem{ref309} Enhancing In-Context Learning with Answer Feedback for Multi-Span Question Answering. arXiv:2306.04508v1, 2023. \url{https://arxiv.org/abs/2306.04508v1}
\bibitem{ref310} The (R)Evolution of Multimodal Large Language Models A Survey. arXiv:2402.12451v1, 2024. \url{https://arxiv.org/abs/2402.12451v1}
\bibitem{ref311} MaPLe Multi-modal Prompt Learning. arXiv:2210.03117v3, 2022. \url{https://arxiv.org/abs/2210.03117v3}
\bibitem{ref312} Joint Visual and Text Prompting for Improved Object-Centric Perception with Multimodal Large Language Models. arXiv:2404.04514v1, 2024. \url{https://arxiv.org/abs/2404.04514v1}
\bibitem{ref313} Causal-CoG A Causal-Effect Look at Context Generation for Boosting Multi-modal Language Models. arXiv:2312.06685v1, 2023. \url{https://arxiv.org/abs/2312.06685v1}
\bibitem{ref314} Browse and Concentrate Comprehending Multimodal Content via prior-LLM Context Fusion. arXiv:2402.12195v1, 2024. \url{https://arxiv.org/abs/2402.12195v1}
\bibitem{ref315} LION Empowering Multimodal Large Language Model with Dual-Level Visual Knowledge. arXiv:2311.11860v2, 2023. \url{https://arxiv.org/abs/2311.11860v2}
\bibitem{ref316} NPHardEval4V A Dynamic Reasoning Benchmark of Multimodal Large Language Models. arXiv:2403.01777v2, 2024. \url{https://arxiv.org/abs/2403.01777v2}
\bibitem{ref317} Evaluating the Capabilities of Multi-modal Reasoning Models with Synthetic Task Data. arXiv:2306.01144v1, 2023. \url{https://arxiv.org/abs/2306.01144v1}
\bibitem{ref318} ReForm-Eval Evaluating Large Vision Language Models via Unified Re-Formulation of Task-Oriented Benchmarks. arXiv:2310.02569v2, 2023. \url{https://arxiv.org/abs/2310.02569v2}
\bibitem{ref319} Cantor Inspiring Multimodal Chain-of-Thought of MLLM. arXiv:2404.16033v1, 2024. \url{https://arxiv.org/abs/2404.16033v1}
\bibitem{ref320} Cognitive bias in large language models Cautious optimism meets anti-Panglossian meliorism. arXiv:2311.10932v1, 2023. \url{https://arxiv.org/abs/2311.10932v1}
\bibitem{ref321} Investigating Bias Representations in Llama 2 Chat via Activation Steering. arXiv:2402.00402v1, 2024. \url{https://arxiv.org/abs/2402.00402v1}
\bibitem{ref322} How Susceptible are Large Language Models to Ideological Manipulation. arXiv:2402.11725v2, 2024. \url{https://arxiv.org/abs/2402.11725v2}
\bibitem{ref323} A Survey on Fairness in Large Language Models. arXiv:2308.10149v2, 2023. \url{https://arxiv.org/abs/2308.10149v2}
\bibitem{ref324} Detecting Bias in Large Language Models Fine-tuned KcBERT. arXiv:2403.10774v1, 2024. \url{https://arxiv.org/abs/2403.10774v1}
\bibitem{ref325} Use of LLMs for Illicit Purposes Threats, Prevention Measures, and Vulnerabilities. arXiv:2308.12833v1, 2023. \url{https://arxiv.org/abs/2308.12833v1}
\bibitem{ref326} Break the Breakout Reinventing LM Defense Against Jailbreak Attacks with Self-Refinement. arXiv:2402.15180v2, 2024. \url{https://arxiv.org/abs/2402.15180v2}
\bibitem{ref327} Speak Out of Turn Safety Vulnerability of Large Language Models in Multi-turn Dialogue. arXiv:2402.17262v1, 2024. \url{https://arxiv.org/abs/2402.17262v1}
\bibitem{ref328} MM-SafetyBench A Benchmark for Safety Evaluation of Multimodal Large Language Models. arXiv:2311.17600v2, 2023. \url{https://arxiv.org/abs/2311.17600v2}
\bibitem{ref329} Safety-Tuned LLaMAs Lessons From Improving the Safety of Large Language Models that Follow Instructions. arXiv:2309.07875v3, 2023. \url{https://arxiv.org/abs/2309.07875v3}
\bibitem{ref330} GrounDial Human-norm Grounded Safe Dialog Response Generation. arXiv:2402.08968v1, 2024. \url{https://arxiv.org/abs/2402.08968v1}
\bibitem{ref331} Alignment is not sufficient to prevent large language models from generating harmful information A psychoanalytic perspective. arXiv:2311.08487v1, 2023. \url{https://arxiv.org/abs/2311.08487v1}
\bibitem{ref332} Towards Uncovering How Large Language Model Works An Explainability Perspective. arXiv:2402.10688v2, 2024. \url{https://arxiv.org/abs/2402.10688v2}
\bibitem{ref333} Language Models Don't Always Say What They Think Unfaithful Explanations in Chain-of-Thought Prompting. arXiv:2305.04388v2, 2023. \url{https://arxiv.org/abs/2305.04388v2}
\bibitem{ref334} Are Large Language Models Post Hoc Explainers. arXiv:2310.05797v3, 2023. \url{https://arxiv.org/abs/2310.05797v3}
\bibitem{ref335} Charting New Territories Exploring the Geographic and Geospatial Capabilities of Multimodal LLMs. arXiv:2311.14656v3, 2023. \url{https://arxiv.org/abs/2311.14656v3}
\bibitem{ref336} A Survey of Resource-efficient LLM and Multimodal Foundation Models. arXiv:2401.08092v1, 2024. \url{https://arxiv.org/abs/2401.08092v1}
\bibitem{ref337} Characterization of Large Language Model Development in the Datacenter. arXiv:2403.07648v2, 2024. \url{https://arxiv.org/abs/2403.07648v2}
\bibitem{ref338} Pre-gated MoE An Algorithm-System Co-Design for Fast and Scalable Mixture-of-Expert Inference. arXiv:2308.12066v2, 2023. \url{https://arxiv.org/abs/2308.12066v2}
\bibitem{ref339} Towards Better Parameter-Efficient Fine-Tuning for Large Language Models A Position Paper. arXiv:2311.13126v1, 2023. \url{https://arxiv.org/abs/2311.13126v1}
\bibitem{ref340} Modality Plug-and-Play Elastic Modality Adaptation in Multimodal LLMs for Embodied AI. arXiv:2312.07886v1, 2023. \url{https://arxiv.org/abs/2312.07886v1}
\bibitem{ref341} The Efficiency Spectrum of Large Language Models An Algorithmic Survey. arXiv:2312.00678v2, 2023. \url{https://arxiv.org/abs/2312.00678v2}
\bibitem{ref342} Small Language Models Fine-tuned to Coordinate Larger Language Models improve Complex Reasoning. arXiv:2310.18338v2, 2023. \url{https://arxiv.org/abs/2310.18338v2}
\bibitem{ref343} Domain Specialization as the Key to Make Large Language Models Disruptive A Comprehensive Survey. arXiv:2305.18703v7, 2023. \url{https://arxiv.org/abs/2305.18703v7}
\bibitem{ref344} Beyond the Imitation Game Quantifying and extrapolating the capabilities of language models. arXiv:2206.04615v3, 2022. \url{https://arxiv.org/abs/2206.04615v3}
\bibitem{ref345} LLMs as Potential Brainstorming Partners for Math and Science Problems. arXiv:2310.10677v1, 2023. \url{https://arxiv.org/abs/2310.10677v1}
\bibitem{ref346} Scientists' Perspectives on the Potential for Generative AI in their Fields. arXiv:2304.01420v1, 2023. \url{https://arxiv.org/abs/2304.01420v1}
\bibitem{ref347} Scientific Language Modeling A Quantitative Review of Large Language Models in Molecular Science. arXiv:2402.04119v1, 2024. \url{https://arxiv.org/abs/2402.04119v1}
\bibitem{ref348} Understanding the Capabilities, Limitations, and Societal Impact of Large Language Models. arXiv:2102.02503v1, 2021. \url{https://arxiv.org/abs/2102.02503v1}
\bibitem{ref349} MarineGPT Unlocking Secrets of Ocean to the Public. arXiv:2310.13596v1, 2023. \url{https://arxiv.org/abs/2310.13596v1}
\bibitem{ref350} Driving with LLMs Fusing Object-Level Vector Modality for Explainable Autonomous Driving. arXiv:2310.01957v2, 2023. \url{https://arxiv.org/abs/2310.01957v2}
\bibitem{ref351} Univalent polymorphism. arXiv:1803.10113v2, 2018. \url{https://arxiv.org/abs/1803.10113v2}
\bibitem{ref352} SpeechGPT-Gen Scaling Chain-of-Information Speech Generation. arXiv:2401.13527v2, 2024. \url{https://arxiv.org/abs/2401.13527v2}
\bibitem{ref353} A Note on Time Measurements in LAMMPS. arXiv:1602.05566v1, 2016. \url{https://arxiv.org/abs/1602.05566v1}
\bibitem{ref354} mPLUG-Owl2 Revolutionizing Multi-modal Large Language Model with Modality Collaboration. arXiv:2311.04257v2, 2023. \url{https://arxiv.org/abs/2311.04257v2}
\bibitem{ref355} Multiplying Matrices Without Multiplying. arXiv:2106.10860v1, 2021. \url{https://arxiv.org/abs/2106.10860v1}
\bibitem{ref356} The Universe of Minds. arXiv:1410.0369v1, 2014. \url{https://arxiv.org/abs/1410.0369v1}
\bibitem{ref357} Multimodal Densenet. arXiv:1811.07407v1, 2018. \url{https://arxiv.org/abs/1811.07407v1}
\bibitem{ref358} DDC A Vision for a Disaggregated Datacenter. arXiv:2402.12742v1, 2024. \url{https://arxiv.org/abs/2402.12742v1}
\bibitem{ref359} Kosmos-2.5 A Multimodal Literate Model. arXiv:2309.11419v1, 2023. \url{https://arxiv.org/abs/2309.11419v1}
\bibitem{ref360} LLMBind A Unified Modality-Task Integration Framework. arXiv:2402.14891v5, 2024. \url{https://arxiv.org/abs/2402.14891v5}
\bibitem{ref361} Lost in Interpreting Speech Translation from Source or Interpreter. arXiv:2106.09343v1, 2021. \url{https://arxiv.org/abs/2106.09343v1}
\bibitem{ref362} Research Re search \& Re-search. arXiv:2403.13705v1, 2024. \url{https://arxiv.org/abs/2403.13705v1}
\bibitem{ref363} SMS in PACE 2020. arXiv:2006.07302v1, 2020. \url{https://arxiv.org/abs/2006.07302v1}
\bibitem{ref364} Stable Voting. arXiv:2108.00542v9, 2021. \url{https://arxiv.org/abs/2108.00542v9}
\bibitem{ref365} Science in the Era of ChatGPT, Large Language Models and Generative AI Challenges for Research Ethics and How to Respond. arXiv:2305.15299v4, 2023. \url{https://arxiv.org/abs/2305.15299v4}
\bibitem{ref366} You Are What You Write Preserving Privacy in the Era of Large Language Models. arXiv:2204.09391v1, 2022. \url{https://arxiv.org/abs/2204.09391v1}
\bibitem{ref367} Unveiling Security, Privacy, and Ethical Concerns of ChatGPT. arXiv:2307.14192v1, 2023. \url{https://arxiv.org/abs/2307.14192v1}
\bibitem{ref368} AI Ethics and Governance in Practice An Introduction. arXiv:2403.15403v1, 2024. \url{https://arxiv.org/abs/2403.15403v1}
\bibitem{ref369} Denevil Towards Deciphering and Navigating the Ethical Values of Large Language Models via Instruction Learning. arXiv:2310.11053v3, 2023. \url{https://arxiv.org/abs/2310.11053v3}
\bibitem{ref370} Ever heard of ethical AI Investigating the salience of ethical AI issues among the German population. arXiv:2207.14086v1, 2022. \url{https://arxiv.org/abs/2207.14086v1}
\bibitem{ref371} Worldwide AI Ethics a review of 200 guidelines and recommendations for AI governance. arXiv:2206.11922v7, 2022. \url{https://arxiv.org/abs/2206.11922v7}
\bibitem{ref372} The Social Impact of Generative AI An Analysis on ChatGPT. arXiv:2403.04667v1, 2024. \url{https://arxiv.org/abs/2403.04667v1}
\bibitem{ref373} ChatGPT Needs SPADE (Sustainability, PrivAcy, Digital divide, and Ethics) Evaluation A Review. arXiv:2305.03123v3, 2023. \url{https://arxiv.org/abs/2305.03123v3}
\bibitem{ref374} A Survey on Integration of Large Language Models with Intelligent Robots. arXiv:2404.09228v1, 2024. \url{https://arxiv.org/abs/2404.09228v1}
\bibitem{ref375} Emergent autonomous scientific research capabilities of large language models. arXiv:2304.05332v1, 2023. \url{https://arxiv.org/abs/2304.05332v1}
\bibitem{ref376} A Survey for Foundation Models in Autonomous Driving. arXiv:2402.01105v1, 2024. \url{https://arxiv.org/abs/2402.01105v1}
\bibitem{ref377} Large Language Models for Business Process Management Opportunities and Challenges. arXiv:2304.04309v1, 2023. \url{https://arxiv.org/abs/2304.04309v1}
\bibitem{ref378} Leveraging Large Language Model for Automatic Evolving of Industrial Data-Centric R\&D Cycle. arXiv:2310.11249v1, 2023. \url{https://arxiv.org/abs/2310.11249v1}
\end{thebibliography}

\end{document}
